\documentclass{situationist-preprint}

\usepackage{longtable}
\usepackage{array}
\usepackage{needspace}

\newcolumntype{P}[1]{>{\raggedright\arraybackslash}p{#1}}
\newenvironment{manualtable}{\par\footnotesize\setlength{\tabcolsep}{3pt}\setlength{\LTleft}{0pt}\setlength{\LTright}{0pt}%
  \begin{longtable}{@{}>{\ttfamily}P{1.1cm}P{5.0cm}P{3.75cm}P{2.25cm}P{2.95cm}@{}}
  \toprule \normalfont\textbf{ID} & \textbf{$T_1$: term and role / token context} & \textbf{$T_2$: image / token referent} & \textbf{status} & \textbf{inferences}\\ \midrule\endhead
  \bottomrule\endfoot}%
  {\end{longtable}\addtocounter{table}{-1}}
\newenvironment{inftable}{\par\footnotesize\setlength{\tabcolsep}{3pt}\setlength{\LTleft}{0pt}\setlength{\LTright}{0pt}%
  \begin{longtable}{@{}>{\ttfamily}P{1.05cm}P{4.75cm}P{6.15cm}>{\ttfamily\raggedright\arraybackslash}P{2.35cm}c@{}}
  \toprule \normalfont\textbf{ID} & \textbf{inference of $T_1$} & \textbf{translation $M(I)$, or interpretation} & \normalfont\textbf{readings} & \textbf{v.}\\ \midrule\endhead}%
  {\end{longtable}\addtocounter{table}{-1}}
\newenvironment{crittable}{\par\footnotesize\setlength{\tabcolsep}{3pt}\setlength{\LTleft}{0pt}\setlength{\LTright}{0pt}%
  \begin{longtable}{@{}>{\ttfamily}P{1.2cm}P{2.0cm}cP{10.7cm}@{}}
  \toprule \normalfont\textbf{Rev.} & \textbf{criterion} & & \textbf{documented outcome}\\ \midrule\endhead}%
  {\end{longtable}\addtocounter{table}{-1}}
\newenvironment{evidtable}{\par\footnotesize\setlength{\tabcolsep}{3pt}\setlength{\LTleft}{0pt}\setlength{\LTright}{0pt}%
  \begin{longtable}{@{}>{\ttfamily}P{1.05cm}P{8.4cm}>{\ttfamily\raggedright\arraybackslash}P{3.2cm}P{2.6cm}@{}}
  \caption{Evidential claims behind the favourable outcomes, and the readings they use.}\label{tab:evid}\\
  \toprule \normalfont\textbf{ID} & \textbf{evidential claim} & \normalfont\textbf{readings} & \textbf{criteria}\\ \midrule\endfirsthead
  \toprule \normalfont\textbf{ID} & \textbf{evidential claim} & \normalfont\textbf{readings} & \textbf{criteria}\\ \midrule\endhead}%
  {\end{longtable}}

\title{Commensurability Across Paradigms\\[0.4em]\large Volume I: Translation and Theory Choice}
\msid{SCI-06-a}
\msfamily{SCI-06}
\msdiscipline{Philosophy of Science}
\mstype{Principal manuscript, Volume I of III}
\mskeywords{incommensurability; Kuhn; Feyerabend; translation; reference potential; partial denotation; theory choice; scientific revolutions; translation manuals}

\newcommand{\nrevs}{10}
\newcommand{\ntermsall}{221}
\newcommand{\ntokensall}{294}
\newcommand{\nrefall}{241}
\newcommand{\ninfall}{300}
\newcommand{\nstatE}{60}
\newcommand{\nstatP}{109}
\newcommand{\nstatT}{36}
\newcommand{\nstatN}{16}
\newcommand{\nverdH}{149}
\newcommand{\nverdR}{34}
\newcommand{\nverdX}{41}
\newcommand{\nverdU}{76}
\newcommand{\ndominant}{8}
\newcommand{\nevidall}{45}
\newcommand{\nclusterall}{47}
\newcommand{\ncoreall}{30}
\newcommand{\timin}{0.783}
\newcommand{\timax}{0.925}
\newcommand{\kappamin}{0.633}
\newcommand{\kappamax}{0.839}
\newcommand{\sensmin}{0.467}
\newcommand{\sensmax}{0.937}
\newcommand{\tiChem}{0.879}
\newcommand{\kappaChem}{0.710}
\newcommand{\piChem}{0.645}
\newcommand{\ntermsChem}{23}
\newcommand{\ninfChem}{31}
\newcommand{\nclChem}{3}
\newcommand{\nXChem}{2}
\newcommand{\nUChem}{9}
\newcommand{\nplusChem}{2}
\newcommand{\tiAstr}{0.887}
\newcommand{\kappaAstr}{0.824}
\newcommand{\piAstr}{0.647}
\newcommand{\ntermsAstr}{25}
\newcommand{\ninfAstr}{34}
\newcommand{\nclAstr}{5}
\newcommand{\nXAstr}{6}
\newcommand{\nUAstr}{6}
\newcommand{\nplusAstr}{4}
\newcommand{\tiMech}{0.845}
\newcommand{\kappaMech}{0.743}
\newcommand{\piMech}{0.714}
\newcommand{\ntermsMech}{23}
\newcommand{\ninfMech}{35}
\newcommand{\nclMech}{7}
\newcommand{\nXMech}{1}
\newcommand{\nUMech}{9}
\newcommand{\nplusMech}{4}
\newcommand{\tiHeat}{0.897}
\newcommand{\kappaHeat}{0.700}
\newcommand{\piHeat}{0.633}
\newcommand{\ntermsHeat}{23}
\newcommand{\ninfHeat}{30}
\newcommand{\nclHeat}{3}
\newcommand{\nXHeat}{2}
\newcommand{\nUHeat}{9}
\newcommand{\nplusHeat}{3}
\newcommand{\tiQuan}{0.925}
\newcommand{\kappaQuan}{0.833}
\newcommand{\piQuan}{0.600}
\newcommand{\ntermsQuan}{21}
\newcommand{\ninfQuan}{30}
\newcommand{\nclQuan}{3}
\newcommand{\nXQuan}{7}
\newcommand{\nUQuan}{5}
\newcommand{\nplusQuan}{4}
\newcommand{\tiBiol}{0.908}
\newcommand{\kappaBiol}{0.839}
\newcommand{\piBiol}{0.581}
\newcommand{\ntermsBiol}{24}
\newcommand{\ninfBiol}{31}
\newcommand{\nclBiol}{4}
\newcommand{\nXBiol}{8}
\newcommand{\nUBiol}{5}
\newcommand{\nplusBiol}{3}
\newcommand{\tiPhys}{0.783}
\newcommand{\kappaPhys}{0.633}
\newcommand{\piPhys}{0.433}
\newcommand{\ntermsPhys}{26}
\newcommand{\ninfPhys}{30}
\newcommand{\nclPhys}{8}
\newcommand{\nXPhys}{6}
\newcommand{\nUPhys}{11}
\newcommand{\nplusPhys}{3}
\newcommand{\tiOpti}{0.812}
\newcommand{\kappaOpti}{0.720}
\newcommand{\piOpti}{0.640}
\newcommand{\ntermsOpti}{18}
\newcommand{\ninfOpti}{25}
\newcommand{\nclOpti}{3}
\newcommand{\nXOpti}{2}
\newcommand{\nUOpti}{7}
\newcommand{\nplusOpti}{3}
\newcommand{\tiMias}{0.793}
\newcommand{\kappaMias}{0.645}
\newcommand{\piMias}{0.581}
\newcommand{\ntermsMias}{21}
\newcommand{\ninfMias}{31}
\newcommand{\nclMias}{7}
\newcommand{\nXMias}{2}
\newcommand{\nUMias}{11}
\newcommand{\nplusMias}{3}
\newcommand{\tiHere}{0.895}
\newcommand{\kappaHere}{0.826}
\newcommand{\piHere}{0.609}
\newcommand{\ntermsHere}{17}
\newcommand{\ninfHere}{23}
\newcommand{\nclHere}{4}
\newcommand{\nXHere}{5}
\newcommand{\nUHere}{4}
\newcommand{\nplusHere}{3}


\begin{document}
\maketitle

\begin{abstract}
Are the paradigms on either side of a scientific revolution incommensurable? We argue that they are not: successive paradigms are partially translatable, and theory choice remains rational across revolutions. We define a translation manual as a partial map from expressions of the earlier theory to expressions of the later one that preserves the reference of term tokens, in Kitcher's sense, and a specified set of inferential relations, and we define a translation index measuring the share of the earlier vocabulary that the manual translates, weighted by inferential role. We prove that every revolution whose untranslatable cluster is non-empty and whose vocabulary contains a referring term has an index strictly between 0 and 1, and that Kuhn's own thesis of local incommensurability entails this condition. We then show that translation is what makes disagreement possible, and that theory choice is rational when a criterion shared by both communities favours the successor on evidence stated in translatable vocabulary. Part~II gives translation manuals for ten revolutions, from phlogiston to oxygen and from Galen to Harvey, with \ntermsall{} terms, \ntokensall{} token readings and \ninfall{} inferences, and checks mechanically that each revolution satisfies both conclusions.
\end{abstract}

\section*{Plan of the work}

\begin{center}\small
\begin{tabular}{@{}rlr@{}}
\toprule
Vol. & Title & Pages \\
\midrule
I & Translation and Theory Choice & 1--89 \\
II & Reference, Taxonomy and the Lexicon & 90--152 \\
III & Values, Weighting and Consensus & 153--214 \\
\bottomrule
\end{tabular}
\end{center}

\paragraph{Volume I} (the present volume) states the question, defines translation manuals and the translation index, proves that partial translatability follows from Kuhn's local thesis and holds of every revolution in a table of ten, and argues that theory choice across these revolutions was rational in a sense that does not require a neutral algorithm. Part~II of the volume contains the manuals and the computed statistics.

\paragraph{Volume II: Reference, Taxonomy and the Lexicon} takes up the semantic question in the form Kuhn gave it after 1983: that a revolution changes the structured lexicon of kind terms, and that two lexicons whose taxonomies violate his no-overlap principle cannot be mapped into one another. It (i) states the no-overlap principle precisely and proves that it constrains type-level translation only, so that token-by-token translation of the kind defined in Volume~I is untouched by it; (ii) develops the theory of reference potential for natural-kind terms, comparing Kitcher's account with the causal-descriptive accounts that descend from Putnam and with Bird's naturalized treatment of incommensurability; (iii) treats partial denotation in Field's sense as a semantic relation in its own right and gives its logic; and (iv) extends the manuals of Volume~I to four further revolutions: continental drift to plate tectonics, the classical to the molecular gene, Daltonian to structural chemistry, and Maxwellian electrodynamics to quantum electrodynamics.

\paragraph{Volume III: Values, Weighting and Consensus} takes up the methodological question. It (i) gives a formal theory of choice under shared criteria and answers Okasha's argument from Arrow's theorem by identifying the informational assumptions under which aggregation of the criteria is possible; (ii) compares the qualitative account of Volume~I with Bayesian reconstructions of Kuhnian choice in the line of Salmon; (iii) studies how communities converge, using the time from the first statement of a successor to its general acceptance in each revolution of Volumes~I and~II; and (iv) connects the norms that SCI-05 identifies as uniting the scientific methods with the institutions of error-correction that SCI-01 takes to distinguish science, and draws the consequences for the rationality of the community, as distinct from the individual.

\section*{Preface}

The thesis that scientific paradigms are incommensurable is more than sixty years old, and the arguments for and against it have reached a familiar equilibrium. Defenders point to the failure of every proposed translation of the central terms of a superseded theory; critics reply that the failure is local, that reference can be preserved where meaning is not, and that scientists in fact compare rival theories. Both sides argue from examples, and both choose the examples that suit them. What has been missing is a measure: a way of stating how much of a superseded theory can be translated into its successor, by a definite standard, across a body of cases chosen before the measuring began.

This volume supplies the measure and applies it. The written part (Part~I) gives the definitions and the arguments. Part~II gives translation manuals for ten revolutions in the history of the natural sciences and medicine. The manuals are compiled by hand, from the historical record and the secondary literature; the statistics are computed from them by a program that also checks, for every revolution, the claims that the written part makes about it. The manuals are long because the vocabularies are, and because a manual that is to be checked must state, for every term and every inference, what it claims.

\section{Introduction}\label{s:intro}

\subsection{The question}

In \emph{The Structure of Scientific Revolutions} Kuhn argued that the history of a mature science consists of long periods of normal research, governed by a paradigm, separated by revolutions in which one paradigm is replaced by another \citep{Kuhn1962}. He claimed that the paradigms on either side of a revolution are \emph{incommensurable}: there is no common measure by which they can be compared. Feyerabend, independently and in the same year, argued that the meanings of the terms of a theory are fixed by the theory as a whole, so that a change of fundamental theory changes the meaning of every term, and successive theories cannot be compared by their consequences \citep{Feyerabend1962}.

If the thesis were true, the consequences would be serious. Scientists who chose the successor could not have chosen it because it was better by any standard both sides accepted, since there would be no such standard; nor because it said more truly what the predecessor said, since the predecessor's claims would not be expressible in the successor's terms. Progress through revolutions would be, at best, progress by the lights of the winners. Lakatos expressed the worry by saying that on Kuhn's account a revolution would be a matter of mob psychology \citep{Lakatos1970}. Kuhn rejected the charge, but the debate that followed did not settle whether he was entitled to reject it.

We argue that paradigms are not incommensurable. Successive paradigms are partially translatable, and theory choice remains rational across revolutions. Both conjuncts are substantive. The first says more than that \emph{some} expressions are shared, which no one has denied, and less than that everything said by the predecessor can be said by the successor, which is false. The second says that the choice of the successor could be supported by reasons available to both communities; it does not say that the choice was dictated by an algorithm, or that every scientist who resisted the successor was irrational.

\subsection{Three theses}

\citet{HoyningenHueneSankey2001} and \citet{Sankey1994} distinguish three strands in the incommensurability literature.

The \emph{semantic} thesis says that the central terms of successive theories differ in meaning, or in reference, in a way that prevents translation of the claims of one into the language of the other. It is the strand that Feyerabend emphasized and that Kuhn, after 1970, made central \citep{Kuhn1983,Kuhn1991}.

The \emph{methodological} thesis says that the standards by which theories are evaluated change with the paradigm, so that there is no neutral standard by which successive paradigms can be compared. It is prominent in \emph{Structure}, where Kuhn writes of paradigms carrying their own standards and of the problems each regards as legitimate.

The \emph{perceptual} thesis says that scientists working in different paradigms see the world differently, so that there is no neutral observational basis for comparison. It draws on the gestalt-switch analogy of \emph{Structure} and on Hanson's account of theory-laden observation \citep{Hanson1958}.

Our first conclusion, partial translatability, answers the semantic thesis. Our second, the rationality of choice, answers the methodological thesis, and, through the requirement that the evidence for a choice be stated in translatable vocabulary, the perceptual thesis as far as it bears on choice.

\subsection{The approach}

The argument has three components.

The first is a definition of translation that is neither too strong nor too weak. Kuhn's later thesis was that translation in the strict sense, a systematic substitution of expressions preserving meaning, fails for a local cluster of terms \citep{Kuhn1983}. We define a translation manual as a partial map from expressions of the earlier theory to expressions of the later one, preserving the reference of term tokens in the sense of \citet{Kitcher1978} and a specified set of inferential relations. A manual so defined can be applied term by term and token by token; it can say, of every term, whether it is translated exactly, partially, token by token, or not at all; and it can say, of every inference of the earlier theory, whether its translation holds in the later theory, holds in a restricted domain, is translated and denied, or cannot be translated. From the manual we define a \emph{translation index}: the share of the earlier theory's vocabulary that the manual translates, weighted by the number of inferences in which each term figures.

The second is a pair of general results. The first shows that a revolution in which some readings of terms have referents in the successor and some do not has a translation index strictly between 0 and 1, and that Kuhn's own thesis of local incommensurability entails these conditions. The second shows that the choice of the successor is rational, in a sense we define, when a criterion shared by both communities favours the successor on evidence that both can state.

The third is the application of the definitions to ten revolutions, chosen before the computation for the reason that they are the revolutions most discussed in the literature on incommensurability, scientific realism and theory change: phlogiston to oxygen, Ptolemy to Kepler, Newton to Einstein, caloric to thermodynamics, classical to quantum mechanics, fixed species to Darwin, Galen to Harvey, emission to wave optics, miasma to germs, and blending inheritance to Mendel. For each we give a translation manual, compute the indices, record the outcome on Kuhn's five shared criteria, and check mechanically that the index lies strictly between 0 and 1 and that at least one shared criterion favoured the successor on evidence stated in translatable vocabulary.

\subsection{Relation to other work in the collection}

SCI-06-b argues, from Kuhn's later texts alone, that his account of local incommensurability and of the distinction between translation and interpretation already contains the premises of our conclusion. The present volume gives that argument a definite form and tests it against the record. SCI-05-a argues that the scientific methods form a family united by four norms, answerability, public checkability, proportion and severity, and that these shared norms are part of what keeps theory choice across revolutions rational; we take that result as an assumption (Assumption~\ref{a:norms}) and show where it enters. SCI-02-a defends scientific realism and answers the pessimistic meta-induction by the retention of the essential content of past theories; we show that partial translatability is the semantic condition of that retention (Proposition~\ref{p:retention}). LNG-06 and LNG-16 argue that radical translation is constrained by more than stimulus meaning and that reference is not inscrutable; we rely on their conclusions in replying to the Quinean form of the incommensurability argument (Objection~\ref{o:quine}).

\section{A history of the debate}\label{s:history}

\subsection{Kuhn and Feyerabend, 1962}

In \emph{Structure} incommensurability was introduced with three illustrations \citep[ch.~12]{Kuhn1970}. The proponents of competing paradigms disagree about the list of problems a candidate for paradigm must solve; they use familiar terms in new ways, as when \enquote{mass} acquires in relativity a meaning it did not have for Newton; and they practise their trades in different worlds, seeing different things when they look from the same point in the same direction. Kuhn drew from the illustrations the conclusion that the transition between paradigms cannot be made a step at a time, forced by logic and neutral experience, and that it must occur all at once, like a gestalt switch, or not at all.

Feyerabend's version was narrower and more radical \citep{Feyerabend1962}. He attacked the account of reduction and explanation according to which a later theory explains an earlier one by entailing it, or entailing an approximation to it. Newton's theory does not entail Galileo's law of free fall, which says that the acceleration of a falling body is constant, since on Newton's theory it increases as the body approaches the Earth. More generally, the meanings of the descriptive terms of a theory depend on the theory, so that when the fundamental theory changes, the meanings change, and the earlier and later theories do not share the terms in which entailment, or contradiction, could be stated. In \emph{Against Method} he extended the argument to observational terms and drew the conclusion that no methodological rule survives the history of science unviolated \citep{Feyerabend1975}.

\subsection{The first replies}

The first replies attacked the meaning-variance premise. \citet{Shapere1966} argued that if every change of theory changes the meanings of all its terms, then no two theories can be compared, rivals included, and the history of science, which records theories competing to explain the same phenomena, becomes unintelligible; Kuhn's own historical narratives presuppose that the competitors were addressing a common subject matter. \citet{Scheffler1967} argued that sameness of reference, not of meaning, is what comparison needs: two theories that use terms with different meanings can still conflict if the terms pick out the same things, as \enquote{the morning star} and \enquote{the evening star} do. Scheffler's point anticipated the causal theory of reference, which supplied a mechanism by which reference could be preserved through changes of belief.

\citet{Davidson1974} attacked the idea of a conceptual scheme. If incommensurability is the impossibility of translation, then to recognize another theory as incommensurable with ours we must recognize it as a language that we cannot translate. But the only evidence that something is a language is that it can be translated. A totally untranslatable scheme is therefore not one we could identify, and the thesis of total incommensurability has no content. Davidson's argument tells against total incommensurability, which few held; it leaves open partial untranslatability, which is the thesis we must address.

\subsection{Reference}

The causal theory of reference, developed by \citet{Putnam1975} among others, held that the reference of a natural-kind term is fixed not by the descriptions its users associate with it but by its causal connection to samples of the kind, through an original baptism and a chain of transmission. If this were right, the reference of \enquote{electron} or \enquote{gene} would be preserved through large changes in theory, and incommensurability of meaning would not imply incommensurability of reference. \citet{Putnam1981} drew the further conclusion that the incommensurability thesis is self-refuting: if Kuhn could not translate the phlogiston theorists, he could not tell us what they believed, yet his histories do just that.

The causal theory has a defect for the cases that matter. It assigns a referent to \enquote{phlogiston} only if something was baptized, and if nothing was, it assigns none; it cannot explain how a theorist who used \enquote{dephlogisticated air} could on some occasions have referred to oxygen, which is what the historical record suggests. Two proposals addressed the defect.

\citet{Field1973} argued that a term of a superseded theory may \emph{partially denote} several entities of the later theory, without determinately denoting any. His example was Newtonian \enquote{mass}, which partially denotes relativistic rest mass and relativistic mass (total energy divided by the square of the speed of light). Some Newtonian claims about mass are true of the one and false of the other: the claim that mass is invariant is true of rest mass, and the claim that momentum is the product of mass and velocity is true of relativistic mass. Field's semantics assigns truth to a sentence if it is true on every admissible assignment, falsity if false on every one, and neither otherwise.

\citet{Kitcher1978} argued that a term has a \emph{reference potential}: a set of ways in which its tokens can have their reference fixed. A token of \enquote{dephlogisticated air} in Priestley's description of the gas he had just prepared from the red calx of mercury refers to what he had in his apparatus, which was oxygen. A token in a statement of the theory, saying that dephlogisticated air is air from which phlogiston has been removed, refers through the description, and since nothing satisfies it, refers to nothing. The term is heterogeneous: its tokens are fixed in different ways and some succeed where others fail. Translation is then possible token by token, by substituting for each token an expression of the later theory with the same referent; what fails is a single substitution for the type. Kitcher developed the account in \emph{The Advancement of Science}, where it underwrites a notion of conceptual progress \citep{Kitcher1993}. \citet{Devitt1979} argued on similar lines that incommensurability theses confuse meaning with reference and that the comparison of theories needs only the latter.

\subsection{Kuhn's later account}

Kuhn responded by restating incommensurability as a thesis about translation and narrowing its scope \citep{Kuhn1983}. Two theories are incommensurable if there is no language into which both can be translated without residue or loss. The incommensurability is \emph{local}: most of the terms shared by the two theories function in the same way in both, and only a small group of interdefined terms resists translation. Kuhn's examples were the Newtonian cluster \enquote{mass}, \enquote{force} and \enquote{weight}, and the eighteenth-century cluster \enquote{phlogiston}, \enquote{element} and \enquote{principle}.

Against Kitcher, Kuhn argued that token-by-token reference is not translation \citep{Kuhn1983}. A translation substitutes, systematically, words for words; Kitcher's technique renders one term by different expressions on different occasions, and when it renders a passage it does not tell the reader why the author said what he said, since it obliterates the unity of the term in the author's usage. What Kitcher describes is interpretation: the historian's learning of the old language, which is possible even where translation is not. Kuhn's distinction between translation and interpretation is the hinge on which the later debate turned. If interpretation is possible, then the claims of the earlier theory can be understood, and the scientists who learned both languages could evaluate both theories, even if the evaluation could not be conducted in a single language.

In the 1990s Kuhn developed the local thesis into an account of the lexicon \citep{Kuhn1991,Kuhn1993}. Kind terms are organized into taxonomies governed by a \emph{no-overlap principle}: two kind terms of the same lexicon may not overlap in their referents unless one is subordinate to the other. A revolution changes the taxonomy so that the kinds of the new lexicon overlap those of the old without being nested in them, and no translation can respect both taxonomies. \citet{Hacking1993} called this the taxonomic solution and argued that it captures the sense in which, after a revolution, the scientist works in a new world. Volume~II takes up the taxonomic account in detail; here we note that it constrains the translation of types and leaves untouched the translation of tokens.

On the methodological side, \citet{Kuhn1977} named five characteristics of a good scientific theory: accuracy, consistency (internal and with other accepted theories), broad scope, simplicity and fruitfulness. These are, he held, shared by the scientists on both sides of a revolution and provide the shared basis for theory choice; what is not shared is a rule for applying and weighing them. Two scientists may agree on the criteria and differ on the choice, because the criteria are imprecise and because they may conflict. Kuhn took this to show that theory choice is not algorithmic; he denied that it showed choice to be irrational. \citet{McMullin1983} and \citet{Laudan1984} developed the point in different directions, Laudan arguing that the aims and methods of science, too, change, but piecemeal and by reasons, so that rational comparison is possible across changes of each.

\subsection{Reconstructions and recent work}

\citet{HoyningenHuene1993} reconstructed Kuhn's philosophy as a form of Kantianism with a dynamic, historically changing set of categories, and gave the authoritative account of the development of the incommensurability thesis. \citet{Sankey1994,Sankey1997} argued that the semantic thesis is either too strong to be true or too weak to threaten rational comparison, and that the referential approach of Scheffler, Putnam, Kitcher and Devitt shows how comparison survives meaning change. \citet{Bird2000} argued that Kuhn's views are best understood as a naturalistic account of scientific cognition, and later that incommensurability is a phenomenon of the psychology of trained perception and judgment, real but not a threat to rationality \citep{Bird2008}. \citet{AndersenBarkerChen2006} studied the cognitive structure of revolutions with frame models of concepts and argued that the incommensurability of conceptual structures is local and graded. \citet{Thagard1992} treated conceptual revolutions as changes in networks of concepts and propositions and evaluated them by explanatory coherence. \citet{Chang2012} reopened the chemical revolution, arguing that the phlogiston theory was abandoned prematurely and had resources the oxygen theory lacked, and drew a pluralist conclusion. \citet{Okasha2011} argued that Kuhn's account of theory choice as the weighing of several criteria falls under Arrow's impossibility theorem, so that no rule of choice satisfies minimal conditions of rationality.

\subsection{Research programmes and research traditions}

Two rivals of Kuhn's account bear on the question. \citet{Lakatos1970} proposed that the unit of appraisal is a research programme, a succession of theories sharing a hard core, and that a programme is progressive if its successive theories predict novel facts, some of which are corroborated, and degenerating if they only accommodate. On Lakatos's account the rationality of choice between programmes is retrospective: one programme is to be preferred when it has been progressive and its rival degenerating. \citet{Musgrave1976} applied the account to the chemical revolution and argued that the oxygen programme was progressive in the 1770s and 1780s while the phlogiston programme degenerated. Lakatos's account presupposes that the predictions of rival programmes can be compared, which is what partial translatability secures: a novel fact predicted by one programme is a fact the other can state and fail to predict. \citet{Laudan1977} replaced programmes by research traditions, evaluated by their problem-solving effectiveness, and in \citet{Laudan1984} argued that theories, methods and aims change at different times, so that each can be rationally evaluated by the others that are held fixed. Laudan's \enquote{reticulated} model is close to ours in spirit: it makes rational comparison possible across change by finding, at each moment, something shared. The present account specifies what is shared in the vocabulary: the referring readings, which are the medium in which problems and their solutions are stated.

\subsection{Recent assessments}

Recent assessments of the incommensurability debate tend to agree on its outcome while disagreeing about its lessons. \citet{Sankey1997} and \citet{HoyningenHuene1993} concur that incommensurability, in Kuhn's later sense, is local and does not entail the impossibility of rational comparison; Hoyningen-Huene emphasizes that Kuhn never held otherwise, and that the received image of Kuhn as an irrationalist is a misreading. \citet{Bird2000,Bird2008} accepts the phenomenon but relocates it in the cognitive habits of trained scientists, so that what is incommensurable is the way two communities see and judge, not the content of what they say. \citet{AndersenBarkerChen2006} offer a cognitive model on which concepts are represented by frames with attributes and values, and incommensurability is a mismatch between frames that can be partial and graded. \citet{Chang2012} accepts that the chemical revolution was a comparison of theories, but argues that the comparison was made on incomplete grounds and that the elimination of phlogiston lost a valuable line of inquiry. The present volume adds to this literature a definite measure and a body of cases; it agrees with Hoyningen-Huene and Sankey on the outcome, and with Andersen, Barker and Chen that incommensurability is graded, and it locates the grade in a quantity, the translation index, that can be computed from a manual.


\subsection{Lessons}

Each defeated proposal teaches something we must respect. The meaning-variance thesis of \citet{Feyerabend1962} fails by proving too much, as Shapere saw; a criterion of commensurability must make rival theories in the same domain comparable. Davidson's argument shows that total untranslatability is not a coherent possibility; a criterion must leave room for partial failure. The causal theory shows that reference can be preserved through changes of belief, but by itself it cannot handle heterogeneous terms; a criterion must assign reference by token. Field shows that some terms are indeterminate between several successors; a criterion must allow a term to be translated by more than one expression. Kuhn's reply to Kitcher shows that token reference is not translation of the type, and a criterion must state what kind of translation it measures and not trade on the ambiguity. Kuhn's account of values shows that shared criteria need not determine a unique choice; a criterion of rationality must not require that they do. And Okasha shows that the weighing of criteria is formally delicate; a criterion of rationality must say what it requires of the weighing.

\section{Framework}\label{s:framework}

\subsection{Theories, terms and readings}

\begin{definition}[Vocabulary, token, reading]\label{d:reading}
The \emph{vocabulary} of a theory $T$ is the set of its descriptive terms. A \emph{token} of a term is a use of it on an occasion, in a context. A \emph{reading} of a term is a class of tokens whose reference is fixed in the same way: for example, the tokens of \enquote{dephlogisticated air} in Priestley's reports of the gas obtained from the red calx of mercury form one reading, and the tokens in statements of the definition of the term form another.
\end{definition}

Readings, rather than individual tokens, are the units of the manuals in Part~II. Two tokens belong to the same reading when the same mode of reference-fixing applies to both: the same sample, the same experimental procedure, or the same description.

\begin{definition}[Reference potential]\label{d:refpot}
The \emph{reference potential} of a term is the set of its readings, each paired with the mode in which the reference of its tokens is fixed \citep{Kitcher1978}. A term is \emph{homogeneous} if all its readings have the same referent, and \emph{heterogeneous} otherwise.
\end{definition}

\begin{assumption}[Token reference]\label{a:token}
The referent of a token is fixed by the mode of its reading: by the sample or apparatus at hand when the reading is fixed by ostension or experiment, and by the satisfier of the description when it is fixed by description. Where a description is satisfied by nothing, the token refers to nothing. Where the evidence does not decide between two referents in the successor theory, the token \emph{partially denotes} both \citep{Field1973}. Referents are identified in the vocabulary of the successor theory, as since developed.
\end{assumption}

The assumption combines Kitcher's reference potential with Field's partial denotation. It is the semantic assumption on which the manuals rest, and it is the assumption Kuhn contested, not as a theory of reference but as a theory of translation. We accept the contest: what Assumption~\ref{a:token} delivers is a token-level translation, and we measure that, distinguishing it throughout from type-level translation (Definition~\ref{d:status}).

The identification of referents in the vocabulary of the successor, as since developed, is not a bias in favour of the successor. It is what the question requires. To ask whether a claim of Priestley's can be translated into Lavoisier's language is to ask whether something in Lavoisier's ontology is what Priestley was talking about; that question is answered by the best account we have of what there was in Priestley's apparatus, which is the account of the successor theory and its continuation. Where the successor's vocabulary at the time lacked a term for something its continuation names (methane, Vibrio cholerae), we use the later term and say so.

\subsection{Translation manuals}

\begin{definition}[Translation manual]\label{d:manual}
Let $T_1$ and $T_2$ be theories and let $\mathcal{I}$ be a specified set of inferences of $T_1$, each an inference that $T_1$ licenses and that is stated with readings of terms of $T_1$. A \emph{translation manual} from $T_1$ to $T_2$ with respect to $\mathcal{I}$ is a partial map $M$ from readings of terms of $T_1$ to expressions of $T_2$, together with its extension to the inferences in $\mathcal{I}$, such that:
\begin{enumerate}[label=(\roman*)]
\item \emph{(referential preservation)} for every reading $r$ in the domain of $M$, the referent of $M(r)$ in $T_2$ is the referent of $r$ under Assumption~\ref{a:token} (or, for a reading that partially denotes several entities, $M(r)$ is the list of expressions for them);
\item \emph{(domain)} a reading is in the domain of $M$ if and only if it has a referent in $T_2$;
\item \emph{(inferential extension)} an inference $I \in \mathcal{I}$ is translated, $M(I)$, if and only if every reading in it is in the domain of $M$, and $M(I)$ is the result of replacing each reading $r$ by $M(r)$ and rephrasing in the idiom of $T_2$.
\end{enumerate}
A manual has a \emph{type-level image} for a term $t$ if it maps all readings of $t$ to one expression of $T_2$ (or one list, under partial denotation).
\end{definition}

A manual is a partial map. It is not required to translate everything, and the measure of its partiality is the object of study. It preserves two kinds of relation. Reference is preserved by condition~(i), which is a constraint on the map. Inferential relations are not guaranteed to be preserved; whether they are is a fact about $T_2$, recorded by the verdicts of the next definition.

\begin{definition}[Verdicts]\label{d:verdict}
For each inference $I \in \mathcal{I}$, exactly one of the following holds.
\begin{description}
\item[H] $I$ is translated and $M(I)$ holds in $T_2$.
\item[R] $I$ is translated and $M(I)$ holds in $T_2$ only within a restricted domain of parameter values (such as $v \ll c$, or $h\nu \ll kT$) or on one of the readings of a partially denoting term.
\item[X] $I$ is translated and $T_2$ denies $M(I)$.
\item[U] $I$ is not translated: some reading in it has no referent in $T_2$.
\end{description}
An inference is \emph{preserved} if its verdict is H or R, \emph{comparable} if it is H, R or X.
\end{definition}

The verdict X is the one on which the argument of Section~\ref{s:rational} turns. An inference with verdict X is a claim of $T_1$ that $T_2$ can state and denies. It is a disagreement. Where the verdict is U there is no disagreement, since the successor cannot state what the predecessor claimed; there is only abandonment.

\begin{definition}[Translation status]\label{d:status}
A term $t$ of $T_1$ is translated
\begin{description}
\item[exactly] if $M$ has a type-level image for $t$ consisting of a single expression, and every inference of $\mathcal{I}$ in which $t$ figures has verdict H;
\item[partially] if $M$ has a type-level image for $t$ and either the image is a list (partial denotation) or some inference in which $t$ figures has verdict R, X or U;
\item[token by token] (\emph{token-dependent}) if $M$ has no type-level image for $t$ but some reading of $t$ is in the domain of $M$;
\item[not at all] (\emph{none}) if no reading of $t$ is in the domain of $M$.
\end{description}
\end{definition}

A term translated partially may still be translated, at the type level, with preservation of reference. What fails is the preservation of some of its inferential role. \enquote{Calx} is translated by \enquote{metallic oxide}, and every token of \enquote{calx} refers to an oxide; but the phlogistic inference that a metal is its calx plus phlogiston cannot be translated, so the inferential role of \enquote{calx} is not wholly preserved. The four statuses distinguish what Kuhn ran together under \enquote{untranslatable}: a term may resist translation of its type while every one of its tokens refers (\enquote{inflammable air}, entry CH04), or some of its tokens refer and others do not (\enquote{dephlogisticated air}, CH02), or none does (\enquote{primum mobile}, AS03).

\begin{definition}[Weight and score]\label{d:score}
The \emph{inferential weight} $w(t)$ of a term $t$ is the number of inferences of $\mathcal{I}$ in which some reading of $t$ figures. The \emph{translation score} $s(t)$ is 1 if $t$ is translated exactly or partially, $n_{\mathrm{ref}}(t)/n(t)$ if it is translated token by token, where $n(t)$ is the number of its readings and $n_{\mathrm{ref}}(t)$ the number of these with a referent in $T_2$, and 0 if it is not translated.
\end{definition}

\begin{definition}[Indices]\label{d:index}
The \emph{translation index} of a manual $M$ is
\[
\tau(M) \;=\; \frac{\sum_{t} w(t)\, s(t)}{\sum_t w(t)},
\]
the sum running over the vocabulary of $T_1$ covered by $M$. The \emph{comparability index} $\kappa(M)$ is the share of inferences in $\mathcal{I}$ that are comparable, and the \emph{preservation index} $\pi(M)$ the share that are preserved.
\end{definition}

The translation index is the share of the vocabulary with full or partial translation, weighted by inferential role. Terms that figure in many inferences count for more, which reflects the holism Kuhn and Feyerabend emphasized: a term central to the theory's inferences matters more to its translatability than a peripheral one. The index counts a token-dependent term in proportion to its referring readings, so that it measures the token-level translation Kitcher defended without pretending that it is type-level translation. Alternative scorings are considered in Proposition~\ref{p:robust}.

\begin{definition}[Partial translatability]\label{d:partial}
$T_1$ is \emph{partially translatable} into $T_2$ with respect to $\mathcal{I}$ if there is a manual $M$ from $T_1$ to $T_2$ with $0 < \tau(M) < 1$.
\end{definition}

\begin{definition}[Untranslatable cluster]\label{d:cluster}
The \emph{untranslatable cluster} $C(M)$ of a manual is the set of terms of $T_1$ at least one of whose readings has no referent in $T_2$.
\end{definition}

The cluster is the precise counterpart of the local group of interdefined terms of \citet{Kuhn1983}. Kuhn's thesis of local incommensurability is that, for every revolution, the cluster is non-empty and small. We test the second part, which is not needed for our result, as well as the first.

\subsection{Why this definition of translation}

The term \enquote{translation manual} is Quine's \citep{Quine1960}, and the choice is deliberate: a manual in our sense, like Quine's, is a systematic correlation of expressions of one language with expressions of another, judged by its fit with the evidence about how the expressions are used. It differs from Quine's in three respects. It is partial, where Quine's manuals were total; it is constrained by reference, fixed under Assumption~\ref{a:token}, where Quine's were constrained only by dispositions to assent and dissent; and it is relativized to a specified set of inferences, whose preservation it records rather than requires.

Each difference answers a lesson of Section~\ref{s:history}. Partiality is required by Davidson's argument, which shows that total failure is not an intelligible possibility, and by Kuhn's local thesis, which says that partial failure is actual. The referential constraint is required by the cases of heterogeneous terms, which no assignment of a single meaning to a type can handle. And the relativization to inferences is required by the holism that Kuhn and Feyerabend emphasized: whether a term \enquote{means the same} in two theories is, on any view that takes holism seriously, a matter of how much of its inferential role survives, and the verdicts record exactly that.

Two stronger definitions were available and rejected. On the first, a translation must preserve meaning, in the sense of preserving every inference in which a term figures. On that definition, almost nothing is translatable across a revolution: of the \ntermsall{} terms of Part~II, only the \nstatE{} translated exactly would count, and of these several (\enquote{fixed air}, \enquote{nitrous air}) would fail if the inference set were extended to the theory of their composition. The definition makes incommensurability trivially true and uninteresting, since it makes every pair of theories that differ in any inference incommensurable, rivals in normal science included; this was Shapere's objection to Feyerabend. On the second, a translation must be a type-level map, but need not preserve inferences. That is the definition Kuhn used against Kitcher, and we report the index it yields as scheme S1 (Proposition~\ref{p:robust}). It too gives indices strictly between 0 and 1 in every case. The conclusion of partial translatability therefore does not depend on choosing the weakest of the available definitions.

\subsection{Criteria, norms and rational choice}

\begin{assumption}[Shared criteria]\label{a:criteria}
The scientific communities on both sides of each revolution accept accuracy, consistency, scope, simplicity and fruitfulness as criteria for the evaluation of theories \citep{Kuhn1977}, though they may differ in how they apply and weigh them.
\end{assumption}

\begin{assumption}[Shared norms]\label{a:norms}
The communities on both sides of each revolution accept the norms of answerability to evidence, public checkability, proportioning confidence to evidence and severity of testing, which unite the scientific methods (SCI-05-a).
\end{assumption}

The two assumptions do different work. Assumption~\ref{a:criteria} supplies the dimensions on which theories are compared. Assumption~\ref{a:norms} constrains how the comparison may be conducted: the evidence on which a criterion is applied must be public and checkable, the confidence placed in the outcome must be proportioned to the evidence, and the outcome must be answerable to evidence that could have gone the other way. Kuhn's criteria without the norms could be applied by private judgment, and a choice made by such judgment would be rational only by courtesy; the norms are what make the application of the criteria a matter of reasons that the other side can check. This is the sense in which, as SCI-05-a holds, the shared norms are part of what keeps theory choice across revolutions rational.

\begin{definition}[Outcome]\label{d:favour}
A criterion \emph{favoured} $T_2$ over $T_1$ in a revolution if the historical record documents that, on the evidence available to the communities at the time of the choice, $T_2$ did better than $T_1$ by that criterion. It favoured $T_1$ if the record documents the converse, and the outcome was \emph{mixed} if it documents respects in which each did better.
\end{definition}

\begin{definition}[Evidential base]\label{d:evbase}
An \emph{evidential claim} for an outcome is a claim about observations or experiments on which the outcome rests. It is \emph{stated in translatable readings} if every reading of a term of $T_1$ that it uses has a referent in $T_2$.
\end{definition}

\begin{principle}[Disagreement presupposes translation]\label{p:disagree}
A claim of $T_1$ conflicts with $T_2$ only if it can be stated in the vocabulary of $T_2$ and $T_2$ denies it.
\end{principle}

\begin{principle}[Interpretation]\label{p:interp}
A reading with no referent in $T_2$ can be learned by a speaker of $T_2$, who can then understand the claims of $T_1$ that contain it \citep{Kuhn1983}.
\end{principle}

\begin{principle}[Shared reasons]\label{p:reasons}
The choice of one theory over another is supported by a reason available to both communities if a criterion that both accept favours the chosen theory on an evidential base that both can state and check.
\end{principle}

Principle~\ref{p:disagree} is a logical point: to deny what another says, one must be able to say it. It is the premise of the arguments of \citet{Davidson1974} and \citet{Putnam1981}, and Kuhn accepted it, which is why he insisted that incommensurability is not incompatibility. Principle~\ref{p:interp} is Kuhn's own; we use it only to note that the untranslatable remainder is not unintelligible. Principle~\ref{p:reasons} states a sufficient condition for a reason to be shared; it does not say that only such reasons count.

\begin{definition}[Rational choice across a revolution]\label{d:rational}
The choice of $T_2$ over $T_1$ is \emph{rational across the revolution} if (i) some criterion of Assumption~\ref{a:criteria} favoured $T_2$, (ii) an evidential base for that outcome is stated in translatable readings, so that the outcome rests on claims that both communities could state, and (iii) the comparison was conducted under the norms of Assumption~\ref{a:norms}. The support is \emph{dominant} if, in addition, no criterion favoured $T_1$.
\end{definition}

The definition is deliberately modest. It does not require that the choice be dictated by the criteria, that every scientist make it, or that it be made at once. It requires that there be a reason for it that a member of the losing community could recognize as a reason by his own standards, stated in terms he could use. That is what the methodological incommensurability thesis denies, and what we assert.

\section{Partial translatability}\label{s:partial}

\begin{lemma}[Bounds]\label{l:bounds}
For any manual $M$: $\tau(M) > 0$ if and only if some term of $T_1$ has a reading with a referent in $T_2$; and $\tau(M) < 1$ if and only if the untranslatable cluster $C(M)$ is non-empty.
\end{lemma}

\begin{proof}
Every term covered by $M$ figures in at least one inference, so every weight is positive. The index is a weighted mean of scores in $[0,1]$ with positive weights, so it is positive if and only if some score is positive, and less than 1 if and only if some score is less than 1. A score is positive if and only if the term is translated exactly or partially (score 1) or token by token (score $n_{\mathrm{ref}}/n > 0$); in each case some reading has a referent, and conversely a term with a referring reading has one of these statuses by Definition~\ref{d:status}. A score is less than 1 if and only if the term is translated token by token with some non-referring reading, or not at all. By Definition~\ref{d:manual}(ii) and Definition~\ref{d:status}, a term with a type-level image has all its readings in the domain of $M$, so these are exactly the terms with a non-referring reading, that is, the members of $C(M)$.
\end{proof}

\begin{lemma}[Disagreement]\label{l:disagree}
Every inference of $T_1$ that $T_2$ denies is translated by $M$. Hence, if $T_2$ denies any inference of $T_1$, then $\kappa(M) > 0$.
\end{lemma}

\begin{proof}
By Principle~\ref{p:disagree}, $T_2$ denies a claim only if it can state it, that is, only if every reading in it has a referent in $T_2$. By Definition~\ref{d:manual}(iii) such an inference is translated, and its verdict is X, which is comparable.
\end{proof}

The lemma turns the usual argument around. The incommensurability theorist points to the depth of the disagreement between the parties to a revolution as evidence that they cannot understand each other. But the depth of a disagreement is measured by what is denied, and what is denied must be stated. The inference CH.I19, that water is an element, was denied by the oxygen chemists; the denial was a disagreement about water, which both sides could identify, and about the predicate \enquote{simple substance}, which both sides could apply. By contrast the inference CH.I01, that combustion releases phlogiston, was not denied in the same sense; it was dropped, since the successor had nothing to say about phlogiston.

\begin{proposition}[Kuhn's local thesis]\label{p:local}
If incommensurability across a revolution is local in Kuhn's sense, that is, if there is a non-empty group of terms of $T_1$ some of whose uses cannot be rendered in $T_2$, while other terms of $T_1$ function in $T_2$ as they did in $T_1$, then every manual from $T_1$ to $T_2$ that covers both groups has $0 < \tau < 1$.
\end{proposition}

\begin{proof}
A term some of whose uses cannot be rendered in $T_2$ has, under Assumption~\ref{a:token}, some reading with no referent in $T_2$: if every reading had a referent, each use could be rendered by an expression for that referent, which is what a token-level rendering is. So $C(M) \neq \emptyset$. A term that functions in $T_2$ as it did in $T_1$ has its readings referring to what the corresponding term of $T_2$ refers to. By Lemma~\ref{l:bounds}, $0<\tau<1$.
\end{proof}

Proposition~\ref{p:local} is the analytic core of the first conclusion. Kuhn's later account says exactly what its antecedent says: incommensurability is confined to a small group of interdefined terms and most terms function the same way in both theories \citep{Kuhn1983}. The proposition shows that this account, taken at its word, entails partial translatability in the sense of Definition~\ref{d:partial}. Kuhn would have objected that what $\tau$ measures is not translation; we take up the objection in Section~\ref{s:objections} (Objection~\ref{o:tokens}).

\begin{theorem}[Partial translatability]\label{t:partial}
Let $R$ be a revolution and $M$ a manual for it satisfying Definition~\ref{d:manual} under Assumption~\ref{a:token}. If $C(M)$ is non-empty and some term of $T_1$ has a reading with a referent in $T_2$, then $T_1$ is partially translatable into $T_2$. In particular, for each of the \nrevs{} revolutions of Part~II, $0 < \tau < 1$; the indices range from \timin{} to \timax{}.
\end{theorem}

\begin{proof}
The general claim is Lemma~\ref{l:bounds} with Definition~\ref{d:partial}. For the revolutions of Part~II, the program computes $C(M)$ and the scores from the manuals and verifies that the cluster is non-empty and that some term has a referring reading (check C6, and check C4 for the inequalities); Table~\ref{tab:stats} gives the values.
\end{proof}

The first conclusion of the volume, that successive paradigms are partially translatable, is the general claim that the hypothesis of Theorem~\ref{t:partial} holds of every revolution. It is supported in two ways. The first is Proposition~\ref{p:local}: anyone who accepts Kuhn's later account must accept it. The second is inductive: the ten revolutions of Part~II include the cases on which the incommensurability thesis was founded (the chemical, Copernican, relativistic and quantum revolutions), and the hypothesis holds of each. The hypothesis has two parts. That the cluster is non-empty is the part of incommensurability that is true. That some readings refer is guaranteed for any successor that replaces its predecessor in the same domain: the two theories must at least share the means of identifying the phenomena both claim to explain, and the terms for these refer in both. Proposition~\ref{p:nototal} makes this exact.

\begin{corollary}[Locality]\label{c:local}
In each revolution of Part~II the cluster is a proper part of the vocabulary. Over all ten, \nclusterall{} of \ntermsall{} terms belong to clusters; the largest share in any one revolution is that of MIAS, with 7 of 21 terms.
\end{corollary}

\begin{proof}
Computed (check C6 and Table~\ref{tab:further}).
\end{proof}

The corollary confirms the second part of Kuhn's local thesis, which our argument does not need: the cluster is small. It also corrects Kuhn's examples in one respect. In the chemical revolution the cluster computed from the manual is \{phlogiston, dephlogisticated air, phlogisticated air\}, not \{phlogiston, element, principle\}. \enquote{Element} is translated, partially, by Lavoisier's \enquote{simple substance}: the extension changed (water and air ceased to be elements), but tokens of \enquote{element} applied to particular substances are tokens of a predicate that the successor could state, and the phlogistic claim that water is an element (CH.I19) is a disagreement, not an untranslatability. Kuhn's own diagnosis confused the two, which Lemma~\ref{l:disagree} separates. In the relativistic revolution the computed cluster includes \enquote{force} and \enquote{weight}, as Kuhn said, but not \enquote{mass}, which partially denotes two successors and is translated partially (entry ME01).

\begin{proposition}[No total incommensurability]\label{p:nototal}
If $T_2$ supersedes $T_1$ by explaining phenomena that $T_1$ also explained, then some term of $T_1$ has a reading with a referent in $T_2$, and $\tau(M) > 0$ for every manual.
\end{proposition}

\begin{proof}
For $T_2$ to explain a phenomenon that $T_1$ explained, it must be the same phenomenon. So there is a description of the phenomenon, available to the community of $T_1$, which picks out something that $T_2$ also explains. The terms of that description, in the readings in which they are used to identify the phenomenon, refer to what $T_2$ refers to. By Lemma~\ref{l:bounds}, $\tau>0$.
\end{proof}

This is Shapere's point, made precise. A theory that shared no referring vocabulary with its predecessor would not be its successor, but a theory of something else.

\begin{proposition}[Robustness]\label{p:robust}
For each revolution of Part~II, $0<\tau<1$ also under the scoring schemes S1 (token-dependent terms scored 0, so that only type-level translation counts), S3 (partially translated terms scored $1/2$) and S4 (all weights equal to 1). The values range from \sensmin{} to \sensmax{}.
\end{proposition}

\begin{proof}
Computed (check C8, Table~\ref{tab:sens}).
\end{proof}

The robustness matters for the objection that the index is an artefact of our choices. S1 is the scheme that Kuhn's criticism of Kitcher recommends: count a term as translated only if its type is. Under S1 the indices fall, most for ASTR (to 0.560) and MIAS (to 0.556), whose heterogeneous terms (\enquote{planet}, \enquote{deferent}, \enquote{epicycle}; \enquote{miasma}, \enquote{contagium}, \enquote{ferment}) are many; but they stay well above 0, because most of the vocabulary of every revolution is translated at the type level. S3 penalizes partial translation; the indices fall to between \sensmin{} and 0.667 but none reaches 0 or 1.

\begin{remark}\label{r:s2}
Under the most generous scheme, S2, which scores a token-dependent term 1 if any of its readings refers, the index of CHEM and of QUAN is exactly 1. Every term of phlogistic chemistry in the manual has at least one referring reading: even \enquote{phlogiston} refers, in Kirwan's identification of it with inflammable air, to hydrogen (CH01c). The scheme therefore loses the information that matters, namely that some \emph{readings} of the term have no referent. This is why Definition~\ref{d:score} scores token-dependent terms by the share of referring readings, and why the untranslatable cluster is defined by readings. Untranslatability, the computation shows, is a property of readings, not of terms.
\end{remark}

\subsection{A worked computation}

To show how the index is computed and what it is sensitive to, we compute it for CHEM by hand from the entries of Part~II. The manual has 23 terms. Their weights, the numbers of inferences in which they figure, sum to 81: \enquote{phlogiston} (CH01) and \enquote{combustion} (CH08) have weight 10 each, \enquote{metal} (CH07) 7, \enquote{calx} (CH06) and \enquote{common air} (CH11) 6 each, \enquote{inflammable air} (CH04) and \enquote{charcoal} (CH20) 5 each, \enquote{dephlogisticated air} (CH02) 4, and the remaining fifteen terms between 1 and 3. Nineteen terms are translated exactly or partially and score 1. The four token-dependent terms score by the share of their referring readings: \enquote{phlogiston} has four readings of which one (CH01c, Kirwan's identification with inflammable air) refers, and scores $1/4$; \enquote{dephlogisticated air} scores $2/3$; \enquote{phlogisticated air} $1/2$; \enquote{inflammable air}, all of whose three readings refer, to hydrogen, methane and carbon monoxide, scores 1. The weighted shortfall from 1 is therefore
\[
10\cdot\tfrac34 + 4\cdot\tfrac13 + 2\cdot\tfrac12 \;=\; 7.5 + 1.333 + 1 \;=\; 9.833,
\]
and $\tau = (81 - 9.833)/81 = 0.879$, the value in Table~\ref{tab:stats}. Three features of the computation deserve comment. First, three quarters of the shortfall comes from \enquote{phlogiston} alone, because it is both the most heterogeneous and among the most heavily weighted terms: the weighting by inferential role does what it was designed to do, which is to make the central term of the cluster count for more than a peripheral one. Secondly, if \enquote{phlogiston} had no referring reading at all, if Kirwan's identification were set aside, its score would be 0 and $\tau$ would fall to $0.848$; the conclusion would be unchanged. Thirdly, the term \enquote{combustion}, which has the same weight as \enquote{phlogiston}, contributes nothing to the shortfall, because every one of its readings refers: the phlogistonists and the oxygen chemists burned the same candles.

\subsection{Revolutionary and cumulative change}

The index also yields a characterization of revolutionary change. Kuhn distinguished revolutions from the cumulative growth of normal science, in which new results are added without revising the vocabulary of the old \citep{Kuhn1970}. The distinction has been criticized as too sharp, and Kuhn himself later allowed that there are small revolutions, affecting a specialty, as well as large ones.

\begin{proposition}[Cumulative and revolutionary change]\label{p:cumulative}
Let $T_2$ succeed $T_1$ in a domain, in the sense of Proposition~\ref{p:nototal}. Then exactly one of the following holds for any manual $M$ that covers the vocabulary of $T_1$: $\tau(M) = 1$ and the cluster is empty, in which case every reading of every term of $T_1$ refers in $T_2$; or $0<\tau(M)<1$ and the cluster is non-empty.
\end{proposition}

\begin{proof}
By Proposition~\ref{p:nototal}, $\tau(M) > 0$. By Lemma~\ref{l:bounds}, $\tau(M)<1$ if and only if the cluster is non-empty, and $\tau(M)=1$ otherwise; when the cluster is empty, no term has a non-referring reading by Definition~\ref{d:cluster}.
\end{proof}

We may call a succession \emph{cumulative} if its cluster is empty and \emph{revolutionary} if it is not. On this definition, the defining mark of a revolution is not the falsification of the predecessor, which happens in cumulative change too (an X verdict is compatible with an empty cluster), but the loss of reference by some reading of its vocabulary: the discovery that something the predecessor talked about does not exist. This captures what is right in Kuhn's idea that revolutions change the world the scientist works in. After the chemical revolution there was no phlogiston, after the relativistic revolution no absolute space, after the Harveian revolution no vital spirit. The size of a revolution can then be measured by the size and weight of its cluster, and the distinction between small and large revolutions becomes a matter of degree, as Kuhn later held it to be. What the characterization rules out is a revolution with $\tau = 0$: a change in which nothing the predecessor talked about survives would not be a succession in a domain at all.

The characterization also explains why revolutions are recognized retrospectively. Whether a reading refers is settled by the successor, so that the cluster of a change can be computed only when the successor is established. While a revolution is in progress, the parties disagree about which readings refer, which is one reason, beyond the weighing of criteria, why the choice is contested at the time and clear in retrospect.

\section{Interpretation and the untranslatable remainder}\label{s:interp}

Theorem~\ref{t:partial} says that the untranslatable remainder is non-empty. What is its status? Kuhn's answer was that it is accessible to interpretation: a historian, or a scientist trained in the successor, can learn the old language and understand what was said in it, though he cannot say it in his own \citep{Kuhn1983}. We make the answer precise and examine what it implies for comparison.

\begin{definition}[Interpretive extension]\label{d:interp}
An \emph{interpretive extension} of $T_2$ for a manual $M$ from $T_1$ is the language obtained by adding to the vocabulary of $T_2$ every term of the cluster $C(M)$, with its non-referring readings governed by the inferences of $T_1$ in which they figure. A sentence of $T_1$ is \emph{interpretable} in $T_2$ if it is a sentence of an interpretive extension.
\end{definition}

Every sentence of $T_1$ is interpretable in $T_2$ by construction: the extension simply adds what was missing. This is the formal content of Principle~\ref{p:interp}. What the extension does not do is assign referents to the added readings; they remain, from the standpoint of $T_2$, terms that purport to refer and do not. A chemist who has learned the phlogistic language can say, in his extended language, that phlogiston is released in combustion; he cannot believe it, because he believes that there is no phlogiston, and so he takes the sentence to be either false or without truth value, according as he treats empty terms.

\begin{proposition}[Interpretation and assessment]\label{p:assess}
Every untranslatable inference of $T_1$ is interpretable in $T_2$, and the speaker of the interpretive extension who accepts $T_2$ assigns it a definite status: he rejects it, or treats it as resting on a false presupposition of existence.
\end{proposition}

\begin{proof}
The first claim is immediate from Definition~\ref{d:interp}. For the second, an untranslatable inference contains a reading with no referent in $T_2$, that is, a reading whose mode of reference-fixing picks out nothing by the lights of $T_2$. By Assumption~\ref{a:token}, a speaker who accepts $T_2$ holds that nothing satisfies the reading's description, or that the sample it was fixed by does not exist. He therefore holds the existential presupposition of the inference false, and on either standard treatment of empty terms he does not accept the inference.
\end{proof}

The proposition shows that untranslatability does not prevent assessment; it changes the form of the assessment. A translated disagreement (verdict X) is a disagreement over a claim both sides can state: the oxygen chemist says that water is not an element, and the phlogiston chemist understands him. An untranslatable claim is assessed by the denial of its presupposition: the oxygen chemist says that there is no phlogiston, and the phlogiston chemist understands that too. What the two cannot do is conduct the argument over the presupposition in the successor's vocabulary alone, since that vocabulary has no term for what is denied to exist. They conduct it in the interpretive extension, which is to say in a language that both sides of the revolution came, by the time it ended, to share.

This is why the record of revolutions is not a record of incomprehension. The participants were bilingual, or became so. Lavoisier had been trained in the phlogistic chemistry; Kirwan, the last major defender of phlogiston, published his \emph{Essay on Phlogiston} in 1787, and its French translation of 1788 appeared with a commentary by Lavoisier and his colleagues refuting it section by section; Kirwan answered the notes in an English second edition of 1789 and later conceded. Darwin had read Paley closely enough to be able to say what each of Paley's contrivances would be on his own theory. Harvey's teacher was Fabricius, whose account of the venous valves Harvey reinterpreted. In each case the successor's advocate could state the predecessor's claims, in the predecessor's terms, and say what was wrong with them.

\begin{remark}\label{r:presupposition}
The distinction between translated denial and presupposition failure corresponds to the distinction, within the cluster, between the readings of a heterogeneous term. When Priestley said that dephlogisticated air supports combustion better than common air (CH.I11), the oxygen chemist translated and agreed. When he said that dephlogisticated air is air deprived of phlogiston (CH.I12), the oxygen chemist could only interpret and deny the presupposition. The same term, in the same author, was the vehicle of both claims. The heterogeneity that Kitcher identified is thus the reason why a revolution can be at once a disagreement over shared claims and the abandonment of a theory.
\end{remark}

\section{The rationality of theory choice}\label{s:rational}

\begin{lemma}[Evidence]\label{l:evidence}
For each revolution of Part~II and each criterion recorded as favouring the successor, there is an evidential claim for that outcome stated in translatable readings.
\end{lemma}

\begin{proof}
Computed (check C9). The \nevidall{} evidential claims are listed in Table~\ref{tab:evid}, with the readings they use; the program verifies that every reading used has a referent in the successor and that every favourable outcome has at least one claim.
\end{proof}

The lemma is the heart of the answer to the methodological and perceptual theses. The evidence on which the successor won was not stated in the successor's untranslatable vocabulary, which would beg the question, nor in the predecessor's, which the successor could not assess. It was stated in vocabulary that both could use: the weight of a metal and of its calx (CHEM.E1), the position of Mars (ASTR.E1), the speed of light in water (OPTI.E4), the number of beats of the heart in half an hour (PHYS.E1). That the evidence was so stated is not surprising; evidence that is to persuade the other side must be stated in terms the other side understands, and scientists who are trying to persuade know this.

\begin{proposition}[Shared criteria are applied under shared norms]\label{p:shared}
Under Assumptions~\ref{a:criteria} and~\ref{a:norms}, if a criterion favoured $T_2$ on an evidential base stated in translatable readings, then the outcome is one that the community of $T_1$ could check by its own standards.
\end{proposition}

\begin{proof}
By Assumption~\ref{a:criteria} the criterion is one the community of $T_1$ accepts. By Definition~\ref{d:evbase} the evidential claims use only readings with referents in $T_2$; by Definition~\ref{d:manual} these readings belong to the vocabulary of $T_1$, which its community can use. By the norm of public checkability in Assumption~\ref{a:norms}, the evidence and the reasoning from it were made available in a form that others could check; by the norm of answerability, the evidence was such as could have told against $T_2$. So the community of $T_1$ could repeat the observations, state their results and apply the criterion to them.
\end{proof}

\begin{theorem}[Rationality]\label{t:rational}
For each of the \nrevs{} revolutions of Part~II, the choice of the successor was rational across the revolution in the sense of Definition~\ref{d:rational}. In \ndominant{} of them the support was dominant.
\end{theorem}

\begin{proof}
Clause~(i) of Definition~\ref{d:rational}: in each revolution at least one criterion is recorded as favouring the successor (check C5, Table~\ref{tab:crit}). Clause~(ii): Lemma~\ref{l:evidence}. Clause~(iii): Assumption~\ref{a:norms}, with Proposition~\ref{p:shared}. Dominance: in all revolutions except ASTR and BIOL no criterion is recorded as favouring the predecessor (Table~\ref{tab:crit}).
\end{proof}

The second conclusion, that theory choice remains rational across revolutions, is the general claim that the conditions of Definition~\ref{d:rational} hold in every revolution. As with the first, the support is partly analytic and partly inductive. The analytic part is that the conditions are of a kind that revolutions tend to satisfy: a successor that did better by no criterion its rivals accepted, on no evidence they could state, would not win converts among them, and the history of science records revolutions in which the converts came, over a generation, from the old community as well as from the young. The inductive part is Theorem~\ref{t:rational}.

\subsection{How the outcomes on the criteria were assigned}

The outcomes of Table~\ref{tab:crit} are historical judgments, and their reliability matters to Theorem~\ref{t:rational}. We assigned an outcome only where three conditions held: the outcome is attested in the standard histories cited in Part~II and is not, to our knowledge, disputed in them; it rests on an identifiable episode, recorded in the note for the entry, in which evidence bearing on the criterion was produced and discussed by members of both communities; and the evidence can be stated, as Table~\ref{tab:evid} states it, in translatable readings. Where any condition failed we recorded \enquote{?} and made no claim. As a result, simplicity is recorded for only three revolutions, since judgments of simplicity in the record are rarely attested without dispute, and consistency for five.

The conservative policy has a cost and a benefit. The cost is that the table understates the reasons available: no doubt Harvey's circulation was found simpler than the Galenic system by some, and the kinetic theory more consistent with mechanics than the caloric theory. The benefit is that Theorem~\ref{t:rational} does not depend on contested judgments. Check C5 requires only one favourable outcome per revolution, and every revolution has at least two, accuracy or scope together with fruitfulness; no single disputed judgment could overturn the theorem for any revolution.

It should be noted that fruitfulness, as Kuhn defined it, is a criterion whose outcome is often known only after the choice: a theory is fruitful if it discloses new phenomena or new relations among known ones. Where the fruitfulness of the successor was manifest only later, it cannot have been a reason for the original choice. We have recorded fruitfulness for all ten revolutions, but Theorem~\ref{t:rational} would hold in each without it, since in each revolution accuracy or scope was also recorded as favouring the successor on evidence available at the time of the choice.

\begin{proposition}[Rational dissent]\label{p:dissent}
Where a criterion favoured the predecessor, adherence to the predecessor for a time was also rational in the sense of Definition~\ref{d:rational}, with the roles of the theories reversed.
\end{proposition}

\begin{proof}
In ASTR, consistency with the accepted physics favoured the geocentric system until the new mechanics was available \citep{Kuhn1957,Kuhn1977}; the evidence (the absence of observed stellar parallax, the behaviour of falling bodies on a supposedly moving Earth) was stated in vocabulary both sides shared. In BIOL, consistency with Kelvin's estimates of the age of the Earth and with Jenkin's argument from blending favoured the fixity of species, or at least told against natural selection as Darwin stated it \citep{Ruse1979}; the evidence was again stated in shared vocabulary. The clauses of Definition~\ref{d:rational} hold with $T_1$ and $T_2$ exchanged.
\end{proof}

Proposition~\ref{p:dissent} expresses what Kuhn meant in saying that the criteria function as values rather than rules \citep{Kuhn1977}: they can favour different theories, and scientists who weigh them differently can make different choices, each with reasons. It is not an embarrassment for the rationality of science that Tycho had reasons to reject the motion of the Earth, or that Darwin's critics had reasons to doubt natural selection. Kuhn held that this distribution of judgment is good for science, since it spreads the risk of a community's commitment \citep{Kuhn1977}. What Definition~\ref{d:rational} adds is that the reasons on each side were stated in the shared, translatable vocabulary, and so were reasons the other side could weigh. And the dissent was resolved, in both cases, by evidence of the same kind: by the new mechanics, which removed the inconsistency, and by the extension of the age of the Earth and Mendelian inheritance, which removed the objections to selection.

\begin{proposition}[Translation and retention]\label{p:retention}
If $S$ retains the essential content of $P$ in the sense of SCI-02-a, then every essential posit of $P$ is stated in readings with referents in $S$. Hence the untranslatable cluster of a revolution contains no term all of whose readings are needed to state the essential posits of the predecessor.
\end{proposition}

\begin{proof}
Retention, as SCI-02-a defines it, requires that for each essential posit $c$ of $P$ there be an essential posit $d$ of $S$ whose truth implies the truth of $c$. For the truth of $d$ to imply the truth of $c$, $c$ must be a claim whose truth conditions $S$ can state; so $c$ must be translatable into the vocabulary of $S$, and every reading in it must have a referent in $S$. A term all of whose readings would be needed to state the essential posits would then have all its readings referring and would not belong to the cluster.
\end{proof}

The proposition shows how the present results fit the realism of SCI-02. Partial translatability is the semantic face of the retention of essential content: what is retained is stated in the translated part of the vocabulary, and what is lost is the untranslatable cluster, whose readings, when they figure in the predecessor's successes at all, figure idly. The manuals bear this out. In HEAT, the inferences on which Laplace's correction of the speed of sound depended (HE.I09, HE.I10, HE.I27) are all preserved, and the untranslatable readings of \enquote{caloric} (HE01a, HE01c) figure only in inferences about the nature of heat that the derivation did not use, in agreement with the analysis of the caloric case in SCI-02-a. In MECH the predictions that made Newton's theory mature, the return of comets and the perturbations of planets, rest on inferences recorded R, preserved in the weak-field limit (ME.I08, ME.I21). SCI-02-a counts the phlogiston theory as borderline in maturity; on either verdict the manual is consistent with it, since the inferences recording the reversibility of calcination and reduction (CH.I05--CH.I07, CH.I10) are preserved.

\section{Choice, consensus and the community}\label{s:community}

Theorem~\ref{t:rational} concerns the existence of shared reasons. A complete account of theory choice across revolutions must also say how the reasons of individuals become the choice of a community, and how long it takes. We do not give that account here; Volume~III does. But the present results constrain it, and we record the constraints.

\subsection{Individual and collective rationality}

Kuhn insisted that the criteria do not determine the choice of an individual scientist, since individuals weigh them differently, and that the community nonetheless converges \citep{Kuhn1977}. The convergence is not explained by an algorithm that every individual applies, but by the accumulation of evidence that bears on several criteria at once, so that individuals with different weightings come, one by one, to the same choice. Definition~\ref{d:rational} fits this picture. It does not say that every individual who chose the successor did so for the same reason, or that every individual who resisted was irrational. It says that there was, at the time, a reason for the choice that both communities could recognize. As the evidence accumulates, the reasons multiply: in MECH, the Michelson--Morley experiment (MECH.E3), the velocity-dependence of electron mass (MECH.E4), the perihelion of Mercury (MECH.E1) and the deflection of light (MECH.E2) bear on different criteria and were obtained over three decades.

The constraint that Theorem~\ref{t:rational} places on a theory of consensus is that the evidence on which consensus forms is stated in translatable readings. A theory of consensus on which the community converged because the young were trained in the new vocabulary and the old died out, Planck's famous remark, which Kuhn quoted, would be compatible with Theorem~\ref{t:rational} only if the training were also training in the evaluation of the shared evidence. The record supports this: converts came from the old community as well as from the young, and the arguments addressed to them were stated in the old community's terms.

\subsection{Bayesian reconstruction}

\citet{Salmon1990} argued that Kuhn's account of choice can be reconstructed in Bayesian terms: the shared criteria enter as factors bearing on prior probabilities (simplicity, consistency with accepted theories) and on likelihoods (accuracy, scope), while the individual differences that Kuhn emphasized appear as differences in priors that are washed out by accumulating evidence. SCI-12 defends Bayesian confirmation as the right account of confirmation. The Bayesian reconstruction presupposes that the evidence is a proposition to which both parties assign probabilities, and that the hypotheses are propositions over which they have probability distributions. The first presupposition is Lemma~\ref{l:evidence}: the evidence is stated in readings both communities can use. The second requires that the hypotheses be stated in a common language; where they contain untranslatable readings, the common language is the interpretive extension of Section~\ref{s:interp}. Partial translatability is thus a precondition of the Bayesian account as well as of the qualitative one.

\subsection{Underdetermination and the timing of choice}

SCI-04 argues that theories are underdetermined by evidence only by empirically equivalent rivals that lack independent motivation. The rivals in a revolution are not of this kind. They differ empirically, as the X verdicts and the evidential claims show, and each has independent motivation. A revolution is therefore not a case of underdetermination; it is a case in which the evidence, at a given time, may not yet decide between rivals that it will eventually decide between. That is the situation Proposition~\ref{p:dissent} describes. SCI-19 argues that science is self-correcting but slowly, with a substantial delay between the appearance of contrary evidence and the correction of the record; the delay between the first statement of a successor and its general acceptance in a revolution is an instance of that delay, and the present account explains why it is a delay and not a failure: the evidence was there to be assessed by both sides, and assessing it took time.

\subsection{The role of institutions}

SCI-01 argues that what distinguishes science from non-science is institutionalized error-correction: peer review, replication and the public criticism of results. The norm of public checkability in Assumption~\ref{a:norms} is the norm these institutions implement. The evidential claims of Table~\ref{tab:evid} are, almost without exception, claims that were published, repeated and criticized: Lavoisier's experiments were repeated in England and Germany; Joule's measurements were repeated and refined; Snow's comparison of the water companies was examined by the medical authorities; the eclipse results of 1919 were disputed and confirmed by later expeditions. Theory choice across revolutions is rational in part because the evidence passes through institutions that make it public and checkable on both sides, and the shared norms of SCI-05 are what these institutions serve.

\section{The ten revolutions}\label{s:revolutions}

This section reads each manual of Part~II as a whole. For each revolution we say where the untranslatable cluster lies, which inferences of the predecessor survive translation, which are translated and denied, and how the evidence on which the successor won was stated. The statistics are those of Tables~\ref{tab:stats} and~\ref{tab:further}.

\subsection{CHEM: phlogiston to oxygen}

The manual covers 23 terms and 31 inferences; $\tau = \tiChem$, $\kappa = \kappaChem$. The cluster has three members, \enquote{phlogiston} (CH01), \enquote{dephlogisticated air} (CH02) and \enquote{phlogisticated air} (CH03), and every one of them has a referring reading as well as a non-referring one. This is the pattern of heterogeneity that \citet{Kitcher1978} described: the terms of the cluster were used both to name what was in the chemist's vessel and to express a theory of what it was. The names succeeded; the theory failed. Nine of the 31 inferences are untranslatable, and every one of them contains a descriptive reading of \enquote{phlogiston} or of one of the airs defined from it.

The rest of the vocabulary translates, most of it at the type level. The gases isolated by the pneumatic chemists, fixed air, nitrous air, marine acid air, alkaline air (CH05, CH13, CH15, CH17), translate exactly: their identification depended on preparation and properties, not on phlogistic theory, and the oxygen chemists took them over under new names. The terms for processes and kinds that the phlogistic theory explained, \enquote{calx}, \enquote{metal}, \enquote{combustion}, \enquote{calcination}, \enquote{reduction} (CH06--CH10), translate partially: their referents are preserved, and so are the experimental inferences in which they figure (CH.I03, CH.I05--CH.I07), but the inferences that state the phlogistic account of the processes are lost. Only two inferences are translated and denied: that water is an element (CH.I19) and that air is an element (CH.I27). Both are claims that Lavoisier's chemistry could state and refute, and both were refuted by experiments, the synthesis and analysis of water and the analysis of air, that the phlogistonists themselves performed or accepted.

The chemical revolution is often presented as the paradigm of incommensurability, since the two sides described the same experiments in incompatible terms: what one called the loss of phlogiston the other called the gain of oxygen. The manual shows how much of this incompatibility was a matter of the cluster. The experiments were described by both sides in shared terms (CHEM.E1--E4), and the phlogistic descriptions in terms of the cluster added a theory to the description. When Lavoisier weighed the tin and the retort, both sides could read the balance. The criteria outcome is correspondingly clear on accuracy, where the gravimetric evidence decided, and mixed on scope, where Chang's point about the properties of metals holds.

\subsection{ASTR: Ptolemy to Kepler}

The manual covers 25 terms and 34 inferences; $\tau = \tiAstr$, $\kappa = \kappaAstr$. The cluster has five members: the primum mobile, the planetary orbs, the aether of which they were made, the equant (in its physical reading) and the centre of the universe (in its cosmological reading). All of them belong to the cosmology that Ptolemaic astronomy shared with Aristotelian physics, not to its mathematical models. The models themselves, deferents, epicycles, eccentrics and equants, are translated: exactly or partially where one expression of the successor serves, and token by token where, as with the deferent and the epicycle, different planets require different expressions.

The revolution has the highest number of R verdicts outside physics (six). They mark the inferences of the Ptolemaic models that hold, in the successor, to first order in the eccentricities: the composition of epicycle and deferent (AS.I01), the uniform motion about the equant (AS.I15), the bisection of the eccentricity (AS.I16), the eccentric circle of the Sun (AS.I25). These are the content of the Ptolemaic theory that Kepler retained, and the manual records the retention precisely: what was retained was a mathematical structure, preserved as an approximation, and not the cosmology in which it was embedded. The six X verdicts include the two central disagreements, the immobility of the Earth (AS.I06) and the reality of retrograde motion (AS.I30), and both are disagreements about bodies and motions that both sides could identify.

ASTR is one of the two revolutions in which a criterion favoured the predecessor. Consistency with Aristotelian physics was on Ptolemy's side until Galileo and Newton supplied a physics of a moving Earth \citep{Kuhn1957}. The successor won on accuracy, with Kepler's tables, and on scope and simplicity, with the determination of planetary distances and the replacement of compounded circles by ellipses. All the evidence (ASTR.E1--E5) concerns the positions and periods of the planets, which both sides observed with the same instruments and recorded in the same coordinates.

\subsection{MECH: Newton to Einstein}

The manual covers 23 terms and 35 inferences; $\tau = \tiMech$, $\kappa = \kappaMech$. The cluster is the largest of the physical revolutions, with seven members: \enquote{force} and \enquote{weight} in their gravitational readings, absolute space and absolute time in their metaphysical readings, the ether as a mechanical medium and a privileged frame, perfect rigidity, and action at a distance. Kuhn's cluster of mass, force and weight is partly confirmed: force and weight are there, mass is not, being partially denoting rather than empty (Section~\ref{s:cases}).

The striking feature of MECH is the number of R verdicts: 13 of 35. Most Newtonian inferences are preserved in the domain of low velocities and weak fields, which was the domain in which all the evidence for Newton's theory was gathered. Relativity does not deny that velocities add; it says that they add to an approximation that is excellent for $v \ll c$. It does not deny that a planet describes a fixed ellipse; it says that the ellipse precesses by an amount that, for Mercury, is 43 seconds of arc per century. Only one inference is translated and denied outright: that the measured speed of light depends on the observer's motion (ME.I19). This is the disagreement on which special relativity was founded, and it was decided by experiments, culminating in that of Michelson and Morley, that both sides designed and interpreted in shared terms (MECH.E3).

The untranslatable inferences concern gravitation as a force (ME.I07, ME.I28) and the structure of space and time (ME.I13, ME.I15, ME.I35). In these the Newtonian claims are not denied but replaced: general relativity has no gravitational force to make claims about, and no absolute space. Yet even here the Newtonian law of gravitation is recovered in the weak-field limit as a law about the metric (ME.I08). The relativistic revolution is thus the clearest case of a pattern that the record shows repeatedly in physics: an untranslatable ontology over a translatable mathematics.

\subsection{HEAT: caloric to energy}

The manual covers 23 terms and 30 inferences; $\tau = \tiHeat$, $\kappa = \kappaHeat$. The cluster consists of \enquote{caloric} (HE01) and the two terms derived from it, free and combined caloric (HE06, HE07). Its structure matches the analysis that SCI-02-a gives of the caloric case. The readings of \enquote{caloric} in which it named a quantity measured by calorimetry (HE01b) refer to heat; the readings in which it named a fluid with properties (HE01a) or a conserved stuff (HE01c) refer to nothing. The inferences that made the theory successful, the method of mixtures (HE.I03), latent heat (HE.I05), Laplace's adiabatic correction (HE.I09, HE.I10, HE.I27) and Carnot's principle (HE.I12), are preserved, and none of them requires an untranslatable reading.

Two disagreements stand out. The claim that in a cycle the heat given to the cold reservoir equals that taken from the hot (HE.I13) is translated and denied: it is false by the first law. The claim that heat cannot be generated without limit from a finite body (HE.I28) is denied by Rumford's and Joule's experiments, which are the evidence on which the revolution turned (HEAT.E1, E2). Both experiments were described, by their authors and by their critics, in the vocabulary of heat, work and temperature, which both theories shared.

\subsection{QUAN: classical to quantum}

The manual covers 21 terms and 30 inferences; $\tau = \tiQuan$, the highest of the ten, and $\kappa = \kappaQuan$. The cluster consists of three terms: the trajectory and the orbit of an electron inside an atom (QU01, QU02), and the classical state as a point in phase space (QU18). It is small because the vocabulary of classical physics, position, momentum, energy, angular momentum, charge, is taken over by quantum mechanics as the vocabulary of observables. What changes is the inferential role of these terms: the claims that they always have definite values (QU.I06), that they can be measured together (QU.I07), that energy and angular momentum vary continuously (QU.I10, QU.I22), are translated and denied. QUAN has seven X verdicts, and more X than U: quantum mechanics contradicts classical physics more than it abandons it.

This finding cuts against a common view, on which the quantum revolution is the paradigm of conceptual discontinuity. The discontinuity is real, but it is a discontinuity of laws stated in a shared vocabulary, which is why the experiments that forced the revolution, on black-body radiation, specific heats, the photoelectric effect and the Stern--Gerlach splitting (QUAN.E1--E5), could be described in classical terms and could refute classical predictions. The verdicts of QU.I06 and QU.I07 are given for the standard formulation of quantum mechanics; on the Everett interpretation, which SCI-09 defends, the denial of definite values holds relative to branches, and the verdicts are unchanged.

\subsection{BIOL: fixed species to descent}

The manual covers 24 terms and 31 inferences; $\tau = \tiBiol$, $\kappa = \kappaBiol$. The cluster consists of the plan of creation (BI05), the archetype (BI06) in its idealist reading, the centre of creation (BI10) in its creationist reading, and the chain of being (BI16). As in the chemical case, three of the four have referring readings: the unity of plan is the pattern of shared ancestral structure, the archetype is the ancestor, the single centre of creation is the area of origin. Darwin's own procedure, as we noted in Section~\ref{s:cases}, was to give these readings referents in his theory and to drop the rest.

BIOL has the largest number of X verdicts, eight: the immutability of species (BI.I01), the permanence of the species--variety distinction (BI.I02), the fixed limits of variation (BI.I03), perfect adaptation (BI.I14), catastrophic extinction (BI.I18), the designed purpose of hybrid sterility (BI.I24), the designed balance of nature (BI.I26), the fixity of instinct (BI.I27). These are the claims the \emph{Origin} argued against, chapter by chapter, and that it could argue against because they were stated in a vocabulary of species, varieties, organs and instincts that Darwin and his opponents shared. The revolution was fought over claims both sides could state; that is why it was fought in print, in the reviews of the \emph{Origin}, with arguments each side answered.

Consistency favoured the predecessor, through Kelvin's age of the Earth and Jenkin's swamping argument; the second appears in HERE as the inference HR.I03, translated and denied by Mendelian genetics. The successor won on scope: one hypothesis explained the natural system, homology, rudimentary organs, embryology, succession and distribution (BIOL.E1--E5), each a domain the predecessor described in the same terms.

\subsection{PHYS: Galen to Harvey}

The manual covers 26 terms and 30 inferences; $\tau = \tiPhys$, the lowest of the ten, and $\kappa = \kappaPhys$, also the lowest. The cluster is the largest, eight terms: the three spirits (GA03--GA05), the septal pores (GA08), innate heat (GA14), sooty vapours (GA15), the faculties (GA18) and the rete mirabile (GA26). Eleven inferences are untranslatable, more than in any other revolution. Galenic physiology explained the functions of the body by agents, spirits and faculties, that the successor does not recognize, and the inferences that state these explanations cannot be translated.

Yet the anatomy translates almost completely. The vessels, the chambers of the heart, the valves, the lungs, the liver and the blood are identified by dissection, and their names translate at the type level: the arterial vein is the pulmonary artery (GA09), the venous artery the pulmonary vein (GA10). The six X verdicts are disagreements about the anatomy and its motions: whether blood is made continuously in the liver (GA.I01), whether it is consumed in the parts (GA.I02), whether the venous artery carries air (GA.I05), whether the valves retard flow to the extremities (GA.I12). Harvey's arguments in \emph{De motu cordis} \citep{French1994} are addressed to exactly these claims, by vivisection, ligature and calculation (PHYS.E1--E3), and his readers could repeat his procedures.

PHYS shows that a low translation index is compatible with a decisive comparison. The cluster was large, but it lay in the explanatory superstructure; the comparison was conducted in the anatomical base.

\subsection{OPTI: particles to waves}

The manual covers 18 terms and 25 inferences; $\tau = \tiOpti$, $\kappa = \kappaOpti$. The cluster consists of the light corpuscle (OP02), the fits (OP03) in their reading as dispositions of corpuscles, and the sides of rays (OP08) in their reading as properties of corpuscles. Newton's measurements of the fits (OP03b) and his description of the sides (OP08b) refer to wavelength and polarization, as Young and Fresnel showed by using Newton's data \citep{Buchwald1989}. The emission theory had a rich phenomenology, of refraction, dispersion, thin-film colours, inflection and double refraction, and the wave theory took over the whole of it; seven untranslatable inferences explain these phenomena by forces on corpuscles, and these the wave theory replaced.

The two denied inferences are the two crucial ones: that light does not bend into the geometrical shadow (OP.I12), denied by the bright spot in the shadow of a disk, and that light travels faster in water than in air (OP.I04), denied by Foucault's measurement of 1850. Both experiments were proposed as tests that would discriminate between the theories, and both were described in the vocabulary of shadows and speeds that both theories used.

\subsection{MIAS: miasma to germs}

The manual covers 21 terms and 31 inferences; $\tau = \tiMias$, $\kappa = \kappaMias$. The cluster is large, seven terms: miasma (MI01), ferment (MI03) and contagium (MI05) in their chemical readings, spontaneous generation (MI07), epidemic constitution (MI08), sewer gas (MI12) and the soil-borne cholera poison (MI13b). Eleven inferences are untranslatable. The miasmatic theory, like Galenic physiology, explained by agents the successor does not recognize.

But MIAS also shows token-level translation most vividly. The miasma of the marshes refers to malarial infection (MI01b); the miasma of cholera districts refers to Vibrio cholerae carried in water (MI01c); Snow's morbid matter refers to the same organism (MI13a); Henle's contagium animatum refers to pathogenic microorganisms (MI05b). The sanitary reformers who acted on the miasmatic theory, by draining marshes and removing filth, did what the germ theory would also have recommended, and the inferences on which they acted (MI.I02, MI.I13) are preserved. The evidence that decided the revolution, Pasteur's flasks, Koch's cultures, Snow's water companies, Lister's dressings (MIAS.E1--E5), was stated in terms of broths, blood, deaths and wounds that both sides used.

\subsection{HERE: blending to Mendel}

The manual covers 17 terms and 23 inferences; $\tau = \tiHere$, $\kappa = \kappaHere$. The cluster consists of gemmules (HR01) in their reading as particles shed by the body, use-inheritance (HR05), telegony (HR06) and Galton's ancestral law (HR13) in its reading as a law of germinal contributions. Of these, gemmules and ancestral heredity have referring readings: dormant gemmules are recessive alleles, and Galton's regressions are what the additive polygenic model predicts \citep{Fisher1918}.

The revolution has five X verdicts, among them the blending of characters (HR.I02), the swamping of rare variants (HR.I03, HR.I23), the uniform intermediacy of hybrid offspring (HR.I11), and the efficacy of selection within a pure line (HR.I14). These are the disagreements on which the Mendelians and the biometricians argued after 1900, and Fisher's reconciliation of 1918 shows how much of the biometricians' vocabulary, regression, correlation, variation, survived translation into Mendelian terms.

\subsection{Patterns}

Three patterns emerge from the ten manuals.

First, in every revolution the untranslatable cluster lies in the predecessor's explanatory posits, not in its descriptive vocabulary. The terms that resist translation are names for agents and mechanisms (phlogiston, spirits, miasmata, caloric, gemmules, orbs) or for frameworks (absolute space, the plan of creation). The terms by which the phenomena were identified translate. This is what makes Proposition~\ref{p:nototal} true of the record and not merely of the logic.

Secondly, the cluster terms are typically heterogeneous: of the \nclusterall{} cluster terms, most have at least one referring reading alongside the non-referring ones. The predecessor's names for its posits were often used to refer to things the successor recognizes, and the untranslatability is confined to the descriptive readings that express the predecessor's theory of those things.

Thirdly, the relation between translated disagreement and untranslatable abandonment differs between the sciences. In physics, where the predecessor survives as a limiting case, R verdicts are frequent and the disagreements are quantitative. In physiology and medicine, where the predecessor posited agents, U verdicts are frequent. In biology and genetics, X verdicts are frequent: the revolutions turned on claims both sides could state. None of these patterns supports incommensurability; each shows a different way in which comparison survived the loss of part of the vocabulary.

\section{Further formal results}\label{s:formal}

We record three structural results about manuals, which bear on the robustness of the conclusions and prepare the ground for Volume~II.

\begin{proposition}[Refinement of readings]\label{p:refine}
Let $M'$ be obtained from $M$ by splitting a reading of a term $t$ into two readings with the same referent status (both referring or both not). Then $\tau(M') = \tau(M)$ if $t$ is translated exactly, partially or not at all; if $t$ is token-dependent the score of $t$ changes, but $\tau(M')$ lies strictly between 0 and 1 if and only if $\tau(M)$ does.
\end{proposition}

\begin{proof}
The status of $t$ depends only on whether it has a type-level image and whether some reading refers, both of which are unchanged by the split. If $t$ is not token-dependent, its score is 1 or 0 and unchanged; the weights are unchanged, since the inferences are. If $t$ is token-dependent, its score $n_{\mathrm{ref}}/n$ changes to $(n_{\mathrm{ref}}+1)/(n+1)$ or $n_{\mathrm{ref}}/(n+1)$, which lies in $(0,1)$ if and only if $n_{\mathrm{ref}}/n$ does, since $0 < n_{\mathrm{ref}} < n$ is preserved by both changes when the split reading has the referent status of the original. The cluster is unchanged, so by Lemma~\ref{l:bounds} the strict inequalities are preserved.
\end{proof}

The proposition answers a natural worry about the manuals: that the index depends on how finely readings are individuated. It does, but the conclusion of Theorem~\ref{t:partial} does not.

\begin{proposition}[Composition]\label{p:compose}
Let $M_1$ be a manual from $T_1$ to $T_2$ and $M_2$ a manual from $T_2$ to $T_3$, where every expression in the range of $M_1$ is a reading in the domain or the cluster of $M_2$. Then the composite $M_2 \circ M_1$ is a manual from $T_1$ to $T_3$, and its cluster contains the cluster of $M_1$. In particular, if $T_1$ is partially translatable into $T_2$ by $M_1$ and the composite has a referring reading, then $T_1$ is partially translatable into $T_3$.
\end{proposition}

\begin{proof}
Reference is transitive along the composite: if $r$ refers in $T_2$ to what $M_1(r)$ refers to, and $M_1(r)$ refers in $T_3$ to what $M_2(M_1(r))$ refers to, then $r$ refers to what $M_2(M_1(r))$ refers to. A reading of $T_1$ not in the domain of $M_1$ is not in the domain of the composite, so the cluster of $M_1$ is contained in that of the composite. The last claim follows by Lemma~\ref{l:bounds}.
\end{proof}

Composition explains why partial translatability does not decay into incommensurability over a sequence of revolutions. The cluster grows along a sequence, since what one revolution abandons is not recovered by the next; but a referring reading of $T_1$ that is translated by the composite remains translated. The descriptive vocabulary of the phenomena, which by Proposition~\ref{p:nototal} each successor shares with its predecessor, is carried along the sequence. This is the semantic counterpart of the corollary of SCI-02-a on chains of theories, on which the retention of essential content propagates backward along a chain of theories.

\begin{proposition}[No overlap constrains types only]\label{p:overlap}
Suppose the kind terms $k$ of $T_1$ and $k'_1, k'_2$ of $T_2$ are such that the extension of $k$ overlaps those of $k'_1$ and $k'_2$ without being included in either, so that no translation of the type $k$ respects the taxonomy of $T_2$ (Kuhn's no-overlap principle). Then $M$ has no single-expression type-level image for $k$; but every reading of $k$ whose tokens are applied to members of the extension of $k'_1$ or of $k'_2$ is in the domain of $M$.
\end{proposition}

\begin{proof}
A single-expression type-level image of $k$ would be a kind term of $T_2$ with the extension of $k$, which the taxonomy of $T_2$ excludes, or a disjunction, which is a list in the sense of Definition~\ref{d:manual}(i). A reading of $k$ applied to members of $k'_1$ refers, under Assumption~\ref{a:token}, to the things it is applied to, which $T_2$ classifies as $k'_1$; so it is in the domain of $M$.
\end{proof}

Kuhn's \enquote{planet} is the example \citep{Kuhn1991}. Its extension in Ptolemaic astronomy included the Sun and the Moon and excluded the Earth; in the Copernican taxonomy the Sun is a star, the Moon a satellite and the Earth a planet. No Copernican kind term has the Ptolemaic extension. The manual records exactly this: \enquote{planet} (AS01) has no type-level image, and each of its four readings refers. Proposition~\ref{p:overlap} shows that this is the general case: the no-overlap principle, which Kuhn made the basis of his later thesis, blocks type-level translation and leaves token-level translation intact. Volume~II develops the point for the taxonomic account as a whole.

\section{Cases}\label{s:cases}

We discuss six entries of Part~II in detail. Each illustrates a different way in which a term can resist type-level translation while the claims containing it remain comparable.

\paragraph{Dephlogisticated air (CH02).} The term has three readings. CH02a is Priestley's description of the air obtained by heating the red calx of mercury, in which a candle burned with a vigorous flame; its tokens refer to what was in his vessel, which was oxygen. CH02c is Cavendish's air exploded with inflammable air in his synthesis of water, which was also oxygen. CH02b is the type description, common air deprived of phlogiston; nothing satisfies it. The manual therefore has no type-level image, and the term is token-dependent with score $2/3$. The inferences in which it figures show what this means in practice. CH.I10 (the calx yields the metal and dephlogisticated air) and CH.I11 (dephlogisticated air supports combustion longer than common air) are preserved; they were the claims on which Priestley and Lavoisier agreed. CH.I12 (dephlogisticated air is avid of phlogiston because deprived of it) is untranslatable; it was the claim on which they differed, and their difference was not a disagreement that either could state in the other's terms, but a difference in what they took the air to be. This is the case on which \citet{Kitcher1978} built the theory of reference potential, and the manual confirms his analysis while recording Kuhn's point that the type is not translated.

\paragraph{Mass (ME01).} Field's example. The term is given a type-level image that is a list, rest mass and relativistic mass, and is translated partially. Three of its inferences are recorded R, preserved on one reading only: invariance (ME.I01) on the rest-mass reading, additivity (ME.I02) and conservation (ME.I03) on the relativistic-mass reading. The proportionality of gravitational and inertial mass (ME.I22) is preserved outright, as the equivalence principle. Field's semantics says that a sentence true on one reading and false on the other is neither true nor false; the verdict R records the same fact without committing us to Field's truth-value gaps. That \enquote{mass} is translated partially, while \enquote{force} and \enquote{weight} are token-dependent with non-referring gravitational readings (ME02b, ME19a), locates the untranslatability of Newtonian mechanics where general relativity places it: in the treatment of gravitation as a force.

\paragraph{Deferent and epicycle (AS05, AS06).} These are the clearest cases of heterogeneity in which every reading refers. The deferent of a superior planet corresponds to the planet's heliocentric orbit (AS05a), while the deferent of Venus or Mercury corresponds to the Earth's orbit (AS05b); the main epicycle stands in the inverse relation (AS06a, AS06b). Both terms are token-dependent with score 1. No single expression of the Copernican vocabulary translates \enquote{epicycle}, because the role the epicycle played was played by different orbits for different planets; but every token refers. The Copernican scope advantage recorded for ASTR rests on exactly this fact: when the epicycles and deferents are translated token by token, the ratios of their radii yield the relative distances of the planets (ASTR.E3), which the Ptolemaic system left undetermined. Here translation, token by token, is not merely compatible with the comparison of the theories; it is the comparison.

\paragraph{Carnot's caloric (HE01d).} The reading in which Carnot speaks of the caloric that falls from the hot to the cold body in a reversible engine has been given two referents. Kelvin read it as heat, on which reading Carnot's claim that the caloric is conserved through the cycle is false. \citet{Callendar1911} read it as what Clausius later called entropy, on which reading the claim is true for a reversible cycle, and the waterwheel analogy (HE.I14) becomes the relation $W = S(T_1-T_2)$. We record the reading as partially denoting both, and the inference as R. The case shows that the assignment of referents is sometimes contested, and that the manual can record the contest rather than settle it. Either way, the reading is in the domain of the manual and the inference is comparable: the question whether Carnot was right about the conservation of what falls was one that his successors could ask and answer, and the answer depends on which quantity is meant.

\paragraph{Rete mirabile (GA26).} Galen described a network of arteries at the base of the brain in which vital spirit was elaborated into animal spirit. The network exists in the ox, which he dissected (GA26b); it does not exist in humans, as Berengario and Vesalius pointed out, and the human reading (GA26a) has no referent. The term is token-dependent with score $1/2$. The case is instructive because the untranslatability has nothing to do with theory-ladenness or meaning change. It arises from an error of extrapolation between species, which the successor can describe exactly. What cannot be translated is the claim about the human brain; what can be translated is the claim about the ox, and the comparison of the two is the comparison of anatomy with anatomy.

\paragraph{Archetype (BI06).} Owen's archetype was the ideal pattern of which every vertebrate skeleton is a modification. On the reading in which it is an idea in the divine mind (BI06a) it has no referent in Darwinian biology. On the reading in which it is the generalized pattern that Owen drew (BI06b), Darwin himself supplied a referent: the structure of the common ancestor of the vertebrates \citep[ch.~13]{Darwin1859}. The inference BI.I08, that all vertebrate skeletons are modifications of one pattern, is preserved under that reading. The case shows a successor theorist doing exactly what the manuals do: taking a term of his predecessor's, distinguishing its readings, and giving one of them a referent in his own theory. Darwin's reinterpretation of the archetype was an act of token-level translation, and he used it to argue that his theory explained what Owen's described.

Two further entries deserve mention. The miasma of the marshes (MI01b) refers to malarial infection by Plasmodium, carried by mosquitoes breeding in the marshes, so that the miasmatists' claim that draining the marshes reduces intermittent fever (MI.I02) is preserved, though their claim that fevers are caused by the emanations of decay (MI.I01) is not. And Darwin's dormant gemmules (HR01b), to which he ascribed reversion, refer to recessive alleles transmitted unexpressed; the inference that reversion is due to them (HR.I06) is preserved in the cases of reversion by recessives.

\section{Objections and replies}\label{s:objections}

\begin{objection}[Token reference is not translation]\label{o:tokens}
Kuhn's reply to Kitcher applies to the translation index. A translation must render a term uniformly; a manual that renders \enquote{dephlogisticated air} as \enquote{oxygen} here and as nothing there is not a translation but a commentary, and the index measures the commentary. What Kuhn denied, the existence of a translation of the cluster, the index does not touch.
\end{objection}

\begin{reply}
The objection is right about what Kuhn denied and wrong about what the conclusion requires. The manuals distinguish type-level from token-level translation (Definition~\ref{d:status}), and Proposition~\ref{p:robust} reports the index computed by type-level translation alone (scheme S1). It lies strictly between 0 and 1 for every revolution: even by Kuhn's standard of translation, most of the vocabulary of each predecessor is translated, and some is not. What remains of the objection is the claim that the comparison of theories requires type-level translation of the cluster. It does not. The comparison requires that the evidence be stated in shared vocabulary (Lemma~\ref{l:evidence}) and that the claims at issue be either stated in it or understood by interpretation (Principle~\ref{p:interp}). Kuhn conceded both \citep{Kuhn1983}. The untranslatability of the type of \enquote{phlogiston} prevented a phlogistic chemist from saying in Lavoisier's language what he believed; it did not prevent him from weighing the calx.
\end{reply}

\begin{objection}[The index is arbitrary]
The value of $\tau$ depends on the choice of terms, readings and inferences. Another investigator, choosing differently, would obtain other values; a defender of incommensurability could choose the theoretical core of each predecessor and obtain values near 0.
\end{objection}

\begin{reply}
The values depend on the choices; the conclusions do not, within wide limits. The conclusions are that $\tau$ lies strictly between 0 and 1 and that some shared criterion favoured the successor on translatable evidence. The first holds whenever the chosen vocabulary includes one term with a non-referring reading and one with a referring reading, and Proposition~\ref{p:nototal} shows that a vocabulary that described the predecessor's evidence would include the second. A defender of incommensurability who restricted the vocabulary to the untranslatable cluster would obtain $\tau=0$, but would then have measured the cluster, not the theory, and would have no vocabulary in which to state what the theory was a theory of. Within the vocabulary actually chosen, the alternative scorings of Proposition~\ref{p:robust} move the indices by as much as a third without moving any of them to a bound. The choice of terms and inferences followed a stated procedure (Part~II, preface), and the manuals are published so that it can be checked.
\end{reply}

\begin{objection}[Methodological incommensurability]
The criteria were not shared. Feyerabend argued that every methodological rule has been violated by the successful science of its day \citep{Feyerabend1975}; the Galenists and Harvey disagreed about the weight to be given to vivisection and to quantity, and the Ptolemaic astronomers did not demand of an astronomical model that it be physically true. Kuhn's list of five is an abstraction that conceals these differences.
\end{objection}

\begin{reply}
Kuhn's list is an abstraction, and Assumption~\ref{a:criteria} claims only what it states: that the criteria are accepted on both sides, not that they are weighed or applied in the same way. The objection's examples concern application and weight. Galenists accepted that a physiology must agree with what dissection shows, and the debate over Harvey turned on what dissection showed; Harvey's quantitative argument (PHYS.E1) used a measure of the ventricle's capacity and a count of the heart's beats, which his opponents could check and some did. The Ptolemaic astronomers sought accuracy above all, and Kepler's tables won on accuracy (ASTR.E1, ASTR.E2). Feyerabend's examples show that particular rules of method were violated by successful science; SCI-05-a argues that his examples violate rules but not the norms of Assumption~\ref{a:norms}, and that conclusion is what we need.
\end{reply}

\begin{objection}[Arrow's theorem]
\citet{Okasha2011} showed that Kuhn's account of theory choice, as an aggregation of rankings by several criteria, falls under Arrow's theorem: if the criteria yield only ordinal rankings, no aggregation rule satisfies unrestricted domain, weak Pareto, independence of irrelevant alternatives and non-dictatorship. So either one criterion dictates, or the choice is not rational.
\end{objection}

\begin{reply}
Definition~\ref{d:rational} does not require an aggregation rule. It requires that a criterion favour the successor on shared evidence; it is a condition on the existence of a shared reason, not on the existence of a function from profiles of criteria to choices. Where the support was dominant, as in eight of the ten revolutions, the Pareto condition alone, which every rule Arrow considers must satisfy, selects the successor among the two theories compared; the impossibility concerns the consistency of a rule over all possible profiles and all sets of alternatives, not the verdict on a profile in which no criterion opposes. Where support was not dominant, as in ASTR and BIOL, the choice depended on weighing, and Proposition~\ref{p:dissent} records that dissent was rational for a time. Okasha himself discusses ways of escaping the theorem by enriching the information the criteria provide; Volume~III takes them up.
\end{reply}

\begin{objection}[Indeterminacy of translation]\label{o:quine}
Quine argued that translation is indeterminate: incompatible manuals fit all the evidence of a speaker's dispositions \citep{Quine1960}. If so, there is no fact about whether a token of \enquote{dephlogisticated air} refers to oxygen, and the manuals record one choice among many.
\end{objection}

\begin{reply}
If the Quinean thesis were true, it would undermine translation between contemporaries as much as between paradigms, and the incommensurability theorist could not use it to mark revolutions out. LNG-06 argues that radical translation is constrained by more than stimulus meaning and that the constraints select determinate translations, and LNG-16 that reference is fixed by the same constraints. On those conclusions, which we accept, the referential assignments of the manuals are determinate where the historical evidence is, and where it is not, as with Carnot's caloric (HE01d) or the continued fever of the London physicians (MI10), the manuals record partial denotation rather than an arbitrary choice.
\end{reply}

\begin{objection}[Theory-laden observation]
The evidential claims are not neutral. \citet{Hanson1958} and \citet{Kuhn1962} argued that observation is laden with theory, so that the phlogistic chemist and the oxygen chemist did not observe the same thing when they watched a metal calcine. The vocabulary that Lemma~\ref{l:evidence} calls translatable is itself the vocabulary of a theory, and its sharing presupposes what is in question.
\end{objection}

\begin{reply}
The evidential claims of Table~\ref{tab:evid} are laden with theory, but not with the theories at issue. The weight of a calx is laden with the theory of the balance, the position of Mars with the theory of the instruments and of refraction, the speed of light in water with the theory of Foucault's rotating mirror; these theories were accepted on both sides of the revolutions in question. What Lemma~\ref{l:evidence} requires, and what check C9 verifies, is that the readings used in stating the evidence have referents in the successor, so that the successor can state the evidence; and, since these readings belong to the predecessor's vocabulary, the predecessor can state it too. Theory-ladenness by a third theory shared by both is no obstacle to comparison.
\end{reply}

\begin{objection}[Whig history]
The manuals identify referents by the successor theory as since developed. This is Whig history: it judges the past by the present and reads into it what the present knows. The phlogistic chemists did not refer to oxygen, since they had no concept of oxygen; to say that they did is to impose our categories on them.
\end{objection}

\begin{reply}
The objection confuses reference with conception. That Priestley had no concept of oxygen does not show that he did not refer to it, any more than the fact that the first observers of a disease lacked the concept of its pathogen shows that they did not refer to the disease. Reference is fixed by the samples and procedures a term is used with, and Priestley's \enquote{dephlogisticated air}, in the reading in which it named the gas in his vessel, was used with oxygen. The manuals do not say that Priestley believed what Lavoisier believed; they record, in the predecessor's own terms, the role each term played, and they record separately what the successor takes it to have referred to. A history that refused to say what the past referred to could not say whether its claims were true, which is what the historian of science, as distinct from the historian of opinion, wants to know.
\end{reply}

\begin{objection}[Holism of meaning]
Feyerabend's argument was that meanings are fixed by the whole theory. If so, a change anywhere in the theory changes the meaning of every term, and the type-level translations that the manuals record, such as \enquote{fixed air} by \enquote{carbon dioxide}, are as illegitimate as the token-level ones.
\end{objection}

\begin{reply}
The manuals preserve reference, not meaning, and holism of meaning, if true, is compatible with preservation of reference; this is the point of \citet{Scheffler1967} and \citet{Devitt1979}. Where the holist objects that comparison needs sameness of meaning, the reply is that it needs only that the claims compared be about the same things, and that disagreements be stated (Principle~\ref{p:disagree}). The manuals also record the residue of holism that is real: the inferential role of a term is not wholly preserved when an inference in which it figures is untranslatable, and Definition~\ref{d:status} makes such a term partially translated. The large number of partial translations (\nstatP{} of \ntermsall{} terms) is the measure of how much holism is true.
\end{reply}

\begin{objection}[Premature abandonment]
\citet{Chang2012} argued that the phlogiston theory was abandoned prematurely: it had explanatory resources, such as the account of the common properties of metals, that the oxygen theory lacked and that were later recovered in the theory of electrons. A choice that discarded these resources was not rational.
\end{objection}

\begin{reply}
The manual records Chang's point: scope is marked as mixed for CHEM, with the common properties of metals on the phlogistic side, and the inference CH.I26, that metals resemble one another because they contain phlogiston, is untranslatable, so that the explanation it offered was lost with the term. But Chang's argument concerns whether it would have been better to keep both theories in play, not whether the reasons for preferring the oxygen theory were shared reasons. They were: the gravimetric evidence (CHEM.E1--E4) was stated in vocabulary the phlogistonists used, and several of them, including Kirwan, were persuaded by it. A rational choice may lose something of value. Whether the loss could have been avoided by a more pluralist policy is a question about the organization of inquiry, on which our conclusion takes no position.
\end{reply}

\section{Consequences}\label{s:consequences}

\paragraph{For scientific realism.} SCI-02 argues that mature theories are approximately true and that the pessimistic meta-induction fails because past theories were partially true. Partial translatability is the semantic precondition of the second claim. If past theories were wholly untranslatable into present ones, the question whether they were partially true by present lights could not be asked. The manuals show that it can be asked and, for most of the vocabulary of each predecessor, answered: the H and R verdicts (\nverdH{} and \nverdR{} of \ninfall{}) record inferences of past theories that hold, outright or in a domain, by the lights of their successors. The realist's selective claim, that the essential content of mature theories is retained, can now be tested against a vocabulary-wide record and not only against the cases realists choose; Proposition~\ref{p:retention} says what the test must find, and the manuals of HEAT and MECH find it.

\paragraph{For relativism.} The incommensurability thesis has been used to support a relativism on which the truth of scientific claims is relative to paradigms. Theorem~\ref{t:partial} removes the support: the claims of a predecessor are, for the most part, claims the successor can state, and where it denies them it denies what the predecessor asserted. A claim that is denied across a revolution is not true relative to one paradigm and false relative to another; it is a claim about water, or Mars, or the speed of light in water, about which one side was right.

\paragraph{For the history of science.} The manuals suggest a method for writing the history of a revolution that is neither Whiggish nor relativist: state, for each term, its role in the old theory, and separately, for each reading, its referent in the new. The first part is the historian's reconstruction of the past in its own terms; the second is the scientist's assessment of it. Keeping them apart, as Definition~\ref{d:manual} does, allows both to be stated without either distorting the other.

\paragraph{For science education.} If successive paradigms were incommensurable, the history of a science could not be taught in the vocabulary of its present, and students would have to learn each past theory as a foreign language. Partial translatability means that most of the history can be taught in present terms, with the cluster taught by interpretation. The manuals indicate what the cluster is in each case, and so what must be taught as a foreign language: in chemistry, three terms; in physiology, eight.

\section{What the result does not establish}\label{s:not}

The result does not establish that every choice made in a revolution was correct, or that the successor was in every respect better than the predecessor. ASTR and BIOL show criteria favouring the predecessor, and CHEM, PHYS and the remaining entries marked mixed show respects in which the predecessor did better.

It does not provide an algorithm of theory choice. Definition~\ref{d:rational} states a condition that rational choices satisfy, not a procedure that generates them, and it is compatible with Kuhn's view that individual scientists weigh the criteria differently.

It does not establish that the translation index measures closeness to truth, or progress. A successor with a higher index relative to its predecessor is one that says more of what its predecessor said, not one that is truer; the relation between the two is the subject of SCI-02.

It does not establish partial translatability for revolutions not in the table, except through Proposition~\ref{p:local}, whose antecedent must then be established case by case, and through the inductive support the table provides. In particular, we have not examined revolutions in the social sciences, where the conditions for Proposition~\ref{p:nototal} may fail.

Finally, the manuals are compiled from the historical record by judgment. The judgments are stated entry by entry, and each can be contested; what the program checks is that the conclusions follow from the judgments, not that the judgments are right. The most contestable are flagged in the manuals by partial denotation (HE01d, MI10, HR10) and in the criteria table by the outcome \enquote{?}, which marks a criterion on which we make no claim.

\section{Conclusion}\label{s:conclusion}

Successive paradigms are partially translatable: a translation manual that preserves the reference of term tokens and records which inferences of the predecessor survive translation yields, in every revolution examined, a translation index strictly between 0 and 1, and Kuhn's thesis of local incommensurability entails that it must. Theory choice remains rational across revolutions: in every revolution examined, a criterion that both communities accepted favoured the successor, on evidence both could state, under norms both observed. Incommensurability, in the only form in which it is true, is the existence of an untranslatable cluster of readings; it marks what was abandoned, not what was compared.

\clearpage
\appendix
\section*{Part II: Translation manuals for ten revolutions}\label{s:appendix}
\addcontentsline{toc}{section}{Part II: Translation manuals for ten revolutions}

\section{Preface to the manuals}\label{s:preface}

\paragraph{What the manuals contain.} For each of ten revolutions, identified as CHEM (phlogistic to oxygen chemistry), ASTR (Ptolemaic to Copernican--Keplerian astronomy), MECH (Newtonian mechanics to relativity), HEAT (caloric theory to thermodynamics and kinetic theory), QUAN (classical to quantum mechanics), BIOL (fixed species to Darwinian descent), PHYS (Galenic physiology to the circulation), OPTI (emission to wave optics), MIAS (miasmatic to germ theory) and HERE (blending inheritance and pangenesis to Mendelian genetics), Section~\ref{s:manuals} gives a translation manual in the sense of Definition~\ref{d:manual}. Each term of the predecessor $T_1$ has an identifier (CH01, AS05, \ldots), its role in $T_1$, its type-level image in $T_2$ (a single expression; a list marked (pd) for partial denotation; or a dash for none), its status, weight $w$ and score $s$, and the inferences in which it figures, grouped by verdict. Under it are its readings (CH02a, CH02b, \ldots), each with the context of use that fixes it and its referent in $T_2$. For terms without a type-level image the inferences are listed by reading. Section~\ref{s:inferences} lists the inferences of each predecessor (CH.I01, \ldots), each with its translation $M(I)$ or, for verdict U, an interpretation in Kuhn's sense, prefixed \enquote{Interpretation}; the readings it uses, with $^{\ast}$ marking a non-referring reading and $^{\circ}$ a term translated other than exactly; and its verdict. Section~\ref{s:criteria} gives the outcome on the five shared criteria and the evidential claims behind every favourable outcome. Section~\ref{s:stats} gives the computed statistics.

\paragraph{How the manuals were compiled.} The predecessor vocabulary of each revolution consists of the terms that figure in the central claims of the predecessor as presented in the primary sources and the standard histories cited in the written part \citep{Kuhn1957,Fox1971,Buchwald1989,Jammer1966,Ruse1979,French1994,Musgrave1976}. Readings were distinguished where the historical record shows a term's reference fixed in more than one way. The inferences are the claims of the predecessor that its textbooks and its critics treated as characteristic, together with the inferences needed to give every term at least one inferential role. Referents and verdicts were assigned by the standard of Assumption~\ref{a:token}. Every criterion outcome recorded as $+$, $-$ or $\pm$ rests on a documented episode named in the table; where we were not confident of the outcome we recorded \enquote{?} and made no claim.

\paragraph{The program.} The generator (\texttt{gen/generate.py}, Python standard library only) reads the hand-curated data and computes, for every term, its status by Definition~\ref{d:status}, its weight and score by Definition~\ref{d:score}, and, for every revolution, $\tau$, $\kappa$, $\pi$, the untranslatable cluster, the \emph{neutral core} (inferences all of whose terms are translated exactly), and $\tau$ under the scoring schemes S1--S4 of Proposition~\ref{p:robust}. It asserts nine checks, and the tables are written only if all succeed:
\begin{description}[leftmargin=2.2em,style=sameline]
\item[C1] integrity: identifiers unique; a token reading is cited for a term exactly when the term lacks a type-level image; every reading of a term with an image has a referent; every non-referring reading figures in some inference; every inference has a rendering, which is an interpretation exactly when the verdict is U;
\item[C2] the verdict U is recorded exactly for inferences containing a reading with no referent;
\item[C3] every term figures in at least one inference;
\item[C4] $0<\tau<1$ for every revolution (Theorem~\ref{t:partial});
\item[C5] in every revolution at least one shared criterion favoured the successor;
\item[C6] the cluster is non-empty and a proper part of the vocabulary;
\item[C7] every revolution contains at least one inference translated and denied (verdict X);
\item[C8] $0<\tau<1$ under S1, S3 and S4;
\item[C9] every favourable outcome is supported by an evidential claim stated in referring readings, and every claim is cited only for outcomes recorded as favourable (Lemma~\ref{l:evidence}).
\end{description}
Check C2 is the mechanical counterpart of Lemma~\ref{l:disagree}: since the verdict X is recorded only for inferences all of whose readings refer, every recorded disagreement is a translated one.

\paragraph{Scope and exhaustiveness.} The manuals are not enumerations of all possible structures, and nothing in them is exhaustive in the sense of a combinatorial search. What is exhaustive is the checking: every entry of every manual is checked by the program against the definitions, and every claim the written part makes about a revolution is computed from the entries. The number of revolutions is limited by our confidence in the historical record; Volume~II extends the table.

\paragraph{Summary.} The manuals contain \ntermsall{} terms with \ntokensall{} readings, of which \nrefall{} have a referent in the successor, and \ninfall{} inferences. Of the terms, \nstatE{} are translated exactly, \nstatP{} partially, \nstatT{} token by token and \nstatN{} not at all. Of the inferences, \nverdH{} are preserved outright (H), \nverdR{} in a restricted domain or on one reading (R), \nverdX{} are translated and denied (X), and \nverdU{} are untranslatable (U). The translation indices range from \timin{} (PHYS) to \timax{} (QUAN), the comparability indices from \kappamin{} (PHYS) to \kappamax{} (BIOL). The clusters together contain \nclusterall{} terms. All checks C1--C9 succeed: the conclusions of Theorems~\ref{t:partial} and~\ref{t:rational} hold in every revolution of the table, with dominant support in \ndominant{} of the ten. The philosophically interesting entries are discussed in Section~\ref{s:cases}; we note here two patterns the statistics reveal. First, the revolutions with the lowest indices, PHYS and MIAS, are those whose predecessors posited the most occult agents (spirits and faculties; miasmata and catalytic poisons); the revolutions in physics, where the predecessor is retained as a limiting case, have many R verdicts (13 of 35 in MECH) and few untranslatable terms. Second, the number of X verdicts, translated disagreements, exceeds the number of U verdicts in QUAN, BIOL and HERE, and is largest in BIOL: in these revolutions the successor contradicted its predecessor more than it abandoned it.

\section{Computed statistics}\label{s:stats}

\begin{table}[!htb]\centering\footnotesize\setlength{\tabcolsep}{3.2pt}
\caption{Per-revolution statistics computed from the manuals.}\label{tab:stats}
\begin{tabular}{@{}lrrrrrrrrrrrrr@{}}\toprule
Rev. & terms & tok. & E & P & T & N & inf. & H & R & X & U & $\tau$ & $\kappa$\\\midrule
CHEM & 23 & 35 & 8 & 11 & 4 & 0 & 31 & 20 & 0 & 2 & 9 & 0.879 & 0.710\\
ASTR & 25 & 35 & 6 & 10 & 6 & 3 & 34 & 16 & 6 & 6 & 6 & 0.887 & 0.824\\
MECH & 23 & 37 & 3 & 13 & 6 & 1 & 35 & 12 & 13 & 1 & 9 & 0.845 & 0.743\\
HEAT & 23 & 30 & 8 & 12 & 1 & 2 & 30 & 16 & 3 & 2 & 9 & 0.897 & 0.700\\
QUAN & 21 & 27 & 2 & 16 & 3 & 0 & 30 & 11 & 7 & 7 & 5 & 0.925 & 0.833\\
BIOL & 24 & 28 & 10 & 10 & 3 & 1 & 31 & 18 & 0 & 8 & 5 & 0.908 & 0.839\\
PHYS & 26 & 31 & 3 & 15 & 4 & 4 & 30 & 13 & 0 & 6 & 11 & 0.783 & 0.633\\
OPTI & 18 & 21 & 8 & 7 & 2 & 1 & 25 & 15 & 1 & 2 & 7 & 0.812 & 0.720\\
MIAS & 21 & 29 & 6 & 8 & 5 & 2 & 31 & 17 & 1 & 2 & 11 & 0.793 & 0.645\\
HERE & 17 & 21 & 6 & 7 & 2 & 2 & 23 & 11 & 3 & 5 & 4 & 0.895 & 0.826\\
\midrule
All & 221 & 294 & 60 & 109 & 36 & 16 & 300 & 149 & 34 & 41 & 76 & & \\\bottomrule\end{tabular}\end{table}

\begin{table}[!htb]\centering\footnotesize
\caption{Further indices: unweighted translation share, inferential preservation $\pi$, disagreement share $\delta$, the untranslatable cluster and the neutral core.}\label{tab:further}
\begin{tabular}{@{}lrrrp{6.2cm}r@{}}\toprule
Rev. & $\tau_{\mathrm{unw}}$ & $\pi$ & $\delta$ & untranslatable cluster $C$ & core\\\midrule
CHEM & 0.931 & 0.645 & 0.065 & CH01 (phlogiston), CH02 (dephlogisticated air), CH03 (phlogisticated air) & 3\\
ASTR & 0.840 & 0.647 & 0.176 & AS03 (primum mobile), AS04 (planetary orb), AS08 (equant), AS15 (centre of the universe), AS24 (aether) & 3\\
MECH & 0.819 & 0.714 & 0.029 & ME02 (force (impressed)), ME06 (absolute space), ME07 (absolute time), ME16 (ether), ME19 (weight), ME22 (rigid body), ME23 (action at a distance) & 2\\
HEAT & 0.891 & 0.633 & 0.067 & HE01 (caloric), HE06 (free caloric), HE07 (combined caloric) & 1\\
QUAN & 0.937 & 0.600 & 0.233 & QU01 (trajectory), QU02 (electron orbit), QU18 (state (mechanical)) & 0\\
BIOL & 0.896 & 0.581 & 0.258 & BI05 (plan of creation), BI06 (archetype), BI10 (centre of creation), BI16 (chain of being) & 7\\
PHYS & 0.769 & 0.433 & 0.200 & GA03 (vital spirit), GA04 (natural spirit), GA05 (animal spirit), GA08 (septal pores), GA14 (innate heat), GA15 (sooty vapours), GA18 (faculties), GA26 (rete mirabile) & 2\\
OPTI & 0.889 & 0.640 & 0.080 & OP02 (light corpuscle), OP03 (fits of easy reflection and transmission), OP08 (sides of rays) & 4\\
MIAS & 0.810 & 0.581 & 0.065 & MI01 (miasma), MI03 (ferment), MI05 (contagium), MI07 (spontaneous generation), MI08 (epidemic constitution), MI12 (sewer gas), MI13 (cholera poison) & 5\\
HERE & 0.824 & 0.609 & 0.217 & HR01 (gemmule), HR05 (use-inheritance), HR06 (telegony), HR13 (ancestral heredity) & 3\\
\bottomrule\end{tabular}\end{table}

\begin{table}[!htb]\centering\footnotesize
\caption{Translation index under the standard scoring S0 and the alternative schemes S1--S4.}\label{tab:sens}
\begin{tabular}{@{}lrrrrrc@{}}\toprule
Rev. & S0 & S1 & S2 & S3 & S4 & C8\\\midrule
CHEM & 0.879 & 0.741 & 1.000 & 0.588 & 0.931 & pass\\
ASTR & 0.887 & 0.560 & 0.920 & 0.667 & 0.840 & pass\\
MECH & 0.845 & 0.724 & 0.983 & 0.509 & 0.819 & pass\\
HEAT & 0.897 & 0.838 & 0.956 & 0.566 & 0.891 & pass\\
QUAN & 0.925 & 0.828 & 1.000 & 0.529 & 0.937 & pass\\
BIOL & 0.908 & 0.837 & 0.980 & 0.633 & 0.896 & pass\\
PHYS & 0.783 & 0.717 & 0.850 & 0.467 & 0.769 & pass\\
OPTI & 0.812 & 0.750 & 0.875 & 0.604 & 0.889 & pass\\
MIAS & 0.793 & 0.556 & 0.933 & 0.615 & 0.810 & pass\\
HERE & 0.895 & 0.842 & 0.947 & 0.566 & 0.824 & pass\\
\bottomrule\end{tabular}\end{table}

\begin{table}[!htb]\centering\footnotesize
\caption{Outcome of each revolution on the five shared criteria ($+$ favoured the successor, $-$ the predecessor, $\pm$ mixed or contested, ? no claim made), with the mechanical checks C4 and C5.}\label{tab:crit}
\begin{tabular}{@{}lcccccccc@{}}\toprule
Rev. & acc. & cons. & scope & simp. & fruit. & C4: $0<\tau<1$ & C5: some $+$ & no $-$\\\midrule
CHEM & $+$ & ? & $\pm$ & ? & $+$ & pass (0.879) & pass (2) & yes\\
ASTR & $+$ & $-$ & $+$ & $+$ & $+$ & pass (0.887) & pass (4) & no\\
MECH & $+$ & $+$ & $+$ & $\pm$ & $+$ & pass (0.845) & pass (4) & yes\\
HEAT & $+$ & ? & $+$ & ? & $+$ & pass (0.897) & pass (3) & yes\\
QUAN & $+$ & $+$ & $+$ & ? & $+$ & pass (0.925) & pass (4) & yes\\
BIOL & ? & $-$ & $+$ & $+$ & $+$ & pass (0.908) & pass (3) & no\\
PHYS & $+$ & $+$ & $\pm$ & ? & $+$ & pass (0.783) & pass (3) & yes\\
OPTI & $+$ & ? & $+$ & ? & $+$ & pass (0.812) & pass (3) & yes\\
MIAS & $+$ & ? & $+$ & ? & $+$ & pass (0.793) & pass (3) & yes\\
HERE & $+$ & ? & $+$ & ? & $+$ & pass (0.895) & pass (3) & yes\\
\bottomrule\end{tabular}\end{table}

\begin{crittable}
CHEM & accuracy & $+$ & Gravimetric balance: the weight gained in calcination equals the weight of air absorbed (Lavoisier, 1770s); the phlogistic accounts of the gain were not quantitative.\\
CHEM & scope & $\pm$ & The oxygen theory explained combustion, calcination and respiration together; the phlogistic theory explained why metals are alike and their calxes unlike (Kuhn; Chang).\\
CHEM & fruitfulness & $+$ & The new nomenclature (1787), the conservation of weight as a working principle, and quantitative analysis, on which Dalton built.\\
\midrule
ASTR & accuracy & $+$ & Kepler's Rudolphine Tables (1627) were much more accurate than earlier tables; Kepler predicted the transit of Mercury of 1631, observed by Gassendi. (Copernicus's own tables were not more accurate than Ptolemaic ones; Kuhn 1957, 1977.)\\
ASTR & consistency & $-$ & External consistency with Aristotelian physics and cosmology favoured the geocentric system until the new mechanics (Kuhn 1977).\\
ASTR & scope & $+$ & Fixed the order and relative distances of the planets from observation and explained the link of the second anomaly to the Sun.\\
ASTR & simplicity & $+$ & With Kepler, one ellipse per planet governed by the area law replaced compounded circles and equants.\\
ASTR & fruitfulness & $+$ & Led to Kepler's laws and, through them, to Newton's mechanics.\\
\midrule
MECH & accuracy & $+$ & General relativity accounted for the anomalous 43 arcseconds per century of Mercury's perihelion advance and predicted the deflection of starlight reported in 1919.\\
MECH & consistency & $+$ & Special relativity removed the conflict between the Galilean invariance of mechanics and the Lorentz invariance of Maxwell's equations.\\
MECH & scope & $+$ & Mechanics, electrodynamics and optics of moving bodies brought under one kinematics; gravitation and inertia unified.\\
MECH & simplicity & $\pm$ & Conceptual unity was gained; the mathematics of general relativity is far more complex.\\
MECH & fruitfulness & $+$ & Mass--energy equivalence, relativistic cosmology, and the relativistic theories of the electron and of fields.\\
\midrule
HEAT & accuracy & $+$ & Joule's determinations of the mechanical equivalent of heat (1840s) showed a fixed ratio of work spent to heat produced, which the conservation of caloric excluded.\\
HEAT & scope & $+$ & The kinetic theory derived the gas laws and predicted that the viscosity of a gas is independent of its pressure (Maxwell, 1860), which Maxwell confirmed experimentally (1866).\\
HEAT & fruitfulness & $+$ & The first and second laws, statistical mechanics, and the theory of transport processes.\\
\midrule
QUAN & accuracy & $+$ & Planck's law (1900) fitted the black-body spectrum where the Rayleigh--Jeans law failed; Einstein's theory of specific heats (1907) explained their fall at low temperature; the hydrogen spectrum was derived.\\
QUAN & consistency & $+$ & The classical nuclear atom is unstable by classical electrodynamics; quantum mechanics yields stable stationary states.\\
QUAN & scope & $+$ & Chemical bonding, the periodic table, spectra and the solid state.\\
QUAN & fruitfulness & $+$ & Quantum chemistry, solid-state physics, nuclear and particle physics.\\
\midrule
BIOL & consistency & $-$ & External: Kelvin's estimates of the age of the Earth (1860s) seemed to leave too little time, and Jenkin (1867) argued that blending inheritance would swamp favourable variations.\\
BIOL & scope & $+$ & One hypothesis explained the natural system, homology, rudimentary organs, embryonic resemblance, the succession of types and geographical distribution (Origin, chs.~11--14).\\
BIOL & simplicity & $+$ & One process of descent with modification replaced a separate act of creation for each species.\\
BIOL & fruitfulness & $+$ & Phylogenetic systematics, biogeography and, after the synthesis with genetics, population biology.\\
\midrule
PHYS & accuracy & $+$ & Harvey's quantitative argument: the blood expelled by the heart in a short time far exceeds what food could supply, so the same blood must return (De motu cordis, ch.~9).\\
PHYS & consistency & $+$ & The circulation required no invisible septal pores, whose existence Vesalius had come to doubt in the 1555 edition of the Fabrica.\\
PHYS & scope & $\pm$ & Harvey explained the motion of the blood and the function of the venous valves; the Galenic system offered accounts of nutrition, respiration and heat that the circulation left open.\\
PHYS & fruitfulness & $+$ & Malpighi's capillaries (1661), the Oxford physiologists' work on respiration, transfusion experiments.\\
\midrule
OPTI & accuracy & $+$ & Fresnel's integrals gave the positions and intensities of diffraction fringes; Arago observed the bright spot at the centre of a disk's shadow that Poisson had derived from Fresnel's theory (1819).\\
OPTI & scope & $+$ & Interference, diffraction, polarization and double refraction brought under transverse waves.\\
OPTI & fruitfulness & $+$ & Fresnel's equations for reflection and refraction; the electromagnetic theory of light.\\
\midrule
MIAS & accuracy & $+$ & Specific organisms were isolated and shown to cause specific diseases (Koch on anthrax, 1876; tuberculosis, 1882; the cholera vibrio, 1884), and the routes of transmission were traced.\\
MIAS & scope & $+$ & Fermentation, putrefaction, wound infection and the specificity of epidemic diseases explained by one hypothesis.\\
MIAS & fruitfulness & $+$ & Antiseptic surgery (Lister, 1867), attenuated vaccines (Pasteur on anthrax, 1881; rabies, 1885).\\
\midrule
HERE & accuracy & $+$ & The Mendelian ratios (3:1, 9:3:3:1) were confirmed in crosses after 1900 (Correns, de Vries, Tschermak, Bateson).\\
HERE & scope & $+$ & Segregation, independent assortment, linkage and sex-linked inheritance (Morgan, 1910) brought under one chromosomal mechanism.\\
HERE & fruitfulness & $+$ & The Hardy--Weinberg principle (1908) and population genetics; Fisher (1918) derived the biometric correlations from Mendelian factors.\\
\bottomrule
\end{crittable}

\clearpage

\section{Evidential claims}\label{s:criteria}

Table~\ref{tab:evid} lists the evidential claims behind the favourable outcomes of Table~\ref{tab:crit}, with the readings each uses; check C9 verifies that every reading refers.

\begin{evidtable}
CH.E1 & Tin calcined in a sealed retort gains weight only when the retort is opened and air rushes in, by the weight of the air admitted (Lavoisier, 1774). & CH09, CH07, CH06, CH11 & accuracy\\
CH.E2 & The red calx of mercury heated alone gives back the metal and an air that supports a flame better than common air (Priestley; Lavoisier). & CH06, CH07, CH02a, CH08, CH11 & accuracy, fruitfulness\\
CH.E3 & Inflammable air exploded with dephlogisticated air leaves water, nearly the whole of both airs being consumed (Cavendish, 1784). & CH04a, CH02c, CH21 & accuracy, fruitfulness\\
CH.E4 & Water passed over red-hot iron yields inflammable air, and the iron becomes a calx heavier than the metal (Lavoisier and Meusnier, 1784). & CH21, CH04a, CH06, CH07 & accuracy\\
\midrule
AS.E1 & Mars's positions observed by Tycho Brahe are fitted by Kepler's ellipse to within the errors of observation, after the best circular model failed by about 8 minutes of arc. & AS01a, AS13 & accuracy, simplicity\\
AS.E2 & The transit of Mercury of 1631 occurred close to the time predicted from the Rudolphine Tables (Gassendi). & AS01d, AS13 & accuracy\\
AS.E3 & The ratio of epicycle to deferent radius for each planet yields, on the Copernican arrangement, the relative distances of the planets. & AS06a, AS05a, AS06b, AS05b, AS22 & scope\\
AS.E4 & Superior planets retrograde only at opposition and inferior planets only near inferior conjunction. & AS11, AS01a, AS01d, AS13 & scope\\
AS.E5 & Kepler's third law relates the periods and the mean distances of all the planets. & AS21a, AS22 & fruitfulness, simplicity\\
\midrule
ME.E1 & The perihelion of Mercury advances by about 43 seconds of arc per century more than the perturbations of the other planets account for (Le Verrier; Newcomb). & ME13, ME11 & accuracy\\
ME.E2 & Starlight passing near the Sun is deflected (eclipse expeditions of 1919). & ME15, ME11 & accuracy\\
ME.E3 & No effect of the Earth's motion on the round-trip time of light in perpendicular arms is detected (Michelson and Morley, 1887). & ME15, ME05 & consistency, scope\\
ME.E4 & The mass of fast cathode-ray electrons, measured by deflection, increases with velocity (Kaufmann; Bucherer, 1908). & ME01, ME05 & scope, fruitfulness\\
\midrule
HE.E1 & Boring a cannon with a blunt tool heats the metal and its water bath without limit while the boring continues (Rumford, 1798). & HE15, HE02 & accuracy\\
HE.E2 & A fixed quantity of work spent in friction or in stirring water produces a fixed quantity of heat (Joule, 1840s). & HE15, HE02, HE03 & accuracy, fruitfulness\\
HE.E3 & The viscosity of a gas is independent of its density over a wide range (Maxwell, 1866). & HE18, HE19 & scope\\
HE.E4 & Air heats on compression and cools on expansion, and the speed of sound requires the ratio of specific heats. & HE10, HE11, HE12 & scope\\
\midrule
QU.E1 & The spectral density of cavity radiation falls at high frequency, contrary to the Rayleigh--Jeans law (Lummer and Pringsheim; Rubens and Kurlbaum, 1899--1900). & QU13, QU12 & accuracy\\
QU.E2 & The atomic heats of diamond and other solids fall far below the Dulong--Petit value at low temperatures (Weber; Nernst). & QU14, QU12 & accuracy, scope\\
QU.E3 & The maximum energy of photoelectrons depends on the frequency of the light, not on its intensity (Lenard, 1902; Millikan, 1916). & QU15, QU09, QU05 & accuracy, fruitfulness\\
QU.E4 & Atoms are stable and emit sharp spectral lines. & QU07, QU11 & consistency, scope\\
QU.E5 & A beam of silver atoms is split into two discrete beams by an inhomogeneous magnetic field (Stern and Gerlach, 1922). & QU17, QU16 & accuracy\\
\midrule
BI.E1 & The species of the Galapagos are allied to those of the nearest mainland, South America, and differ from island to island. & BI01, BI10b & scope\\
BI.E2 & Homologous bones underlie the forelimbs of man, bat and porpoise. & BI07 & scope, simplicity\\
BI.E3 & Rudimentary teeth occur in foetal whales and are absorbed before birth. & BI11 & scope\\
BI.E4 & Extinct edentates of South America are allied to the living armadillos and sloths of the same continent. & BI14, BI01 & scope, fruitfulness\\
BI.E5 & Domestic breeds descended from one wild species differ more than many wild species do (pigeons). & BI17, BI01, BI23 & simplicity\\
\midrule
GA.E1 & The left ventricle holds a measurable quantity of blood, and the heart beats over a thousand times in half an hour (De motu cordis, ch.~9). & GA07, GA12 & accuracy\\
GA.E2 & A ligature tight enough to stop the arteries makes the hand pale and cold; loosened enough to let the arteries pass, it makes the veins below it swell (De motu cordis, ch.~11). & GA13, GA02, GA01 & accuracy\\
GA.E3 & Blood pressed in a vein of the arm moves readily toward the heart and cannot be pushed past the valves away from it (De motu cordis, ch.~13). & GA17, GA01 & accuracy, fruitfulness\\
GA.E4 & No pores can be found in the interventricular septum on dissection. & GA07 & consistency\\
GA.E5 & Under the microscope blood is seen passing from arteries to veins in the lung of the frog (Malpighi, 1661). & GA22, GA02, GA01 & fruitfulness\\
\midrule
OP.E1 & A bright spot appears at the centre of the shadow of a small opaque disk (Arago, 1819). & OP16, OP07 & accuracy\\
OP.E2 & The positions of the fringes outside the shadow of an edge vary with distance as Fresnel's integrals predict. & OP07, OP16 & accuracy, scope\\
OP.E3 & Two beams polarized at right angles do not form fringes when recombined (Fresnel and Arago, 1819). & OP09 & scope\\
OP.E4 & Light travels more slowly in water than in air (Foucault, 1850). & OP06 & accuracy, fruitfulness\\
\midrule
MI.E1 & Broth boiled in a swan-necked flask stays unchanged until the neck is broken (Pasteur, 1859--1862). & MI06 & scope\\
MI.E2 & The blood of animals dead of anthrax contains rod-shaped bodies, cultures of which reproduce the disease (Koch, 1876). & MI18, MI05b & accuracy\\
MI.E3 & Cholera deaths in 1854 clustered about the Broad Street pump; customers of the Southwark and Vauxhall Company died at a far higher rate than those of the Lambeth Company (Snow, 1855). & MI21, MI13a & accuracy\\
MI.E4 & Compound fractures dressed with carbolic acid rarely suppurated (Lister, 1867). & MI16, MI14 & fruitfulness\\
MI.E5 & Sheep vaccinated with attenuated anthrax cultures survived inoculation that killed unvaccinated sheep (Pouilly-le-Fort, 1881). & MI15, MI05b & fruitfulness\\
\midrule
HR.E1 & In the second generation from crosses of true-breeding peas, the dominant and recessive forms appear in the ratio of about 3:1. & HR07, HR08, HR11, HR03 & accuracy\\
HR.E2 & Characters absent in the first-generation hybrids reappear unchanged in the second. & HR08, HR16, HR11 & accuracy, scope\\
HR.E3 & Selection within a pure line of beans does not shift the mean seed size of its progeny (Johannsen, 1903). & HR07, HR10 & scope\\
HR.E4 & In Drosophila the white-eye character is transmitted with sex (Morgan, 1910). & HR11, HR16 & scope, fruitfulness\\
\bottomrule
\end{evidtable}


\section{The manuals}\label{s:manuals}

\Needspace{10\baselineskip}\subsection*{CHEM: Chemical revolution}
\addcontentsline{toc}{subsection}{CHEM: Chemical revolution}
\noindent\textit{Predecessor} $T_1$: Phlogistic chemistry (Stahl; Black, Priestley, Cavendish, Scheele, Kirwan), c.~1700--1787.\\
\textit{Successor} $T_2$: Oxygen chemistry (Lavoisier, Trait\'{e} \'{e}l\'{e}mentaire de chimie, 1789) and its continuation.\par\smallskip
\begin{manualtable}
\textbf{CH01} & \textbf{phlogiston}. Principle of inflammability, contained in all combustibles and metals and released into the air in combustion and calcination. & --- & \textbf{T} token-dependent; $w$\,=\,10; $s$\,=\,1/4 & by token \\
\quad CH01a & Stahl and his followers: the principle common to all combustibles, given off when they burn & \textit{no referent} & fails & U: I01 I04 I09 I12 I20 I26 \\
\quad CH01b & In reduction: what a calx takes from charcoal when heated with it to restore the metal & \textit{no referent} & fails & U: I08 \\
\quad CH01c & Kirwan (Essay on Phlogiston, 1787) and some statements of Cavendish: phlogiston identified with, or as the main constituent of, the inflammable air obtained from metals and acids & hydrogen & refers & H: I28 \\
\quad CH01d & What air absorbs from burning bodies until it is saturated and combustion stops & \textit{no referent} & fails & U: I02 I22 \\
\midrule
\textbf{CH02} & \textbf{dephlogisticated air}. Air deprived of its phlogiston and therefore unusually avid of it; supports combustion and respiration better than common air. & --- & \textbf{T} token-dependent; $w$\,=\,4; $s$\,=\,2/3 & by token \\
\quad CH02a & Priestley, 1774--1775: the air obtained by heating the red calx of mercury, in which a candle burned with a remarkably vigorous flame & oxygen & refers & H: I10 I11 \\
\quad CH02b & The type description: common air from which phlogiston has been removed & \textit{no referent} & fails & U: I12 \\
\quad CH02c & Cavendish, 1784: the air exploded with inflammable air in his synthesis of water & oxygen & refers & H: I18 \\
\midrule
\textbf{CH03} & \textbf{phlogisticated air}. Air saturated with phlogiston, which no longer supports combustion or respiration. & --- & \textbf{T} token-dependent; $w$\,=\,2; $s$\,=\,1/2 & by token \\
\quad CH03a & Residue of air in which charcoal or a candle has burned, after the fixed air has been absorbed (Rutherford, 1772; Priestley) & nitrogen (with traces of the noble gases, then unknown) & refers & H: I14 \\
\quad CH03b & The type description: common air that has taken up as much phlogiston as it can hold & \textit{no referent} & fails & U: I02 \\
\midrule
\textbf{CH04} & \textbf{inflammable air}. An air that burns; for Cavendish (1766) the air obtained from metals dissolved in acids. & --- & \textbf{T} token-dependent; $w$\,=\,5; $s$\,=\,1 & by token \\
\quad CH04a & Cavendish, 1766: air from iron, zinc or tin dissolved in dilute vitriolic or marine acid & hydrogen & refers & H: I17 I18 I28 \\
\quad CH04b & Volta, 1776: inflammable air collected from the mud of marshes & methane & refers & H: I29 \\
\quad CH04c & The heavy inflammable air from heating charcoal with metallic calxes (Priestley) & carbon monoxide & refers & H: I30 \\
\midrule
\textbf{CH05} & \textbf{fixed air}. Air fixed in chalk and magnesia alba and released on heating or by acids; extinguishes flame. & carbon dioxide (carbonic acid gas) & \textbf{E} exact; $w$\,=\,3; $s$\,=\,1 & H: I14 I15 I16 \\
\quad CH05a & Black, 1756: the air expelled from magnesia alba and chalk by heat & carbon dioxide & refers & as term \\
\quad CH05b & Air given off in fermentation and exhaled in respiration & carbon dioxide & refers & as term \\
\midrule
\textbf{CH06} & \textbf{calx}. The earthy residue of the calcination of a metal; the metal deprived of its phlogiston. & metallic oxide & \textbf{P} partial; $w$\,=\,6; $s$\,=\,1 & H: I06 I07 I10 I30; U: I04 I08 \\
\quad CH06a & Red calx of mercury (mercurius calcinatus per se) & mercuric oxide & refers & as term \\
\quad CH06b & Calx of lead (litharge, minium) & lead oxides & refers & as term \\
\quad CH06c & Calx of tin formed in a sealed retort (Lavoisier, 1774) & tin oxide & refers & as term \\
\midrule
\textbf{CH07} & \textbf{metal}. A compound of a calx with phlogiston; the common properties of metals are due to their common phlogiston. & metal (a simple substance) & \textbf{P} partial; $w$\,=\,7; $s$\,=\,1 & H: I05 I06 I07 I10 I17; U: I04 I26 \\
\quad CH07a & Mercury, lead, tin, iron, zinc as named in calcination experiments & the same metals & refers & as term \\
\quad CH07b & Metals as a class sharing lustre and malleability & the metals & refers & as term \\
\midrule
\textbf{CH08} & \textbf{combustion}. Release of phlogiston from a combustible into the surrounding air. & combustion (combination with oxygen) & \textbf{P} partial; $w$\,=\,10; $s$\,=\,1 & H: I03 I11 I13 I14 I16 I23 I31; U: I01 I02 I09 \\
\quad CH08a & Burning of a candle, charcoal or sulphur in a closed volume of air & combustion in air & refers & as term \\
\midrule
\textbf{CH09} & \textbf{calcination}. Conversion of a metal into its calx by the loss of phlogiston on heating in air. & oxidation of a metal & \textbf{E} exact; $w$\,=\,2; $s$\,=\,1 & H: I05 I06 \\
\quad CH09a & Heating of tin or lead in air until it becomes a calx & oxidation of tin or lead & refers & as term \\
\midrule
\textbf{CH10} & \textbf{reduction (revivification)}. Restoration of a metal from its calx by restoring phlogiston, usually from charcoal. & reduction (removal of oxygen) & \textbf{P} partial; $w$\,=\,2; $s$\,=\,1 & H: I07; U: I08 \\
\quad CH10a & Heating a calx with charcoal to recover the metal & reduction of the oxide by carbon & refers & as term \\
\midrule
\textbf{CH11} & \textbf{common air}. The atmosphere, treated as one kind of air whose goodness varies with the phlogiston it holds. & atmospheric air (a mixture chiefly of oxygen and nitrogen) & \textbf{P} partial; $w$\,=\,6; $s$\,=\,1 & H: I03 I05 I11 I13; X: I27; U: I01 \\
\quad CH11a & The air of a room or of the open atmosphere used in eudiometry & atmospheric air & refers & as term \\
\midrule
\textbf{CH12} & \textbf{goodness of air}. Fitness of an air for respiration, inversely related to its content of phlogiston; measured by the nitrous air test. & proportion of oxygen in an air & \textbf{P} partial; $w$\,=\,2; $s$\,=\,1 & H: I21; U: I22 \\
\quad CH12a & Priestley's nitrous air test: the diminution of volume when nitrous air is mixed with the sample over water & oxygen content of the sample & refers & as term \\
\midrule
\textbf{CH13} & \textbf{nitrous air}. Air obtained from metals and spirit of nitre; diminishes good air when mixed with it over water. & nitric oxide & \textbf{E} exact; $w$\,=\,1; $s$\,=\,1 & H: I21 \\
\quad CH13a & Priestley, 1772: air from copper or brass dissolved in spirit of nitre & nitric oxide & refers & as term \\
\midrule
\textbf{CH14} & \textbf{dephlogisticated nitrous air}. Air obtained by exposing nitrous air to iron and other substances; supports a candle flame. & nitrous oxide & \textbf{E} exact; $w$\,=\,1; $s$\,=\,1 & H: I23 \\
\quad CH14a & Priestley's preparation; later the subject of Davy's researches (1800) & nitrous oxide & refers & as term \\
\midrule
\textbf{CH15} & \textbf{marine acid air}. Air obtained from spirit of salt; very soluble in water. & hydrogen chloride & \textbf{E} exact; $w$\,=\,1; $s$\,=\,1 & H: I24 \\
\quad CH15a & Priestley: air collected over mercury from heated spirit of salt & hydrogen chloride & refers & as term \\
\midrule
\textbf{CH16} & \textbf{dephlogisticated marine acid}. Marine acid deprived of phlogiston by manganese (Scheele, 1774); bleaches vegetable colours. & oxygenated muriatic acid (Lavoisierian name; now chlorine) & \textbf{E} exact; $w$\,=\,1; $s$\,=\,1 & H: I25 \\
\quad CH16a & Scheele: the yellow-green air from marine acid and manganese & chlorine & refers & as term \\
\midrule
\textbf{CH17} & \textbf{alkaline air}. Air obtained from sal ammoniac and lime; very soluble in water, alkaline. & ammoniacal gas & \textbf{E} exact; $w$\,=\,1; $s$\,=\,1 & H: I24 \\
\quad CH17a & Priestley: air collected over mercury from sal ammoniac and slaked lime & ammonia & refers & as term \\
\midrule
\textbf{CH18} & \textbf{vitriolic acid}. Oil of vitriol; in Stahl's chemistry, the acid which with phlogiston composes sulphur. & sulphuric acid & \textbf{P} partial; $w$\,=\,2; $s$\,=\,1 & H: I17; U: I20 \\
\quad CH18a & Oil of vitriol used to dissolve metals & sulphuric acid & refers & as term \\
\midrule
\textbf{CH19} & \textbf{sulphur}. A compound of vitriolic acid and phlogiston (Stahl). & sulphur (a simple substance) & \textbf{P} partial; $w$\,=\,2; $s$\,=\,1 & H: I31; U: I20 \\
\quad CH19a & Common brimstone burned in air & sulphur & refers & as term \\
\midrule
\textbf{CH20} & \textbf{charcoal}. A substance very rich in phlogiston, burning almost without residue. & carbon (as charcoal) & \textbf{P} partial; $w$\,=\,5; $s$\,=\,1 & H: I07 I16 I30; U: I08 I09 \\
\quad CH20a & Charcoal used in reduction and in combustion experiments & impure carbon & refers & as term \\
\midrule
\textbf{CH21} & \textbf{water}. An element; in Cavendish's 1784 interpretation, the base of both dephlogisticated and inflammable air. & water (a compound of hydrogen and oxygen) & \textbf{P} partial; $w$\,=\,3; $s$\,=\,1 & H: I18 I31; X: I19 \\
\quad CH21a & The dew deposited in Cavendish's exploding vessel & water & refers & as term \\
\midrule
\textbf{CH22} & \textbf{element}. One of the few ultimate constituents (earth, water, air, fire, or Stahl's principles), not identified by analysis. & simple substance (Lavoisier's operational sense: not yet decomposed) & \textbf{P} partial; $w$\,=\,2; $s$\,=\,1 & X: I19 I27 \\
\quad CH22a & Water and air counted among the elements & simple substance (the predicate) & refers & as term \\
\midrule
\textbf{CH23} & \textbf{respiration}. Discharge of phlogiston from the lungs into the air. & respiration (slow combustion consuming oxygen) & \textbf{E} exact; $w$\,=\,3; $s$\,=\,1 & H: I11 I13 I16 \\
\quad CH23a & Mice and men breathing confined air until it is unfit & respiration & refers & as term \\
\bottomrule
\end{manualtable}

\Needspace{10\baselineskip}\subsection*{ASTR: Astronomical revolution}
\addcontentsline{toc}{subsection}{ASTR: Astronomical revolution}
\noindent\textit{Predecessor} $T_1$: Ptolemaic astronomy (Almagest, c.~150; Planetary Hypotheses) with Aristotelian cosmology, as taught to 1543.\\
\textit{Successor} $T_2$: Copernican--Keplerian astronomy (De revolutionibus, 1543; Astronomia nova, 1609; Rudolphine Tables, 1627).\par\smallskip
\begin{manualtable}
\textbf{AS01} & \textbf{planet}. A wandering star: Moon, Mercury, Venus, Sun, Mars, Jupiter, Saturn. & --- & \textbf{T} token-dependent; $w$\,=\,13; $s$\,=\,1 & by token \\
\quad AS01a & Mars, Jupiter or Saturn called a planet & the same body, a planet & refers & H: I02 I03 I09 I21; R: I01 I05; U: I11 \\
\quad AS01b & The Sun called a planet & the Sun, a star at the centre of the system & refers & X: I20 I34 \\
\quad AS01c & The Moon called a planet & the Moon, a satellite of the Earth & refers & H: I26 \\
\quad AS01d & Venus or Mercury called a planet & the same body, a planet & refers & H: I19 I22; R: I04 \\
\midrule
\textbf{AS02} & \textbf{fixed star}. A star keeping its position relative to the other stars; carried round daily by the heavens. & star (beyond the planetary system) & \textbf{P} partial; $w$\,=\,3; $s$\,=\,1 & H: I09 I13; X: I14 \\
\quad AS02a & The stars catalogued in Almagest VII--VIII & the same stars & refers & as term \\
\midrule
\textbf{AS03} & \textbf{primum mobile}. The outermost moving sphere, which turns daily from east to west and carries the inner spheres with it. & --- & \textbf{N} none; $w$\,=\,1; $s$\,=\,0 & by token \\
\quad AS03a & The sphere that communicates the daily motion to all the heavens & \textit{no referent} & fails & U: I08 \\
\midrule
\textbf{AS04} & \textbf{planetary orb}. A solid spherical shell of aether carrying a planet; the orbs nest without gaps (Planetary Hypotheses; Peurbach's total and partial orbs). & --- & \textbf{N} none; $w$\,=\,4; $s$\,=\,0 & by token \\
\quad AS04a & The orb of Mars as a solid shell & \textit{no referent} & fails & U: I08 I11 I33 \\
\quad AS04b & The orbs whose nesting fixes absolute distances in the Planetary Hypotheses & \textit{no referent} & fails & U: I29 \\
\midrule
\textbf{AS05} & \textbf{deferent}. The large circle about the Earth (or an eccentric point) on which the centre of a planet's epicycle moves. & --- & \textbf{T} token-dependent; $w$\,=\,2; $s$\,=\,1 & by token \\
\quad AS05a & Deferent of a superior planet (Mars, Jupiter, Saturn) & the planet's orbit about the Sun & refers & R: I01 \\
\quad AS05b & Deferent of an inferior planet (Venus, Mercury), whose epicycle centre stays on the line to the mean Sun & the Earth's orbit about the Sun (seen from the Earth as the Sun's annual motion) & refers & R: I04 \\
\midrule
\textbf{AS06} & \textbf{epicycle}. A circle whose centre moves on the deferent and on which the planet moves; accounts for the second anomaly. & --- & \textbf{T} token-dependent; $w$\,=\,5; $s$\,=\,1 & by token \\
\quad AS06a & Main epicycle of a superior planet & the Earth's orbit about the Sun, as reflected in the planet's geocentric motion & refers & R: I01 I05 I15 \\
\quad AS06b & Main epicycle of Venus or Mercury & the planet's orbit about the Sun & refers & R: I04 \\
\quad AS06c & Epicycle of the Moon (first lunar model) & the eccentricity of the Moon's orbit about the Earth (the equation of centre) & refers & H: I28 \\
\midrule
\textbf{AS07} & \textbf{eccentric}. A circle whose centre is displaced from the Earth. & eccentricity of an orbit (displacement of the Sun from its centre) & \textbf{P} partial; $w$\,=\,2; $s$\,=\,1 & R: I16 I25 \\
\quad AS07a & Eccentric of the Sun's circle (Hipparchus, Almagest III) & the eccentricity of the Earth's orbit & refers & as term \\
\midrule
\textbf{AS08} & \textbf{equant}. The point about which the centre of the epicycle moves with uniform angular speed. & --- & \textbf{T} token-dependent; $w$\,=\,3; $s$\,=\,1/2 & by token \\
\quad AS08a & Equant of a superior planet, at twice the eccentricity of the deferent's centre & the empty focus of the planet's ellipse (to first order in the eccentricity) & refers & R: I15 I16 \\
\quad AS08b & The equant as the centre of a physically uniform rotation of the orb & \textit{no referent} & fails & U: I33 \\
\midrule
\textbf{AS09} & \textbf{first anomaly}. The non-uniformity of a planet's motion that depends on its position in the zodiac. & non-uniform orbital motion due to eccentricity (Kepler's second law) & \textbf{E} exact; $w$\,=\,1; $s$\,=\,1 & H: I18 \\
\quad AS09a & The zodiacal anomaly of Mars & the non-uniform heliocentric motion of Mars & refers & as term \\
\midrule
\textbf{AS10} & \textbf{second anomaly}. The non-uniformity, including retrogradation, that depends on the planet's position relative to the Sun. & the apparent effect of the Earth's orbital motion & \textbf{E} exact; $w$\,=\,2; $s$\,=\,1 & H: I17 I31 \\
\quad AS10a & The synodic anomaly of Jupiter & the parallax of the Earth's orbit in Jupiter's direction & refers & as term \\
\midrule
\textbf{AS11} & \textbf{retrograde motion}. A real backward motion of the planet, produced by its motion on the epicycle. & apparent retrograde motion & \textbf{P} partial; $w$\,=\,3; $s$\,=\,1 & H: I02 I03; X: I30 \\
\quad AS11a & Retrogradation of Mars near opposition & apparent retrograde motion of Mars & refers & as term \\
\midrule
\textbf{AS12} & \textbf{Earth}. The stationary body at the centre of the universe. & the Earth (a planet) & \textbf{P} partial; $w$\,=\,4; $s$\,=\,1 & R: I04 I05 I25; X: I06 \\
\quad AS12a & The Earth on which observers stand & the Earth & refers & as term \\
\midrule
\textbf{AS13} & \textbf{Sun}. A planet moving about the Earth on an eccentric circle. & the Sun & \textbf{P} partial; $w$\,=\,9; $s$\,=\,1 & H: I02 I22 I24 I27; R: I04 I05 I25; X: I20 I34 \\
\quad AS13a & The Sun whose annual motion defines the ecliptic & the Sun & refers & as term \\
\midrule
\textbf{AS14} & \textbf{Moon}. The planet nearest the Earth. & the Moon (the Earth's satellite) & \textbf{P} partial; $w$\,=\,5; $s$\,=\,1 & H: I26 I28 I32; X: I20 I23 \\
\quad AS14a & The Moon of eclipses and lunar theory & the Moon & refers & as term \\
\midrule
\textbf{AS15} & \textbf{centre of the universe}. The place to which heavy bodies tend and about which the heavens rotate; the centre of the Earth. & --- & \textbf{T} token-dependent; $w$\,=\,2; $s$\,=\,1/2 & by token \\
\quad AS15a & The place to which heavy bodies naturally move & the centre of the Earth & refers & H: I07 \\
\quad AS15b & The fixed point about which the heavens make their daily rotation & \textit{no referent} & fails & U: I10 \\
\midrule
\textbf{AS16} & \textbf{diurnal motion}. The daily rotation of the whole heavens from east to west. & apparent diurnal motion (due to the Earth's rotation) & \textbf{P} partial; $w$\,=\,3; $s$\,=\,1 & H: I09; U: I08 I10 \\
\quad AS16a & The daily rising and setting of all celestial bodies & apparent diurnal motion & refers & as term \\
\midrule
\textbf{AS17} & \textbf{ecliptic}. The great circle traced annually by the Sun among the fixed stars. & ecliptic (projection of the plane of the Earth's orbit) & \textbf{E} exact; $w$\,=\,2; $s$\,=\,1 & H: I27 I32 \\
\quad AS17a & The circle of the zodiac in which eclipses occur & ecliptic & refers & as term \\
\midrule
\textbf{AS18} & \textbf{tropical year}. The period of the Sun's return to the same equinox. & tropical year & \textbf{E} exact; $w$\,=\,1; $s$\,=\,1 & H: I24 \\
\quad AS18a & Ptolemy's value, slightly short of 365 1/4 days & tropical year & refers & as term \\
\midrule
\textbf{AS19} & \textbf{precession of the equinoxes}. A slow eastward motion of the sphere of fixed stars relative to the equinoxes. & precession (slow motion of the Earth's axis) & \textbf{P} partial; $w$\,=\,2; $s$\,=\,1 & H: I24; X: I14 \\
\quad AS19a & Hipparchus' and Ptolemy's precession & precession & refers & as term \\
\midrule
\textbf{AS20} & \textbf{synodic period}. The period of the second anomaly: return to the same configuration with the Sun. & synodic period & \textbf{E} exact; $w$\,=\,1; $s$\,=\,1 & H: I17 \\
\quad AS20a & Synodic period of Mars (about 780 days) & synodic period of Mars & refers & as term \\
\midrule
\textbf{AS21} & \textbf{zodiacal period}. The period of the mean motion of the epicycle centre around the zodiac. & --- & \textbf{T} token-dependent; $w$\,=\,2; $s$\,=\,1 & by token \\
\quad AS21a & Zodiacal period of a superior planet (Saturn about 29 1/2 years) & sidereal period of the planet & refers & H: I21 \\
\quad AS21b & Zodiacal period of Venus or Mercury (one year) & sidereal period of the Earth & refers & H: I19 \\
\midrule
\textbf{AS22} & \textbf{order of the planets}. The order by distance from the Earth, conventionally fixed by period: Moon, Mercury, Venus, Sun, Mars, Jupiter, Saturn. & order of distances (from the Earth or the Sun) & \textbf{P} partial; $w$\,=\,1; $s$\,=\,1 & X: I20 \\
\quad AS22a & The order stated in Almagest IX.1 & order of distances & refers & as term \\
\midrule
\textbf{AS23} & \textbf{elongation}. Angular distance of a planet from the Sun. & elongation & \textbf{E} exact; $w$\,=\,2; $s$\,=\,1 & H: I22 I31 \\
\quad AS23a & Greatest elongations of Venus and Mercury & elongation & refers & as term \\
\midrule
\textbf{AS24} & \textbf{aether}. The fifth element, of which the heavens are made: ungenerated, unalterable, moving in circles. & --- & \textbf{N} none; $w$\,=\,1; $s$\,=\,0 & by token \\
\quad AS24a & The substance of the orbs and stars & \textit{no referent} & fails & U: I12 \\
\midrule
\textbf{AS25} & \textbf{comet}. A burning exhalation in the upper sublunary region. & comet (a body beyond the Moon) & \textbf{P} partial; $w$\,=\,1; $s$\,=\,1 & X: I23 \\
\quad AS25a & The comet of 1577 & the comet of 1577 & refers & as term \\
\bottomrule
\end{manualtable}

\Needspace{10\baselineskip}\subsection*{MECH: Relativistic revolution}
\addcontentsline{toc}{subsection}{MECH: Relativistic revolution}
\noindent\textit{Predecessor} $T_1$: Newtonian mechanics and gravitation (Principia, 1687), with the ether-based electrodynamics of the late nineteenth century.\\
\textit{Successor} $T_2$: Special relativity (Einstein, 1905) and general relativity (Einstein, 1915).\par\smallskip
\begin{manualtable}
\textbf{ME01} & \textbf{mass}. Quantity of matter; invariant, additive and conserved; measures inertia and gravitational charge. & rest mass $\mid$ relativistic mass (total energy divided by $c^2$) (pd) & \textbf{P} partial; $w$\,=\,7; $s$\,=\,1 & H: I22; R: I01 I02 I03 I04 I26; U: I07 \\
\quad ME01a & Mass of a planet in the Kepler problem & rest mass $\approx$ relativistic mass (slow motion) & refers & as term \\
\quad ME01b & Mass of a fast electron in the deflection experiments of Kaufmann (1901--1906) & relativistic mass, or rest mass, as the reading is chosen & refers & as term \\
\quad ME01c & Mass of a composite body as the sum of the masses of its parts & relativistic mass (total energy over $c^2$) & refers & as term \\
\midrule
\textbf{ME02} & \textbf{force (impressed)}. Action exerted on a body changing its state of motion; equal to the rate of change of momentum. & --- & \textbf{T} token-dependent; $w$\,=\,3; $s$\,=\,2/3 & by token \\
\quad ME02a & Electric force on a charged particle & the Lorentz force (rate of change of relativistic momentum) & refers & H: I06 \\
\quad ME02b & Gravitational force of the Sun on a planet & \textit{no referent} & fails & U: I07 I28 \\
\quad ME02c & Contact force of a support or a spring & the same force (four-force in the rest frame) & refers & (none listed) \\
\midrule
\textbf{ME03} & \textbf{momentum (quantity of motion)}. Product of mass and velocity; conserved in isolated systems. & $m_0 v$ (low velocities) $\mid$ relativistic momentum $\gamma m_0 v$ (pd) & \textbf{P} partial; $w$\,=\,3; $s$\,=\,1 & H: I05 I06; R: I04 \\
\quad ME03a & Momentum of colliding billiard balls & $m_0 v \approx \gamma m_0 v$ & refers & as term \\
\quad ME03b & Momentum of a cathode-ray electron & relativistic momentum & refers & as term \\
\midrule
\textbf{ME04} & \textbf{kinetic energy (vis viva)}. One half of mass times the square of velocity. & kinetic energy $(\gamma-1)m_0c^2$ & \textbf{P} partial; $w$\,=\,1; $s$\,=\,1 & R: I26 \\
\quad ME04a & Energy of a falling body & kinetic energy & refers & as term \\
\midrule
\textbf{ME05} & \textbf{velocity}. Rate of change of position, relative to a frame; composes by vector addition. & velocity relative to an inertial frame & \textbf{P} partial; $w$\,=\,4; $s$\,=\,1 & R: I04 I10 I26; X: I19 \\
\quad ME05a & Velocity of a ship relative to the shore & velocity relative to a frame & refers & as term \\
\midrule
\textbf{ME06} & \textbf{absolute space}. Space in its own nature, without relation to anything external, always similar and immovable. & --- & \textbf{T} token-dependent; $w$\,=\,3; $s$\,=\,1/3 & by token \\
\quad ME06a & The immovable space of the Scholium to the Definitions & \textit{no referent} & fails & U: I13 \\
\quad ME06b & That relative to which the water in the rotating bucket rotates, as shown by its concave surface & the local inertial (affine) structure of spacetime & refers & H: I14 \\
\quad ME06c & The rest frame of the centre of the world (Principia, Book III, Hypothesis I) & \textit{no referent} & fails & U: I35 \\
\midrule
\textbf{ME07} & \textbf{absolute time}. Time that flows equably without relation to anything external. & --- & \textbf{T} token-dependent; $w$\,=\,2; $s$\,=\,1/2 & by token \\
\quad ME07a & True and mathematical time of the Scholium & \textit{no referent} & fails & U: I15 \\
\quad ME07b & The time variable in the equations of planetary motion & proper time along the observer's worldline (to the accuracy of the Newtonian limit) & refers & R: I16 \\
\midrule
\textbf{ME08} & \textbf{simultaneity}. A relation between events holding or failing absolutely. & simultaneity relative to an inertial frame & \textbf{P} partial; $w$\,=\,2; $s$\,=\,1 & R: I09 I16 \\
\quad ME08a & Two distant lightning strikes judged simultaneous & simultaneity in the observer's rest frame & refers & as term \\
\midrule
\textbf{ME09} & \textbf{inertial frame}. A frame at rest or in uniform motion relative to absolute space. & global inertial frame of special relativity $\mid$ local freely falling frame of general relativity (pd) & \textbf{P} partial; $w$\,=\,4; $s$\,=\,1 & H: I12 I34; R: I11 I30 \\
\quad ME09a & A frame at rest relative to the fixed stars, gravitation ignored & global inertial frame & refers & as term \\
\quad ME09b & A system of bodies falling freely in a uniform field (Principia, Corollary VI) & local freely falling frame & refers & as term \\
\midrule
\textbf{ME10} & \textbf{Galilean transformation}. The transformation between inertial frames: $x' = x - vt$, $t' = t$. & Lorentz transformation (of which it is the limit $v/c \to 0$) & \textbf{P} partial; $w$\,=\,2; $s$\,=\,1 & R: I10 I11 \\
\quad ME10a & Transformation from shore to ship frame & Lorentz transformation & refers & as term \\
\midrule
\textbf{ME11} & \textbf{gravitation}. A universal attractive force between masses, varying inversely as the square of the distance. & spacetime curvature produced by energy and momentum & \textbf{P} partial; $w$\,=\,5; $s$\,=\,1 & H: I22 I31; R: I21; U: I07 I23 \\
\quad ME11a & The gravitation that keeps the Moon in its orbit & curvature of spacetime near the Earth & refers & as term \\
\midrule
\textbf{ME12} & \textbf{gravitational potential}. The scalar whose gradient gives the gravitational acceleration. & the deviation of $g_{00}$ from its flat value (weak field) & \textbf{P} partial; $w$\,=\,1; $s$\,=\,1 & R: I08 \\
\quad ME12a & Potential of the Sun at the orbit of a planet & $(g_{00}+1)c^2/2$ in the weak field & refers & as term \\
\midrule
\textbf{ME13} & \textbf{orbit}. The path of a body under gravitation; a fixed conic for two bodies. & spatial path of a timelike geodesic & \textbf{P} partial; $w$\,=\,1; $s$\,=\,1 & R: I21 \\
\quad ME13a & The orbit of Mercury & the path of Mercury & refers & as term \\
\midrule
\textbf{ME14} & \textbf{length}. The extension of a body, independent of its motion. & length in a frame (proper length in the rest frame) & \textbf{P} partial; $w$\,=\,2; $s$\,=\,1 & H: I25; R: I17 \\
\quad ME14a & Length of a measuring rod & length of the rod in a frame & refers & as term \\
\midrule
\textbf{ME15} & \textbf{speed of light}. The speed of light relative to the ether, about 300\,000 km/s. & the invariant speed $c$ & \textbf{P} partial; $w$\,=\,4; $s$\,=\,1 & H: I20; X: I19; U: I18 I33 \\
\quad ME15a & R\o mer's and Fizeau's measured speed & the invariant speed $c$ & refers & as term \\
\midrule
\textbf{ME16} & \textbf{ether}. The medium of electromagnetic waves, at rest in absolute space. & --- & \textbf{T} token-dependent; $w$\,=\,3; $s$\,=\,1/3 & by token \\
\quad ME16a & A mechanical medium with elastic properties & \textit{no referent} & fails & U: I18 \\
\quad ME16b & Lorentz: the seat of the electric and magnetic fields & the electromagnetic field & refers & H: I20 \\
\quad ME16c & The rest frame relative to which the Earth's motion was sought by Michelson and Morley (1887) & \textit{no referent} & fails & U: I33 \\
\midrule
\textbf{ME17} & \textbf{gravitational constant}. The universal constant $G$ fixing the strength of gravitation. & the constant $G$ of the field equations & \textbf{E} exact; $w$\,=\,1; $s$\,=\,1 & H: I31 \\
\quad ME17a & Cavendish's torsion-balance value (1798) & the same constant & refers & as term \\
\midrule
\textbf{ME18} & \textbf{centre of mass}. The mass-weighted mean position of a system. & centre of energy in a frame & \textbf{E} exact; $w$\,=\,1; $s$\,=\,1 & H: I29 \\
\quad ME18a & Centre of mass of the solar system & centre of energy of the solar system & refers & as term \\
\midrule
\textbf{ME19} & \textbf{weight}. The gravitational force on a body. & --- & \textbf{T} token-dependent; $w$\,=\,2; $s$\,=\,1/2 & by token \\
\quad ME19a & The weight of a stone in the hand & \textit{no referent} & fails & U: I28 \\
\quad ME19b & What a spring balance registers & the supporting force preventing geodesic motion & refers & H: I32 \\
\midrule
\textbf{ME20} & \textbf{energy}. The conserved quantity comprising kinetic and potential energy. & energy (time component of four-momentum) & \textbf{E} exact; $w$\,=\,1; $s$\,=\,1 & H: I27 \\
\quad ME20a & Energy of a pendulum & energy & refers & as term \\
\midrule
\textbf{ME21} & \textbf{acceleration}. Rate of change of velocity; the same in all inertial frames. & coordinate acceleration $\mid$ proper acceleration (pd) & \textbf{P} partial; $w$\,=\,3; $s$\,=\,1 & H: I34; R: I08 I30 \\
\quad ME21a & Acceleration of a falling body & coordinate acceleration in the Earth frame & refers & as term \\
\quad ME21b & Acceleration felt by a passenger & proper acceleration & refers & as term \\
\midrule
\textbf{ME22} & \textbf{rigid body}. A body whose parts keep their distances; transmits impulses instantaneously. & --- & \textbf{T} token-dependent; $w$\,=\,2; $s$\,=\,1/2 & by token \\
\quad ME22a & Perfect rigidity as an idealization of transmission of impulse & \textit{no referent} & fails & U: I24 \\
\quad ME22b & A measuring rod moved slowly & a body in Born-rigid motion & refers & H: I25 \\
\midrule
\textbf{ME23} & \textbf{action at a distance}. Instantaneous action of one body on another without intervening medium. & --- & \textbf{N} none; $w$\,=\,1; $s$\,=\,0 & by token \\
\quad ME23a & Gravitational attraction across the solar system & \textit{no referent} & fails & U: I23 \\
\bottomrule
\end{manualtable}

\Needspace{10\baselineskip}\subsection*{HEAT: Thermodynamic revolution}
\addcontentsline{toc}{subsection}{HEAT: Thermodynamic revolution}
\noindent\textit{Predecessor} $T_1$: The caloric theory (Black's calorimetry as read by Lavoisier and Laplace, Laplace on sound, Carnot 1824), c.~1780--1840.\\
\textit{Successor} $T_2$: Thermodynamics and the kinetic theory of gases (Joule, Clausius, Kelvin, Maxwell, Boltzmann), 1840--1880.\par\smallskip
\begin{manualtable}
\textbf{HE01} & \textbf{caloric}. A subtle, elastic, weightless fluid whose particles repel one another and are attracted by ordinary matter; conserved. & --- & \textbf{T} token-dependent; $w$\,=\,8; $s$\,=\,1/2 & by token \\
\quad HE01a & The fluid of heat as described by Lavoisier and Laplace & \textit{no referent} & fails & U: I01 I08 I15 I18 I21 \\
\quad HE01b & The quantity of caloric given up by a body in the ice calorimeter & heat transferred (energy) & refers & H: I27 \\
\quad HE01c & Caloric as conserved in all processes & \textit{no referent} & fails & U: I02 \\
\quad HE01d & Carnot: the caloric that falls from the hot to the cold body in a reversible engine & entropy transferred (Callendar's reading); heat (Kelvin's) & refers & R: I14 \\
\midrule
\textbf{HE02} & \textbf{heat (quantity of)}. Quantity of caloric transferred. & heat (energy transferred by thermal interaction) & \textbf{P} partial; $w$\,=\,7; $s$\,=\,1 & H: I03 I04 I07 I23 I24; X: I13 I28 \\
\quad HE02a & Heat given by hot water to cold in Black's method of mixtures & heat transferred & refers & as term \\
\quad HE02b & Heat taken from the furnace by a steam engine & heat transferred & refers & as term \\
\midrule
\textbf{HE03} & \textbf{temperature}. Degree of heat; the tension or intensity of free caloric. & temperature & \textbf{P} partial; $w$\,=\,11; $s$\,=\,1 & H: I03 I04 I05 I07 I09 I12 I16 I19 I30; R: I14 I20 \\
\quad HE03a & Reading of a mercury thermometer & empirical temperature & refers & as term \\
\quad HE03b & Temperature of the boiler and condenser in Carnot & temperature (thermodynamic) & refers & as term \\
\midrule
\textbf{HE04} & \textbf{latent heat}. Caloric absorbed by a body in melting or vaporization, combined with it so as not to affect the thermometer. & latent heat (energy of phase change) & \textbf{P} partial; $w$\,=\,4; $s$\,=\,1 & H: I05 I24 I25; U: I06 \\
\quad HE04a & Black's latent heat of fusion of ice & latent heat of fusion & refers & as term \\
\quad HE04b & Latent heat of steam & latent heat of vaporization & refers & as term \\
\midrule
\textbf{HE05} & \textbf{specific heat}. The capacity of a body for caloric, per unit weight. & specific heat capacity & \textbf{P} partial; $w$\,=\,4; $s$\,=\,1 & H: I07 I11; U: I08 I15 \\
\quad HE05a & Specific heats measured by mixtures (Black, Crawford) & specific heat capacity & refers & as term \\
\midrule
\textbf{HE06} & \textbf{free caloric}. Caloric not combined with matter, sensible to the thermometer. & --- & \textbf{N} none; $w$\,=\,1; $s$\,=\,0 & by token \\
\quad HE06a & The free caloric whose density the thermometer measures & \textit{no referent} & fails & U: I26 \\
\midrule
\textbf{HE07} & \textbf{combined caloric}. Caloric held by chemical affinity in a body, insensible to the thermometer. & --- & \textbf{N} none; $w$\,=\,2; $s$\,=\,0 & by token \\
\quad HE07a & Caloric combined in water to form steam & \textit{no referent} & fails & U: I06 I29 \\
\midrule
\textbf{HE08} & \textbf{conduction}. Passage of heat through a body from hotter to colder parts. & heat conduction & \textbf{E} exact; $w$\,=\,1; $s$\,=\,1 & H: I16 \\
\quad HE08a & Fourier's heat flow in a solid (1822) & heat conduction & refers & as term \\
\midrule
\textbf{HE09} & \textbf{radiant heat}. Caloric emitted in free flight from bodies. & thermal radiation & \textbf{P} partial; $w$\,=\,2; $s$\,=\,1 & H: I17; U: I18 \\
\quad HE09a & Prevost's exchanges of radiant heat (1791) & thermal radiation & refers & as term \\
\quad HE09b & Herschel's rays beyond the red (1800) & infrared radiation & refers & as term \\
\midrule
\textbf{HE10} & \textbf{adiabatic compression}. Compression without exchange of caloric with the surroundings, which squeezes caloric out of combination and raises temperature. & adiabatic compression & \textbf{E} exact; $w$\,=\,3; $s$\,=\,1 & H: I09 I10 I27 \\
\quad HE10a & Compression in a sound wave (Laplace) & adiabatic compression & refers & as term \\
\midrule
\textbf{HE11} & \textbf{ratio of specific heats}. The ratio of the specific heat of a gas at constant pressure to that at constant volume. & ratio of specific heats $\gamma$ & \textbf{E} exact; $w$\,=\,2; $s$\,=\,1 & H: I10 I11 \\
\quad HE11a & The ratio inferred by Laplace from the speed of sound & $\gamma$ & refers & as term \\
\midrule
\textbf{HE12} & \textbf{speed of sound}. Speed of a pressure wave in air. & speed of sound & \textbf{E} exact; $w$\,=\,1; $s$\,=\,1 & H: I10 \\
\quad HE12a & The measured speed of sound in air & speed of sound & refers & as term \\
\midrule
\textbf{HE13} & \textbf{heat engine}. A device that produces motive power from the passage of caloric from a hot to a cold body. & heat engine & \textbf{P} partial; $w$\,=\,3; $s$\,=\,1 & H: I12; R: I14; X: I13 \\
\quad HE13a & The steam engine & heat engine & refers & as term \\
\midrule
\textbf{HE14} & \textbf{reversible process}. A process that can be run backward through the same states (Carnot). & reversible process & \textbf{E} exact; $w$\,=\,1; $s$\,=\,1 & H: I12 \\
\quad HE14a & Carnot's ideal cycle & Carnot cycle & refers & as term \\
\midrule
\textbf{HE15} & \textbf{frictional heating}. Release of caloric from bodies whose capacity is reduced by friction or abrasion. & conversion of work into heat by friction & \textbf{P} partial; $w$\,=\,2; $s$\,=\,1 & X: I28; U: I15 \\
\quad HE15a & Rumford's boring of cannon (1798) & conversion of work into heat & refers & as term \\
\midrule
\textbf{HE16} & \textbf{thermal equilibrium}. State in which bodies in contact no longer exchange caloric. & thermal equilibrium & \textbf{E} exact; $w$\,=\,2; $s$\,=\,1 & H: I17 I19 \\
\quad HE16a & Bodies left in contact & thermal equilibrium & refers & as term \\
\midrule
\textbf{HE17} & \textbf{thermal expansion}. Increase of volume due to the repulsion of added caloric. & thermal expansion & \textbf{P} partial; $w$\,=\,1; $s$\,=\,1 & R: I20 \\
\quad HE17a & Expansion of gases measured by Gay-Lussac (1802) & thermal expansion & refers & as term \\
\midrule
\textbf{HE18} & \textbf{elastic fluid (gas)}. A fluid whose particles are held apart by atmospheres of caloric. & gas & \textbf{P} partial; $w$\,=\,5; $s$\,=\,1 & H: I09 I11; R: I20 I22; U: I21 \\
\quad HE18a & Air and the permanent gases & gases & refers & as term \\
\midrule
\textbf{HE19} & \textbf{elasticity (of a gas)}. The spring of a gas, due to the repulsion of caloric atmospheres. & pressure & \textbf{P} partial; $w$\,=\,2; $s$\,=\,1 & R: I22; U: I21 \\
\quad HE19a & The spring of air in Boyle's experiments & pressure & refers & as term \\
\midrule
\textbf{HE20} & \textbf{vapour}. A liquid rendered elastic by combination with caloric. & vapour (the gaseous phase) & \textbf{P} partial; $w$\,=\,2; $s$\,=\,1 & H: I25; U: I29 \\
\quad HE20a & Steam & water vapour & refers & as term \\
\midrule
\textbf{HE21} & \textbf{calorimeter}. An instrument measuring caloric released by a body. & calorimeter & \textbf{E} exact; $w$\,=\,1; $s$\,=\,1 & H: I24 \\
\quad HE21a & Lavoisier and Laplace's ice calorimeter (1783) & ice calorimeter & refers & as term \\
\midrule
\textbf{HE22} & \textbf{thermometer}. An instrument measuring the degree of heat. & thermometer & \textbf{P} partial; $w$\,=\,2; $s$\,=\,1 & H: I30; U: I26 \\
\quad HE22a & Mercury thermometer & thermometer & refers & as term \\
\midrule
\textbf{HE23} & \textbf{animal heat}. Caloric released in respiration. & body heat (produced by metabolism) & \textbf{E} exact; $w$\,=\,1; $s$\,=\,1 & H: I23 \\
\quad HE23a & Heat given off by a guinea pig in the ice calorimeter (Lavoisier and Laplace) & metabolic heat & refers & as term \\
\bottomrule
\end{manualtable}

\Needspace{10\baselineskip}\subsection*{QUAN: Quantum revolution}
\addcontentsline{toc}{subsection}{QUAN: Quantum revolution}
\noindent\textit{Predecessor} $T_1$: Classical mechanics and electrodynamics applied to atoms and radiation, with the nuclear atom (Rutherford, 1911), c.~1895--1912.\\
\textit{Successor} $T_2$: Quantum mechanics (Heisenberg 1925; Born, Jordan; Schr\"odinger 1926; Dirac) in its standard formulation.\par\smallskip
\begin{manualtable}
\textbf{QU01} & \textbf{trajectory}. The continuous path of a particle, fixed by its position and momentum at one time and the forces. & --- & \textbf{T} token-dependent; $w$\,=\,4; $s$\,=\,2/3 & by token \\
\quad QU01a & Trajectory of an electron inside an atom & \textit{no referent} & fails & U: I30 \\
\quad QU01b & Track of an alpha particle in a cloud chamber (C.T.R. Wilson, 1911--1912) & the path of a well-localized wave packet (Mott, 1929) & refers & H: I09 \\
\quad QU01c & Trajectory of a projectile or planet & the path of the centre of a localized state (Ehrenfest's theorem) & refers & R: I08 I29 \\
\midrule
\textbf{QU02} & \textbf{electron orbit}. The closed path of an electron revolving about the nucleus. & --- & \textbf{T} token-dependent; $w$\,=\,4; $s$\,=\,1/2 & by token \\
\quad QU02a & Orbit of the electron in the ground state of hydrogen & \textit{no referent} & fails & U: I01 I02 I04 \\
\quad QU02b & Orbit of an electron in a very highly excited atom & a Rydberg wave packet circulating with the classical orbital period & refers & R: I05 \\
\midrule
\textbf{QU03} & \textbf{position}. A quantity possessed with a definite value by every particle at every time. & position (an observable) & \textbf{P} partial; $w$\,=\,3; $s$\,=\,1 & R: I08; X: I06 I07 \\
\quad QU03a & Position of an electron on a scintillation screen & position observable (measured value) & refers & as term \\
\midrule
\textbf{QU04} & \textbf{momentum}. Mass times velocity, possessed with a definite value at every time. & momentum (an observable) & \textbf{P} partial; $w$\,=\,3; $s$\,=\,1 & R: I08; X: I06 I07 \\
\quad QU04a & Momentum of a cathode-ray electron deflected in a magnetic field & momentum observable & refers & as term \\
\midrule
\textbf{QU05} & \textbf{energy}. A quantity taking any value in a continuous range for a bound system. & energy (the Hamiltonian, with discrete eigenvalues for bound states) & \textbf{P} partial; $w$\,=\,3; $s$\,=\,1 & R: I11; X: I10 I16 \\
\quad QU05a & Energy of the electron in an atom & energy of the bound state & refers & as term \\
\midrule
\textbf{QU06} & \textbf{electron}. J.J. Thomson's corpuscle (1897): a small charged particle with definite charge and mass. & electron & \textbf{P} partial; $w$\,=\,4; $s$\,=\,1 & H: I09 I17; X: I10; U: I01 \\
\quad QU06a & Cathode-ray corpuscles & electrons & refers & as term \\
\quad QU06b & Electrons in Millikan's oil drops & electrons & refers & as term \\
\midrule
\textbf{QU07} & \textbf{atom}. A nucleus with electrons moving about it under electric forces (Rutherford). & atom & \textbf{P} partial; $w$\,=\,4; $s$\,=\,1 & H: I19 I28; X: I10; U: I01 \\
\quad QU07a & The hydrogen atom & hydrogen atom & refers & as term \\
\midrule
\textbf{QU08} & \textbf{nucleus}. A small, massive, positively charged centre of the atom. & nucleus & \textbf{P} partial; $w$\,=\,5; $s$\,=\,1 & H: I19 I20; U: I01 I02 I30 \\
\quad QU08a & Rutherford's nucleus (1911) & nucleus & refers & as term \\
\midrule
\textbf{QU09} & \textbf{light wave}. An electromagnetic wave in which energy is continuously distributed over the wave front. & electromagnetic radiation (the quantized field) & \textbf{P} partial; $w$\,=\,5; $s$\,=\,1 & H: I15 I21 I27; X: I14 I16 \\
\quad QU09a & Light in an interference experiment & electromagnetic radiation & refers & as term \\
\midrule
\textbf{QU10} & \textbf{radiation of an accelerated charge}. Continuous emission of electromagnetic energy by every accelerated charge (Larmor, 1897). & electromagnetic emission by charges & \textbf{P} partial; $w$\,=\,3; $s$\,=\,1 & H: I27; R: I03; U: I02 \\
\quad QU10a & Radiation from currents in an antenna & emission from macroscopic currents & refers & as term \\
\quad QU10b & Radiation expected from an electron revolving in an atom & transition radiation & refers & as term \\
\midrule
\textbf{QU11} & \textbf{spectral line}. Radiation emitted at the frequency of a periodic motion of the electrons or its harmonics. & transition frequency between stationary states & \textbf{P} partial; $w$\,=\,3; $s$\,=\,1 & H: I28; R: I05; U: I04 \\
\quad QU11a & The Balmer lines of hydrogen & transition frequencies of hydrogen & refers & as term \\
\midrule
\textbf{QU12} & \textbf{equipartition}. In thermal equilibrium each degree of freedom has mean energy $kT/2$. & the high-temperature limit of the quantum mean energy & \textbf{P} partial; $w$\,=\,3; $s$\,=\,1 & R: I11 I12 I13 \\
\quad QU12a & Equipartition among vibrational modes of a solid & high-temperature limit & refers & as term \\
\midrule
\textbf{QU13} & \textbf{black-body radiation}. Radiation in thermal equilibrium within an enclosure. & black-body radiation & \textbf{P} partial; $w$\,=\,1; $s$\,=\,1 & R: I12 \\
\quad QU13a & Cavity radiation measured by Lummer and Pringsheim, Rubens and Kurlbaum & black-body radiation & refers & as term \\
\midrule
\textbf{QU14} & \textbf{atomic heat}. The product of specific heat and atomic weight of a solid element, constant by the law of Dulong and Petit. & molar heat capacity of a solid & \textbf{P} partial; $w$\,=\,1; $s$\,=\,1 & R: I13 \\
\quad QU14a & Atomic heat of copper at room temperature & molar heat capacity & refers & as term \\
\midrule
\textbf{QU15} & \textbf{photoelectric emission}. Release of electrons from metals by light (Hertz, 1887; Lenard, 1902). & photoelectric effect & \textbf{P} partial; $w$\,=\,1; $s$\,=\,1 & X: I14 \\
\quad QU15a & Lenard's measurements of emitted electron energies & photoelectric effect & refers & as term \\
\midrule
\textbf{QU16} & \textbf{angular momentum}. Moment of momentum, taking any value. & angular momentum (quantized observable; includes spin) & \textbf{P} partial; $w$\,=\,3; $s$\,=\,1 & H: I23; X: I22 I24 \\
\quad QU16a & Angular momentum of an orbiting electron & orbital angular momentum & refers & as term \\
\midrule
\textbf{QU17} & \textbf{magnetic moment}. The moment of a current loop, taking any direction relative to a field. & magnetic moment (quantized components) & \textbf{P} partial; $w$\,=\,1; $s$\,=\,1 & X: I24 \\
\quad QU17a & Magnetic moment of a silver atom (Stern and Gerlach, 1922) & magnetic moment & refers & as term \\
\midrule
\textbf{QU18} & \textbf{state (mechanical)}. A point in phase space: positions and momenta of all particles. & --- & \textbf{T} token-dependent; $w$\,=\,2; $s$\,=\,1/2 & by token \\
\quad QU18a & The microstate of a system of electrons & \textit{no referent} & fails & U: I25 \\
\quad QU18b & The macroscopic state of a gas given by pressure, volume and temperature & the thermodynamic macrostate & refers & H: I26 \\
\midrule
\textbf{QU19} & \textbf{charge}. Electric charge, conserved and coming in multiples of $e$. & charge & \textbf{P} partial; $w$\,=\,3; $s$\,=\,1 & H: I17 I18; R: I03 \\
\quad QU19a & Millikan's elementary charge (1909--1913) & elementary charge & refers & as term \\
\midrule
\textbf{QU20} & \textbf{X-rays}. Short-wavelength electromagnetic radiation. & X-rays & \textbf{E} exact; $w$\,=\,1; $s$\,=\,1 & H: I21 \\
\quad QU20a & X-rays diffracted by crystals (Laue, 1912) & X-rays & refers & as term \\
\midrule
\textbf{QU21} & \textbf{Rutherford scattering}. Deflection of alpha particles by the Coulomb field of a nucleus. & Coulomb scattering & \textbf{E} exact; $w$\,=\,1; $s$\,=\,1 & H: I20 \\
\quad QU21a & Geiger and Marsden's scattering (1909--1913) & Coulomb scattering & refers & as term \\
\bottomrule
\end{manualtable}

\Needspace{10\baselineskip}\subsection*{BIOL: Darwinian revolution}
\addcontentsline{toc}{subsection}{BIOL: Darwinian revolution}
\noindent\textit{Predecessor} $T_1$: Natural history of fixed species (Linnaeus, Cuvier, Paley, Owen), c.~1750--1859.\\
\textit{Successor} $T_2$: Darwin's theory of descent with modification by natural selection (Origin of Species, 1859) and its continuation.\par\smallskip
\begin{manualtable}
\textbf{BI01} & \textbf{species}. A kind of organism created as such, fixed, with limits of variation set by its type. & species (a lineage of interbreeding populations) & \textbf{P} partial; $w$\,=\,11; $s$\,=\,1 & H: I16 I20 I22 I23 I30 I31; X: I01 I02 I03 I26; U: I17 \\
\quad BI01a & Felis leo in Linnaeus' Systema Naturae & the lion & refers & as term \\
\quad BI01b & The Galapagos finches as named by Gould (1837) & the finch species & refers & as term \\
\midrule
\textbf{BI02} & \textbf{variety}. A form within a species that departs from the type and does not persist. & variety (an incipient species) & \textbf{P} partial; $w$\,=\,1; $s$\,=\,1 & X: I02 \\
\quad BI02a & Varieties of cultivated plants & varieties & refers & as term \\
\midrule
\textbf{BI03} & \textbf{genus}. A group of species sharing a common character, an order of the natural system. & genus (a group of species descended from a common ancestor) & \textbf{E} exact; $w$\,=\,2; $s$\,=\,1 & H: I04 I31 \\
\quad BI03a & The genus Felis & the genus Felis & refers & as term \\
\midrule
\textbf{BI04} & \textbf{natural system}. The arrangement of organisms in groups subordinate to groups, expressing their true affinities. & genealogical classification & \textbf{P} partial; $w$\,=\,3; $s$\,=\,1 & H: I04 I06; U: I05 \\
\quad BI04a & The natural system sought by de Candolle and Cuvier & genealogical classification & refers & as term \\
\midrule
\textbf{BI05} & \textbf{plan of creation}. The design in the Creator's mind of which the natural system is the expression. & --- & \textbf{T} token-dependent; $w$\,=\,3; $s$\,=\,1/2 & by token \\
\quad BI05a & The plan as the cause of the natural system & \textit{no referent} & fails & U: I05 I11 \\
\quad BI05b & Unity of plan: the shared structure of the members of a great class & nested pattern of shared ancestral structure & refers & H: I06 \\
\midrule
\textbf{BI06} & \textbf{archetype}. Owen (1848): the ideal pattern of which all vertebrate skeletons are modifications. & --- & \textbf{T} token-dependent; $w$\,=\,2; $s$\,=\,1/2 & by token \\
\quad BI06a & The archetype as an idea in the divine mind & \textit{no referent} & fails & U: I07 \\
\quad BI06b & The generalized vertebrate pattern drawn by Owen & the structure of the common ancestor of the vertebrates (Darwin, Origin, ch.~13) & refers & H: I08 \\
\midrule
\textbf{BI07} & \textbf{homology}. The same organ in different animals under every variety of form and function (Owen). & homology (similarity due to common descent) & \textbf{E} exact; $w$\,=\,3; $s$\,=\,1 & H: I08 I09 I10 \\
\quad BI07a & The forelimb bones of man, bat and porpoise & homologous forelimbs & refers & as term \\
\midrule
\textbf{BI08} & \textbf{analogy}. Similarity of function in organs of different structure. & analogy (convergent similarity) & \textbf{E} exact; $w$\,=\,1; $s$\,=\,1 & H: I10 \\
\quad BI08a & Wings of birds and of insects & analogous wings & refers & as term \\
\midrule
\textbf{BI09} & \textbf{adaptation (contrivance)}. The fitting of an organism's parts to its conditions of life, evidence of design (Paley, 1802). & adaptation (product of natural selection) & \textbf{P} partial; $w$\,=\,2; $s$\,=\,1 & H: I15; X: I14 \\
\quad BI09a & Paley's eye & the eye as an adaptation & refers & as term \\
\midrule
\textbf{BI10} & \textbf{centre of creation}. The place where a species was created. & --- & \textbf{T} token-dependent; $w$\,=\,2; $s$\,=\,1/2 & by token \\
\quad BI10a & The place and act of separate creation of a species & \textit{no referent} & fails & U: I17 \\
\quad BI10b & The single region from which a species spread & the area of origin of the species (Darwin's single centres) & refers & H: I16 \\
\midrule
\textbf{BI11} & \textbf{rudimentary organ}. An organ present in reduced form, retained for the completion of the plan or for symmetry. & vestigial organ & \textbf{P} partial; $w$\,=\,2; $s$\,=\,1 & H: I12; U: I11 \\
\quad BI11a & Teeth of foetal whales; mammae of male mammals & vestigial organs & refers & as term \\
\midrule
\textbf{BI12} & \textbf{embryonic resemblance}. The greater resemblance of embryos than of adults within a class (von Baer). & embryonic resemblance & \textbf{E} exact; $w$\,=\,1; $s$\,=\,1 & H: I13 \\
\quad BI12a & Von Baer's unlabelled embryos & embryonic resemblance & refers & as term \\
\midrule
\textbf{BI13} & \textbf{extinction}. The disappearance of species, in Cuvier by revolutions of the globe. & extinction & \textbf{P} partial; $w$\,=\,2; $s$\,=\,1 & H: I30; X: I18 \\
\quad BI13a & Extinction of the mammoth and the Megatherium & extinction & refers & as term \\
\midrule
\textbf{BI14} & \textbf{succession of types}. The resemblance of extinct to living forms in the same region (Clift, Owen). & succession of types & \textbf{E} exact; $w$\,=\,1; $s$\,=\,1 & H: I20 \\
\quad BI14a & Fossil edentates of South America and living armadillos & succession of types & refers & as term \\
\midrule
\textbf{BI15} & \textbf{fossil sequence}. The order of appearance of fossil forms in successive strata. & fossil sequence & \textbf{E} exact; $w$\,=\,1; $s$\,=\,1 & H: I19 \\
\quad BI15a & The faunal succession of William Smith and Cuvier & fossil sequence & refers & as term \\
\midrule
\textbf{BI16} & \textbf{chain of being}. A single linear scale of organisms from lowest to highest, without gaps. & --- & \textbf{N} none; $w$\,=\,1; $s$\,=\,0 & by token \\
\quad BI16a & The scala naturae of Bonnet & \textit{no referent} & fails & U: I21 \\
\midrule
\textbf{BI17} & \textbf{domestic race}. A breed of domestic animal or cultivated plant. & breed & \textbf{E} exact; $w$\,=\,2; $s$\,=\,1 & H: I22 I28 \\
\quad BI17a & Breeds of pigeons & breeds & refers & as term \\
\midrule
\textbf{BI18} & \textbf{hybrid sterility}. The infertility of crosses between species, a provision against their confusion. & hybrid sterility & \textbf{P} partial; $w$\,=\,2; $s$\,=\,1 & H: I23; X: I24 \\
\quad BI18a & Sterility of the mule & hybrid sterility & refers & as term \\
\midrule
\textbf{BI19} & \textbf{struggle for existence}. The checks on increase described by Malthus and de Candolle's war of nature. & struggle for existence & \textbf{E} exact; $w$\,=\,1; $s$\,=\,1 & H: I25 \\
\quad BI19a & Malthus's checks on population & struggle for existence & refers & as term \\
\midrule
\textbf{BI20} & \textbf{balance of nature}. A designed equilibrium in which each species keeps its station (Linnaeus' economy of nature). & the network of ecological relations & \textbf{P} partial; $w$\,=\,1; $s$\,=\,1 & X: I26 \\
\quad BI20a & Linnaeus' Oeconomia naturae & ecological relations & refers & as term \\
\midrule
\textbf{BI21} & \textbf{instinct}. A fixed endowment of each species, as perfect as its structures. & instinct (heritable behaviour subject to selection) & \textbf{P} partial; $w$\,=\,1; $s$\,=\,1 & X: I27 \\
\quad BI21a & The cell-making instinct of the hive bee & instinct & refers & as term \\
\midrule
\textbf{BI22} & \textbf{sport}. A sudden large departure from the parental form. & sport (a large discontinuous variation) & \textbf{E} exact; $w$\,=\,1; $s$\,=\,1 & H: I28 \\
\quad BI22a & Ancon sheep & a large heritable variation & refers & as term \\
\midrule
\textbf{BI23} & \textbf{variation}. Individual differences, confined within the limits of the type. & variation (heritable individual differences) & \textbf{P} partial; $w$\,=\,2; $s$\,=\,1 & H: I22; X: I03 \\
\quad BI23a & Differences among individuals of a litter & variation & refers & as term \\
\midrule
\textbf{BI24} & \textbf{higher and lower organisms}. Position on the scale of perfection. & degree of differentiation and specialization of parts & \textbf{E} exact; $w$\,=\,1; $s$\,=\,1 & H: I29 \\
\quad BI24a & Mammals higher than fishes & degree of differentiation & refers & as term \\
\bottomrule
\end{manualtable}

\Needspace{10\baselineskip}\subsection*{PHYS: Harveian revolution}
\addcontentsline{toc}{subsection}{PHYS: Harveian revolution}
\noindent\textit{Predecessor} $T_1$: Galenic physiology (Galen, On the Natural Faculties, On the Use of Parts; as taught to 1628).\\
\textit{Successor} $T_2$: Harvey's theory of the circulation (De motu cordis, 1628) and its continuation (Malpighi's capillaries, 1661).\par\smallskip
\begin{manualtable}
\textbf{GA01} & \textbf{venous blood}. Nutritive blood made in the liver from chyle and distributed by the veins to be consumed by the parts. & venous blood & \textbf{P} partial; $w$\,=\,5; $s$\,=\,1 & H: I08 I26 I29; X: I01 I02 \\
\quad GA01a & Blood in the vena cava & venous blood & refers & as term \\
\quad GA01b & Blood let from a vein of the arm & venous blood & refers & as term \\
\midrule
\textbf{GA02} & \textbf{arterial blood}. Thin, bright blood mixed with vital spirit, distributed by the arteries. & arterial blood & \textbf{P} partial; $w$\,=\,4; $s$\,=\,1 & H: I08 I23 I29; U: I07 \\
\quad GA02a & Blood from a cut artery & arterial blood & refers & as term \\
\midrule
\textbf{GA03} & \textbf{vital spirit}. Pneuma elaborated in the left ventricle from inspired air and blood; carried by the arteries to vivify the body. & --- & \textbf{N} none; $w$\,=\,4; $s$\,=\,0 & by token \\
\quad GA03a & The vital spirit of the left ventricle & \textit{no referent} & fails & U: I06 I07 I14 I17 \\
\midrule
\textbf{GA04} & \textbf{natural spirit}. Spirit elaborated in the liver and carried by the veins, governing nutrition and growth. & --- & \textbf{N} none; $w$\,=\,1; $s$\,=\,0 & by token \\
\quad GA04a & The natural spirit of the liver & \textit{no referent} & fails & U: I20 \\
\midrule
\textbf{GA05} & \textbf{animal spirit}. Spirit elaborated in the brain from vital spirit, carried by the nerves, producing sensation and motion. & --- & \textbf{T} token-dependent; $w$\,=\,2; $s$\,=\,1/2 & by token \\
\quad GA05a & Spirit elaborated in the rete mirabile and ventricles of the brain & \textit{no referent} & fails & U: I17 \\
\quad GA05b & What passes along the nerves from the brain to move the muscles & nerve conduction & refers & H: I18 \\
\midrule
\textbf{GA06} & \textbf{liver}. The origin of the veins and the organ that makes the blood. & liver & \textbf{P} partial; $w$\,=\,3; $s$\,=\,1 & X: I01 I22; U: I20 \\
\quad GA06a & The liver dissected & liver & refers & as term \\
\midrule
\textbf{GA07} & \textbf{heart}. The origin of the arteries and seat of the innate heat; attracts blood in diastole. & heart (a muscular pump) & \textbf{P} partial; $w$\,=\,4; $s$\,=\,1 & H: I30; U: I03 I06 I13 \\
\quad GA07a & The heart dissected and observed beating & heart & refers & as term \\
\midrule
\textbf{GA08} & \textbf{septal pores}. Invisible pores in the septum through which blood passes from the right to the left ventricle. & --- & \textbf{N} none; $w$\,=\,1; $s$\,=\,0 & by token \\
\quad GA08a & The pores of the interventricular septum & \textit{no referent} & fails & U: I03 \\
\midrule
\textbf{GA09} & \textbf{arterial vein}. The vessel carrying blood from the right ventricle to the lungs to nourish them. & pulmonary artery & \textbf{P} partial; $w$\,=\,2; $s$\,=\,1 & H: I04; X: I27 \\
\quad GA09a & Vena arteriosa & pulmonary artery & refers & as term \\
\midrule
\textbf{GA10} & \textbf{venous artery}. The vessel carrying air from the lungs to the left ventricle and sooty vapours back. & pulmonary vein & \textbf{P} partial; $w$\,=\,2; $s$\,=\,1 & X: I05; U: I15 \\
\quad GA10a & Arteria venosa & pulmonary veins & refers & as term \\
\midrule
\textbf{GA11} & \textbf{diastole}. The dilation of the heart, its active motion, by which it attracts blood. & diastole (relaxation and filling) & \textbf{P} partial; $w$\,=\,1; $s$\,=\,1 & U: I09 \\
\quad GA11a & Diastole observed in a vivisected heart & diastole & refers & as term \\
\midrule
\textbf{GA12} & \textbf{systole}. The contraction of the heart. & systole (active contraction and ejection) & \textbf{E} exact; $w$\,=\,2; $s$\,=\,1 & H: I10 I30 \\
\quad GA12a & Systole observed in a vivisected heart & systole & refers & as term \\
\midrule
\textbf{GA13} & \textbf{pulse}. The dilation and contraction of the arteries, produced by a pulsific faculty conveyed from the heart through their coats. & arterial pulse & \textbf{P} partial; $w$\,=\,2; $s$\,=\,1 & H: I24; U: I11 \\
\quad GA13a & The pulse felt at the wrist & arterial pulse & refers & as term \\
\midrule
\textbf{GA14} & \textbf{innate heat}. The heat seated in the heart, source of life, moderated by respiration. & --- & \textbf{T} token-dependent; $w$\,=\,2; $s$\,=\,1/2 & by token \\
\quad GA14a & Innate heat as the principle of life seated in the heart & \textit{no referent} & fails & U: I13 \\
\quad GA14b & The warmth of the living body, felt to the touch & body heat (produced by metabolism) & refers & H: I28 \\
\midrule
\textbf{GA15} & \textbf{sooty vapours}. Waste vapours produced by the innate heat in the heart and expelled through the lungs. & --- & \textbf{T} token-dependent; $w$\,=\,2; $s$\,=\,1/2 & by token \\
\quad GA15a & Vapours passing back from the left ventricle through the venous artery & \textit{no referent} & fails & U: I15 \\
\quad GA15b & Waste carried away by expired air & carbon dioxide and water vapour in expired air & refers & H: I16 \\
\midrule
\textbf{GA16} & \textbf{chyle}. The product of digestion carried from the intestines to the liver. & chyle & \textbf{P} partial; $w$\,=\,2; $s$\,=\,1 & X: I01 I22 \\
\quad GA16a & The milky juice of digestion & chyle & refers & as term \\
\midrule
\textbf{GA17} & \textbf{valves of the veins}. Membranes in the veins that retard the flow of blood to the extremities (Fabricius, 1603). & venous valves & \textbf{P} partial; $w$\,=\,1; $s$\,=\,1 & X: I12 \\
\quad GA17a & Fabricius' ostiola & venous valves & refers & as term \\
\midrule
\textbf{GA18} & \textbf{faculties}. The attractive, retentive, alterative and expulsive powers of the parts. & --- & \textbf{N} none; $w$\,=\,3; $s$\,=\,0 & by token \\
\quad GA18a & The attractive faculty of the heart and of the parts & \textit{no referent} & fails & U: I09 I11 I25 \\
\midrule
\textbf{GA19} & \textbf{nutrition}. The assimilation by each part of the blood it attracts and consumes. & nutrition & \textbf{P} partial; $w$\,=\,4; $s$\,=\,1 & H: I26; X: I02 I27; U: I25 \\
\quad GA19a & Nourishment of the flesh & nutrition & refers & as term \\
\midrule
\textbf{GA20} & \textbf{lungs}. Organs that cool the heart and supply it with air. & lungs & \textbf{P} partial; $w$\,=\,4; $s$\,=\,1 & H: I04; X: I05 I27; U: I14 \\
\quad GA20a & The lungs dissected & lungs & refers & as term \\
\midrule
\textbf{GA21} & \textbf{respiration}. The drawing in of air to cool the innate heat and to supply the matter of vital spirit. & respiration & \textbf{P} partial; $w$\,=\,2; $s$\,=\,1 & H: I16; U: I14 \\
\quad GA21a & Breathing & respiration & refers & as term \\
\midrule
\textbf{GA22} & \textbf{anastomoses}. Connections between arteries and veins at their extremities, through which some blood and spirit pass. & capillary connections & \textbf{E} exact; $w$\,=\,1; $s$\,=\,1 & H: I21 \\
\quad GA22a & Galen's anastomoses & capillary beds & refers & as term \\
\midrule
\textbf{GA23} & \textbf{portal vein}. The vein carrying chyle from the intestines to the liver. & hepatic portal vein & \textbf{P} partial; $w$\,=\,1; $s$\,=\,1 & X: I22 \\
\quad GA23a & Vena portae & hepatic portal vein & refers & as term \\
\midrule
\textbf{GA24} & \textbf{vena cava}. The great vein carrying blood from the liver to the body. & vena cava & \textbf{P} partial; $w$\,=\,1; $s$\,=\,1 & X: I02 \\
\quad GA24a & Vena cava & vena cava & refers & as term \\
\midrule
\textbf{GA25} & \textbf{aorta}. The great artery from the left ventricle. & aorta & \textbf{E} exact; $w$\,=\,2; $s$\,=\,1 & H: I10 I23 \\
\quad GA25a & Aorta & aorta & refers & as term \\
\midrule
\textbf{GA26} & \textbf{rete mirabile}. A network of arteries at the base of the brain in which vital spirit is elaborated into animal spirit. & --- & \textbf{T} token-dependent; $w$\,=\,2; $s$\,=\,1/2 & by token \\
\quad GA26a & The rete mirabile of the human brain & \textit{no referent} & fails & U: I17 \\
\quad GA26b & The rete mirabile dissected in the ox & the carotid rete of ruminants & refers & H: I19 \\
\bottomrule
\end{manualtable}

\Needspace{10\baselineskip}\subsection*{OPTI: Optical revolution}
\addcontentsline{toc}{subsection}{OPTI: Optical revolution}
\noindent\textit{Predecessor} $T_1$: Newtonian emission optics (Opticks, 1704) and its development (Laplace, Biot, Malus), 1704--c.~1815.\\
\textit{Successor} $T_2$: The wave theory of light (Young 1801--1803; Fresnel 1815--1823).\par\smallskip
\begin{manualtable}
\textbf{OP01} & \textbf{ray}. The least part of light, a stream of corpuscles moving in a straight line in a uniform medium. & ray (normal to the wave fronts) & \textbf{P} partial; $w$\,=\,5; $s$\,=\,1 & H: I05 I21 I22; R: I02; X: I12 \\
\quad OP01a & A sunbeam admitted through a hole in the shutter & a narrow beam of light & refers & as term \\
\quad OP01b & Rays of different colours separated by a prism & rays of different wavelengths & refers & as term \\
\midrule
\textbf{OP02} & \textbf{light corpuscle}. A small body emitted by luminous bodies and acted on by forces of matter. & --- & \textbf{N} none; $w$\,=\,6; $s$\,=\,0 & by token \\
\quad OP02a & The particles of Query 29 of the Opticks & \textit{no referent} & fails & U: I01 I03 I07 I11 I17 I20 \\
\midrule
\textbf{OP03} & \textbf{fits of easy reflection and transmission}. Alternating dispositions of a ray, recurring at equal intervals, to be reflected or transmitted at a surface. & --- & \textbf{T} token-dependent; $w$\,=\,4; $s$\,=\,1/2 & by token \\
\quad OP03a & Fits as dispositions of the corpuscles themselves & \textit{no referent} & fails & U: I07 \\
\quad OP03b & The interval of the fits measured from Newton's rings, from which Young (1802) computed wavelengths & half the wavelength of the light in the film & refers & H: I08 I09 I25 \\
\midrule
\textbf{OP04} & \textbf{colour (of a ray)}. An original and unchangeable property of a ray, connected with its degree of refrangibility. & wavelength (frequency) of the light & \textbf{E} exact; $w$\,=\,4; $s$\,=\,1 & H: I06 I09 I23 I24 \\
\quad OP04a & Red and violet light of the prismatic spectrum & long and short wavelengths & refers & as term \\
\midrule
\textbf{OP05} & \textbf{refrangibility}. The disposition of a ray to be refracted more or less, fixed for each colour. & dependence of refractive index on wavelength & \textbf{P} partial; $w$\,=\,4; $s$\,=\,1 & H: I05 I06 I22; U: I03 \\
\quad OP05a & Greater refrangibility of violet than red & larger index at shorter wavelengths & refers & as term \\
\midrule
\textbf{OP06} & \textbf{velocity of light in a medium}. Greater in a denser medium, because the corpuscles are attracted by it. & phase velocity of light in the medium & \textbf{P} partial; $w$\,=\,1; $s$\,=\,1 & X: I04 \\
\quad OP06a & Velocity of light in water & speed of light in water & refers & as term \\
\midrule
\textbf{OP07} & \textbf{inflection}. The bending of rays near the edges of bodies by forces of the edge. & diffraction & \textbf{P} partial; $w$\,=\,2; $s$\,=\,1 & H: I10; U: I11 \\
\quad OP07a & Grimaldi's and Newton's fringes at the edge of a shadow & diffraction fringes & refers & as term \\
\midrule
\textbf{OP08} & \textbf{sides of rays}. Original differences between the sides of a ray, which account for the unusual refraction of Iceland crystal. & --- & \textbf{T} token-dependent; $w$\,=\,2; $s$\,=\,1/2 & by token \\
\quad OP08a & Sides as properties of the corpuscles & \textit{no referent} & fails & U: I14 \\
\quad OP08b & The property of a beam emerging from Iceland crystal that fixes how it is refracted by a second crystal & polarization (the orientation of the transverse vibration) & refers & H: I13 \\
\midrule
\textbf{OP09} & \textbf{polarization}. Malus (1808--1810): the arrangement of the axes of the light particles by reflection or double refraction. & polarization & \textbf{P} partial; $w$\,=\,3; $s$\,=\,1 & H: I15 I16; U: I17 \\
\quad OP09a & Light reflected from a window and viewed through Iceland crystal (Malus) & polarized light & refers & as term \\
\midrule
\textbf{OP10} & \textbf{double refraction}. The division of a beam into two by Iceland crystal. & birefringence & \textbf{E} exact; $w$\,=\,2; $s$\,=\,1 & H: I13 I16 \\
\quad OP10a & Iceland crystal & birefringence of calcite & refers & as term \\
\midrule
\textbf{OP11} & \textbf{Newton's rings}. Coloured rings formed in a thin film of air between two glasses. & thin-film interference rings & \textbf{E} exact; $w$\,=\,3; $s$\,=\,1 & H: I08 I24 I25 \\
\quad OP11a & Rings between a lens and a flat glass & interference rings & refers & as term \\
\midrule
\textbf{OP12} & \textbf{aberration of starlight}. The annual displacement of stars due to the Earth's motion and the finite speed of light (Bradley, 1729). & aberration & \textbf{E} exact; $w$\,=\,1; $s$\,=\,1 & H: I18 \\
\quad OP12a & Bradley's observations of gamma Draconis & aberration & refers & as term \\
\midrule
\textbf{OP13} & \textbf{white light}. A heterogeneous mixture of rays of all colours. & light of many wavelengths superposed & \textbf{E} exact; $w$\,=\,2; $s$\,=\,1 & H: I05 I23 \\
\quad OP13a & Sunlight & sunlight & refers & as term \\
\midrule
\textbf{OP14} & \textbf{homogeneal light}. Light of one colour and one refrangibility. & monochromatic light & \textbf{E} exact; $w$\,=\,2; $s$\,=\,1 & H: I06 I22 \\
\quad OP14a & Light from one part of the prismatic spectrum & nearly monochromatic light & refers & as term \\
\midrule
\textbf{OP15} & \textbf{reflection}. The return of rays from a surface by a repulsive power acting at a small distance from it. & reflection & \textbf{P} partial; $w$\,=\,3; $s$\,=\,1 & H: I15 I21; U: I20 \\
\quad OP15a & Reflection from glass & reflection & refers & as term \\
\midrule
\textbf{OP16} & \textbf{geometrical shadow}. The region a straight-line propagation leaves dark behind an obstacle. & geometrical shadow & \textbf{P} partial; $w$\,=\,2; $s$\,=\,1 & H: I10; X: I12 \\
\quad OP16a & The shadow of a small disk & geometrical shadow & refers & as term \\
\midrule
\textbf{OP17} & \textbf{intensity}. The density of rays. & intensity & \textbf{E} exact; $w$\,=\,1; $s$\,=\,1 & H: I19 \\
\quad OP17a & Brightness of a candle at different distances & intensity & refers & as term \\
\midrule
\textbf{OP18} & \textbf{prismatic spectrum}. The oblong image formed by a prism from sunlight. & dispersion spectrum & \textbf{E} exact; $w$\,=\,1; $s$\,=\,1 & H: I23 \\
\quad OP18a & Newton's spectrum in the experimentum crucis & dispersion spectrum & refers & as term \\
\bottomrule
\end{manualtable}

\Needspace{10\baselineskip}\subsection*{MIAS: Bacteriological revolution}
\addcontentsline{toc}{subsection}{MIAS: Bacteriological revolution}
\noindent\textit{Predecessor} $T_1$: Miasmatic and zymotic theories of epidemic disease (Chadwick, Southwood Smith, Farr; Liebig's chemical theory of ferments), c.~1840--1875.\\
\textit{Successor} $T_2$: The germ theory of disease and fermentation (Pasteur, Lister, Koch), 1857--1890, and its continuation.\par\smallskip
\begin{manualtable}
\textbf{MI01} & \textbf{miasma}. A noxious effluvium from decaying organic matter, marshes and filth, which produces epidemic disease when breathed. & --- & \textbf{T} token-dependent; $w$\,=\,5; $s$\,=\,2/3 & by token \\
\quad MI01a & The poisonous emanation of putrefying matter as such & \textit{no referent} & fails & U: I01 I14 I19 \\
\quad MI01b & The cause of intermittent fevers in marshy districts & infection by Plasmodium, carried by mosquitoes breeding in marshes (Laveran, 1880; Ross, 1897) & refers & H: I02 \\
\quad MI01c & The cause of cholera in low-lying districts with foul air and water (Farr) & Vibrio cholerae carried in contaminated water & refers & R: I31 \\
\midrule
\textbf{MI02} & \textbf{zymotic disease}. Farr's class of epidemic, endemic and contagious diseases, produced by a process in the blood analogous to fermentation. & infectious disease & \textbf{P} partial; $w$\,=\,2; $s$\,=\,1 & H: I24; U: I03 \\
\quad MI02a & Smallpox, measles, scarlatina, typhus and cholera in Farr's nosology & these infectious diseases & refers & as term \\
\midrule
\textbf{MI03} & \textbf{ferment}. An agent that induces fermentation or putrefaction in another substance. & --- & \textbf{T} token-dependent; $w$\,=\,4; $s$\,=\,2/3 & by token \\
\quad MI03a & Yeast in alcoholic fermentation (Cagniard-Latour, Schwann, 1837; Pasteur, 1857--1860) & yeast cells & refers & H: I07 \\
\quad MI03b & Diastase, extracted from malt (Payen and Persoz, 1833) & amylase (an enzyme) & refers & H: I29 \\
\quad MI03c & Liebig: a body in decomposition communicating its molecular motion to the fermentable substance & \textit{no referent} & fails & U: I03 I08 \\
\midrule
\textbf{MI04} & \textbf{contagion}. Communication of a disease from the sick to the healthy by contact or by objects they have touched. & transmission of infection & \textbf{E} exact; $w$\,=\,2; $s$\,=\,1 & H: I04 I27 \\
\quad MI04a & Communication of smallpox to attendants and by clothing & transmission of smallpox & refers & as term \\
\midrule
\textbf{MI05} & \textbf{contagium}. The specific matter by which a contagious disease is communicated. & --- & \textbf{T} token-dependent; $w$\,=\,4; $s$\,=\,2/3 & by token \\
\quad MI05a & Variolous matter used in inoculation & variola virus & refers & H: I04 I24 \\
\quad MI05b & Henle's contagium animatum (1840): a living agent & pathogenic microorganisms & refers & H: I05 \\
\quad MI05c & A chemical poison that multiplies by inducing decomposition in the blood & \textit{no referent} & fails & U: I06 \\
\midrule
\textbf{MI06} & \textbf{putrefaction}. Spontaneous chemical decomposition of dead organic matter, initiated by oxygen. & putrefaction (decomposition by microorganisms) & \textbf{P} partial; $w$\,=\,2; $s$\,=\,1 & H: I30; X: I09 \\
\quad MI06a & Putrefaction of meat broth & microbial decomposition & refers & as term \\
\midrule
\textbf{MI07} & \textbf{spontaneous generation}. Origin of microscopic organisms in infusions without parent organisms. & --- & \textbf{T} token-dependent; $w$\,=\,2; $s$\,=\,1/2 & by token \\
\quad MI07a & The doctrine that living beings arise from non-living matter & \textit{no referent} & fails & U: I10 \\
\quad MI07b & The organisms appearing in boiled hay infusions (Pouchet, 1859) & organisms grown from heat-resistant spores (Cohn; Tyndall, 1876--1877) & refers & H: I11 \\
\midrule
\textbf{MI08} & \textbf{epidemic constitution}. A state of the atmosphere favouring a particular epidemic (Sydenham). & --- & \textbf{N} none; $w$\,=\,2; $s$\,=\,0 & by token \\
\quad MI08a & The atmospheric constitution of an epidemic year & \textit{no referent} & fails & U: I12 I21 \\
\midrule
\textbf{MI09} & \textbf{intermittent fever}. A fever recurring at regular intervals, endemic in marshy districts. & malaria & \textbf{E} exact; $w$\,=\,2; $s$\,=\,1 & H: I02 I28 \\
\quad MI09a & Tertian and quartan fevers & malaria (tertian and quartan forms) & refers & as term \\
\midrule
\textbf{MI10} & \textbf{continued fever}. A fever without regular intermission, treated as one disease varying in form. & typhus $\mid$ typhoid fever (pd) & \textbf{P} partial; $w$\,=\,1; $s$\,=\,1 & X: I25 \\
\quad MI10a & The continued fever of the London Fever Hospital before 1850 & typhus or typhoid (partial denotation) & refers & as term \\
\midrule
\textbf{MI11} & \textbf{filth}. Accumulated organic refuse, excrement and decaying matter, the source of miasma. & organic refuse (a source of pathogens and of their vectors) & \textbf{P} partial; $w$\,=\,2; $s$\,=\,1 & H: I13; U: I14 \\
\quad MI11a & Cesspools and refuse of the courts described in Chadwick's Report (1842) & sources of faecal contamination & refers & as term \\
\midrule
\textbf{MI12} & \textbf{sewer gas}. Gas escaping from drains, a cause of typhoid and other fevers. & --- & \textbf{N} none; $w$\,=\,1; $s$\,=\,0 & by token \\
\quad MI12a & Sewer gas as the cause of typhoid & \textit{no referent} & fails & U: I15 \\
\midrule
\textbf{MI13} & \textbf{cholera poison}. The specific poison producing cholera. & --- & \textbf{T} token-dependent; $w$\,=\,2; $s$\,=\,1/2 & by token \\
\quad MI13a & Snow (1849, 1855): morbid matter that multiplies in the gut and passes in the evacuations into water & Vibrio cholerae & refers & H: I16 \\
\quad MI13b & Pettenkofer: a poison that must mature in porous soil before it can cause disease & \textit{no referent} & fails & U: I17 \\
\midrule
\textbf{MI14} & \textbf{disinfectant}. An agent that destroys miasma and removes offensive smells. & disinfectant (an agent killing microorganisms) & \textbf{P} partial; $w$\,=\,2; $s$\,=\,1 & H: I20; U: I19 \\
\quad MI14a & Chloride of lime & hypochlorite disinfectant & refers & as term \\
\midrule
\textbf{MI15} & \textbf{vaccination}. Inoculation with cowpox matter to protect against smallpox (Jenner, 1798). & vaccination & \textbf{E} exact; $w$\,=\,1; $s$\,=\,1 & H: I22 \\
\quad MI15a & Jennerian vaccination & vaccination with cowpox virus & refers & as term \\
\midrule
\textbf{MI16} & \textbf{hospital gangrene}. A spreading destruction of wounds in crowded hospital wards. & wound infection (necrotizing) & \textbf{E} exact; $w$\,=\,1; $s$\,=\,1 & H: I23 \\
\quad MI16a & Hospital gangrene in military hospitals & necrotizing wound infection & refers & as term \\
\midrule
\textbf{MI17} & \textbf{puerperal fever}. Fever following childbirth, often epidemic in lying-in wards. & puerperal sepsis & \textbf{P} partial; $w$\,=\,2; $s$\,=\,1 & H: I20; U: I21 \\
\quad MI17a & Puerperal fever in the Vienna maternity clinic (Semmelweis, 1847) & puerperal sepsis & refers & as term \\
\midrule
\textbf{MI18} & \textbf{infection}. The taking of a disease from a source, by air, contact or objects. & infection (entry and multiplication of a pathogen) & \textbf{E} exact; $w$\,=\,2; $s$\,=\,1 & H: I27 I30 \\
\quad MI18a & Infection by fomites & infection & refers & as term \\
\midrule
\textbf{MI19} & \textbf{predisposing cause}. A condition of the body or of life that renders a person liable to an epidemic. & host susceptibility & \textbf{E} exact; $w$\,=\,1; $s$\,=\,1 & H: I26 \\
\quad MI19a & Poverty, intemperance, debility & factors of susceptibility & refers & as term \\
\midrule
\textbf{MI20} & \textbf{fermentation}. Decomposition of a substance induced by a ferment. & fermentation (microbial or enzymatic metabolism) & \textbf{P} partial; $w$\,=\,2; $s$\,=\,1 & H: I07; U: I08 \\
\quad MI20a & Alcoholic fermentation of grape juice & alcoholic fermentation by yeast & refers & as term \\
\midrule
\textbf{MI21} & \textbf{cholera}. Asiatic cholera, an epidemic disease with violent purging. & cholera (infection by Vibrio cholerae) & \textbf{P} partial; $w$\,=\,3; $s$\,=\,1 & H: I16 I18; R: I31 \\
\quad MI21a & The London epidemics of 1849 and 1854 & cholera & refers & as term \\
\bottomrule
\end{manualtable}

\Needspace{10\baselineskip}\subsection*{HERE: Mendelian revolution}
\addcontentsline{toc}{subsection}{HERE: Mendelian revolution}
\noindent\textit{Predecessor} $T_1$: Pre-Mendelian theories of heredity: blending inheritance, Darwin's pangenesis (1868), Galton's ancestral heredity, the inheritance of acquired characters, c.~1850--1900.\\
\textit{Successor} $T_2$: Mendelian genetics (Mendel 1866, rediscovered 1900) and the chromosome theory (Morgan, 1910--1915).\par\smallskip
\begin{manualtable}
\textbf{HR01} & \textbf{gemmule}. Darwin (1868): minute particles thrown off by every cell, circulating through the body and collecting in the reproductive elements; may lie dormant for generations. & --- & \textbf{T} token-dependent; $w$\,=\,2; $s$\,=\,1/2 & by token \\
\quad HR01a & Gemmules as particles shed by the cells of the body & \textit{no referent} & fails & U: I01 \\
\quad HR01b & Dormant gemmules transmitted unexpressed through generations, to which Darwin ascribed reversion & recessive alleles carried unexpressed & refers & R: I06 \\
\midrule
\textbf{HR02} & \textbf{blending inheritance}. The mixing of the parents' hereditary natures in the offspring, so that their characters are averaged. & intermediate inheritance of quantitative characters & \textbf{P} partial; $w$\,=\,5; $s$\,=\,1 & R: I22; X: I02 I03 I11 I23 \\
\quad HR02a & Children of tall and short parents intermediate in stature & quantitative (polygenic) inheritance & refers & as term \\
\quad HR02b & Jenkin (1867): the dilution of a rare variation by crossing & the frequency of a rare variant in a population & refers & as term \\
\midrule
\textbf{HR03} & \textbf{prepotency}. The stronger power of one parent, or one breed, to transmit its characters to the offspring. & dominance & \textbf{E} exact; $w$\,=\,1; $s$\,=\,1 & H: I04 \\
\quad HR03a & Darwin on prepotent breeds in crosses & dominance of an allele & refers & as term \\
\midrule
\textbf{HR04} & \textbf{reversion}. The reappearance of an ancestral character after one or more generations. & reappearance of a recessive or recombined phenotype & \textbf{P} partial; $w$\,=\,2; $s$\,=\,1 & H: I05; R: I06 \\
\quad HR04a & Crosses of distinct breeds of pigeon yielding the blue colouring of the rock pigeon & recombination of complementary alleles & refers & as term \\
\midrule
\textbf{HR05} & \textbf{use-inheritance}. The transmission to offspring of characters acquired by use or disuse. & --- & \textbf{N} none; $w$\,=\,1; $s$\,=\,0 & by token \\
\quad HR05a & Inheritance of effects of use and disuse & \textit{no referent} & fails & U: I07 \\
\midrule
\textbf{HR06} & \textbf{telegony}. The influence of a previous sire on offspring of a female by later sires. & --- & \textbf{N} none; $w$\,=\,1; $s$\,=\,0 & by token \\
\quad HR06a & Lord Morton's mare and quagga (1820) & \textit{no referent} & fails & U: I08 \\
\midrule
\textbf{HR07} & \textbf{pure breed}. A breed or race that reproduces its characters without variation. & line homozygous for the relevant characters & \textbf{P} partial; $w$\,=\,3; $s$\,=\,1 & H: I05 I09; X: I14 \\
\quad HR07a & Mendel's true-breeding pea varieties & homozygous lines & refers & as term \\
\midrule
\textbf{HR08} & \textbf{hybrid}. The offspring of a cross between distinct breeds or species. & heterozygote (first-generation hybrid) & \textbf{P} partial; $w$\,=\,2; $s$\,=\,1 & H: I10; X: I11 \\
\quad HR08a & Hybrid peas of the first generation & heterozygotes & refers & as term \\
\midrule
\textbf{HR09} & \textbf{sport}. A sudden marked variation appearing in a single individual. & mutation (with visible effect) & \textbf{E} exact; $w$\,=\,1; $s$\,=\,1 & H: I12 \\
\quad HR09a & Sports in cultivated plants & mutations & refers & as term \\
\midrule
\textbf{HR10} & \textbf{individual variation}. The differences among individuals of one breed or species. & heritable (genetic) variation $\mid$ non-heritable variation induced by the environment (pd) & \textbf{P} partial; $w$\,=\,4; $s$\,=\,1 & R: I13; X: I03 I14 I23 \\
\quad HR10a & Differences in seed size within a bean variety & genetic or environmental variation (partial denotation; Johannsen, 1909) & refers & as term \\
\midrule
\textbf{HR11} & \textbf{character}. A feature of an organism transmitted in heredity. & phenotypic trait & \textbf{P} partial; $w$\,=\,6; $s$\,=\,1 & H: I04 I10 I20; R: I22; X: I02; U: I07 \\
\quad HR11a & Seed colour and shape in peas & phenotypic traits & refers & as term \\
\midrule
\textbf{HR12} & \textbf{hereditary tendency}. The disposition, received from the parents, to develop certain characters. & genotype & \textbf{E} exact; $w$\,=\,1; $s$\,=\,1 & H: I15 \\
\quad HR12a & The tendencies of an individual to produce a character in its offspring & genotype & refers & as term \\
\midrule
\textbf{HR13} & \textbf{ancestral heredity}. Galton (1897): each ancestor contributes a fixed fraction of the heritage, one half from the parents, one quarter from the grandparents, and so on. & --- & \textbf{T} token-dependent; $w$\,=\,2; $s$\,=\,1/2 & by token \\
\quad HR13a & The law as a statement of contributions from ancestral germs & \textit{no referent} & fails & U: I17 \\
\quad HR13b & The measured regression of offspring on parents in stature & regression predicted by the additive polygenic model (Fisher, 1918) & refers & H: I18 \\
\midrule
\textbf{HR14} & \textbf{regression}. The tendency of offspring of exceptional parents to be less exceptional (Galton, 1886). & regression to the mean & \textbf{E} exact; $w$\,=\,2; $s$\,=\,1 & H: I16 I18 \\
\quad HR14a & Galton's stature data & regression to the mean & refers & as term \\
\midrule
\textbf{HR15} & \textbf{germ-plasm}. Weismann: the hereditary substance, continuous across generations and not affected by the body. & germ line (chromosomal hereditary material) & \textbf{E} exact; $w$\,=\,1; $s$\,=\,1 & H: I19 \\
\quad HR15a & Weismann's germ-plasm (1892) & germ-line hereditary material & refers & as term \\
\midrule
\textbf{HR16} & \textbf{latent character}. A character carried but not expressed in an individual, and expressed in its descendants. & recessive or sex-limited character & \textbf{E} exact; $w$\,=\,1; $s$\,=\,1 & H: I20 \\
\quad HR16a & Characters latent in one sex and expressed in the other & sex-limited characters & refers & as term \\
\quad HR16b & Characters skipping a generation & recessive characters & refers & as term \\
\midrule
\textbf{HR17} & \textbf{reproductive element}. The ovule or pollen grain, egg or spermatozoon, by which each parent contributes to the offspring. & gamete & \textbf{P} partial; $w$\,=\,3; $s$\,=\,1 & H: I15 I21; U: I01 \\
\quad HR17a & Pollen and ovules of peas & gametes & refers & as term \\
\bottomrule
\end{manualtable}



\section{Inferences, translations and verdicts}\label{s:inferences}

\Needspace{10\baselineskip}\subsection*{CHEM: inferences}
\addcontentsline{toc}{subsection}{CHEM: inferences}
\begin{inftable}
CH.I01 & Combustion releases phlogiston into the air. & Interpretation: a substance present in every combustible leaves it and passes into the air when it burns. No such substance exists; the air instead loses oxygen. & CH08$^{\circ}$, CH01a$^{\ast}$, CH11$^{\circ}$ & \textbf{U} \\
CH.I02 & Combustion in a closed volume of air stops when the air is saturated with phlogiston. & Interpretation: the air can take up only a limited amount of a substance given off by the burning body. In T2 combustion stops because the oxygen of the enclosed air is used up. & CH08$^{\circ}$, CH01d$^{\ast}$, CH03b$^{\ast}$ & \textbf{U} \\
CH.I03 & Combustion requires air. & Combustion requires air (more exactly, the oxygen in it). & CH08$^{\circ}$, CH11$^{\circ}$ & \textbf{H} \\
CH.I04 & A metal is a compound of its calx and phlogiston. & Interpretation: the metal is its oxide plus a further constituent. T2 has no such constituent; the oxide is the metal plus oxygen. & CH07$^{\circ}$, CH06$^{\circ}$, CH01a$^{\ast}$ & \textbf{U} \\
CH.I05 & Calcination of a metal requires air. & Oxidation of a metal by heating requires air. & CH09, CH07$^{\circ}$, CH11$^{\circ}$ & \textbf{H} \\
CH.I06 & The calx weighs more than the metal from which it was made. & The oxide weighs more than the metal from which it was made. & CH06$^{\circ}$, CH07$^{\circ}$, CH09 & \textbf{H} \\
CH.I07 & Heating a calx with charcoal restores the metal. & Heating an oxide with carbon reduces it to the metal. & CH06$^{\circ}$, CH20$^{\circ}$, CH10$^{\circ}$, CH07$^{\circ}$ & \textbf{H} \\
CH.I08 & In reduction, the calx takes phlogiston from the charcoal. & Interpretation: the oxide gains something from the carbon. In T2 the carbon gains oxygen from the oxide. & CH10$^{\circ}$, CH06$^{\circ}$, CH20$^{\circ}$, CH01b$^{\ast}$ & \textbf{U} \\
CH.I09 & Charcoal burns almost without residue because it is nearly pure phlogiston. & Interpretation: charcoal is almost wholly the principle of inflammability. In T2 charcoal is nearly pure carbon, which burns to a gas. & CH20$^{\circ}$, CH01a$^{\ast}$, CH08$^{\circ}$ & \textbf{U} \\
CH.I10 & The red calx of mercury heated alone yields the metal and dephlogisticated air. & Mercuric oxide heated alone yields mercury and oxygen. & CH06$^{\circ}$, CH07$^{\circ}$, CH02a & \textbf{H} \\
CH.I11 & Dephlogisticated air supports combustion and respiration longer than common air. & Oxygen supports combustion and respiration longer than atmospheric air. & CH02a, CH08$^{\circ}$, CH23, CH11$^{\circ}$ & \textbf{H} \\
CH.I12 & Dephlogisticated air is avid of phlogiston because it has been deprived of it. & Interpretation: the new air is common air emptied of a substance it therefore seeks. In T2 it is oxygen, the active part of air, not a residue. & CH02b$^{\ast}$, CH01a$^{\ast}$ & \textbf{U} \\
CH.I13 & Respiration and combustion both render air unfit for further combustion. & Respiration and combustion both render air unfit for further combustion (both consume its oxygen). & CH23, CH08$^{\circ}$, CH11$^{\circ}$ & \textbf{H} \\
CH.I14 & Air in which a candle has burned out, washed with lime water, extinguishes flame and kills animals. & Nitrogen, the residue after combustion and absorption of carbon dioxide, extinguishes flame and kills animals. & CH03a, CH08$^{\circ}$, CH05 & \textbf{H} \\
CH.I15 & Fixed air is released when chalk is heated or treated with acids. & Carbon dioxide is released when calcium carbonate is heated or treated with acids. & CH05 & \textbf{H} \\
CH.I16 & Fixed air is produced by respiration and by burning charcoal. & Carbon dioxide is produced by respiration and by burning carbon. & CH05, CH23, CH20$^{\circ}$, CH08$^{\circ}$ & \textbf{H} \\
CH.I17 & Inflammable air is obtained by dissolving iron or zinc in dilute vitriolic acid. & Hydrogen is obtained by dissolving iron or zinc in dilute sulphuric acid. & CH04a, CH07$^{\circ}$, CH18$^{\circ}$ & \textbf{H} \\
CH.I18 & Inflammable air exploded with dephlogisticated air leaves water, and nearly the whole of both airs is consumed. & Hydrogen exploded with oxygen yields water, nearly the whole of both gases being consumed. & CH04a, CH02c, CH21$^{\circ}$ & \textbf{H} \\
CH.I19 & Water is an element. & Water is a simple substance. [Denied: water is a compound of hydrogen and oxygen.] & CH21$^{\circ}$, CH22$^{\circ}$ & \textbf{X} \\
CH.I20 & Sulphur is composed of vitriolic acid and phlogiston. & Interpretation: sulphur is sulphuric acid combined with the principle of inflammability. In T2 sulphur is simple and sulphuric acid contains sulphur and oxygen. & CH19$^{\circ}$, CH18$^{\circ}$, CH01a$^{\ast}$ & \textbf{U} \\
CH.I21 & The goodness of an air is measured by its diminution with nitrous air over water. & The oxygen content of an air is measured by its contraction with nitric oxide over water. & CH12$^{\circ}$, CH13 & \textbf{H} \\
CH.I22 & The goodness of an air is its capacity to receive phlogiston. & Interpretation: an air is good in proportion as it can absorb what burning bodies give off. In T2 it is good in proportion to its oxygen. & CH12$^{\circ}$, CH01d$^{\ast}$ & \textbf{U} \\
CH.I23 & Dephlogisticated nitrous air supports the flame of a candle. & Nitrous oxide supports the flame of a candle. & CH14, CH08$^{\circ}$ & \textbf{H} \\
CH.I24 & Marine acid air and alkaline air combine to form a white salt (sal ammoniac). & Hydrogen chloride and ammonia combine to form ammonium chloride. & CH15, CH17 & \textbf{H} \\
CH.I25 & Dephlogisticated marine acid bleaches vegetable colours. & Chlorine (oxygenated muriatic acid) bleaches vegetable colours. & CH16 & \textbf{H} \\
CH.I26 & Metals resemble one another because they all contain phlogiston. & Interpretation: the metals share a common constituent that explains their resemblance. T2 recognizes no such constituent and in 1789 had no explanation of the resemblance. & CH07$^{\circ}$, CH01a$^{\ast}$ & \textbf{U} \\
CH.I27 & Air is an element. & Atmospheric air is a simple substance. [Denied: it is a mixture of oxygen and nitrogen.] & CH11$^{\circ}$, CH22$^{\circ}$ & \textbf{X} \\
CH.I28 & Inflammable air is phlogiston, or nearly so. & Hydrogen is hydrogen, or nearly so (true on this reading: the token of phlogiston refers to hydrogen). & CH04a, CH01c & \textbf{H} \\
CH.I29 & Marsh air burns with a flame. & Methane burns with a flame. & CH04b & \textbf{H} \\
CH.I30 & Heavy inflammable air is produced when charcoal is heated with a metallic calx. & Carbon monoxide is produced when carbon is heated with a metal oxide. & CH04c, CH20$^{\circ}$, CH06$^{\circ}$ & \textbf{H} \\
CH.I31 & Sulphur burned in air yields an acid with water. & Sulphur burned in air yields an acid with water. & CH19$^{\circ}$, CH08$^{\circ}$, CH21$^{\circ}$ & \textbf{H} \\
\bottomrule
\end{inftable}

\Needspace{10\baselineskip}\subsection*{ASTR: inferences}
\addcontentsline{toc}{subsection}{ASTR: inferences}
\begin{inftable}
AS.I01 & Each planet moves on an epicycle whose centre moves on a deferent. & The geocentric motion of a planet is the composition of its heliocentric orbit and the Earth's orbit. [Holds to first order: circular approximations to ellipses.] & AS01a, AS06a, AS05a & \textbf{R} \\
AS.I02 & A superior planet retrogrades when it is in opposition to the Sun. & A superior planet appears to retrograde when it is in opposition to the Sun. & AS11$^{\circ}$, AS01a, AS13$^{\circ}$ & \textbf{H} \\
AS.I03 & A superior planet is brightest when it retrogrades. & A superior planet is brightest when it appears to retrograde (it is then nearest the Earth). & AS11$^{\circ}$, AS01a & \textbf{H} \\
AS.I04 & The epicycle centre of Venus and Mercury always lies on the line from the Earth to the mean Sun. & The Sun, centre of the orbits of Venus and Mercury, lies on the line from the Earth to the Sun. [Holds exactly with the true Sun; the mean Sun is a first-order substitute.] & AS06b, AS05b, AS01d, AS13$^{\circ}$, AS12$^{\circ}$ & \textbf{R} \\
AS.I05 & The radius of a superior planet's epicycle stays parallel to the line from the Earth to the mean Sun. & The component of a superior planet's geocentric motion due to the Earth's orbit stays parallel to the Earth--Sun line. [Holds with circular orbits.] & AS06a, AS01a, AS13$^{\circ}$, AS12$^{\circ}$ & \textbf{R} \\
AS.I06 & The Earth is at rest. & The Earth is at rest. [Denied: it rotates daily and revolves annually.] & AS12$^{\circ}$ & \textbf{X} \\
AS.I07 & Heavy bodies move naturally toward the centre of the universe. & Heavy bodies move naturally toward the centre of the Earth. & AS15a & \textbf{H} \\
AS.I08 & The primum mobile communicates the daily motion to the inner orbs. & Interpretation: an outermost body turns daily and drags the inner heavens with it. T2 recognizes no such body; the daily motion is the Earth's rotation. & AS03a$^{\ast}$, AS16$^{\circ}$, AS04a$^{\ast}$ & \textbf{U} \\
AS.I09 & All celestial bodies share a westward diurnal motion with a period of one sidereal day. & All celestial bodies appear to move westward once each sidereal day. & AS16$^{\circ}$, AS02$^{\circ}$, AS01a & \textbf{H} \\
AS.I10 & The diurnal motion is a real rotation of the heavens about the centre of the universe. & Interpretation: the heavens physically rotate about a fixed central point. In T2 nothing rotates about such a point; the Earth rotates. & AS16$^{\circ}$, AS15b$^{\ast}$ & \textbf{U} \\
AS.I11 & The planets are carried by solid orbs that nest without gaps. & Interpretation: solid shells of a celestial matter carry the planets. T2 has no such shells; Kepler sought instead a force from the Sun. & AS04a$^{\ast}$, AS01a & \textbf{U} \\
AS.I12 & The heavens consist of aether, which is ungenerated and unalterable. & Interpretation: the heavens are made of a special unalterable matter. T2 (with Tycho's parallaxes of the nova of 1572 and the comet of 1577) recognizes no such matter. & AS24a$^{\ast}$ & \textbf{U} \\
AS.I13 & The fixed stars keep their mutual positions throughout the year. & The stars keep their mutual positions throughout the year (no annual parallax is detectable). & AS02$^{\circ}$ & \textbf{H} \\
AS.I14 & Precession is a slow eastward motion of the sphere of the fixed stars. & Precession is a slow motion of the stars relative to the equinoxes. [Denied: it is a slow motion of the Earth's axis.] & AS19$^{\circ}$, AS02$^{\circ}$ & \textbf{X} \\
AS.I15 & The epicycle centre moves uniformly as seen from the equant. & The planet moves uniformly in angle about the empty focus of its ellipse. [Holds to first order in the eccentricity.] & AS08a, AS06a & \textbf{R} \\
AS.I16 & The equant lies at twice the distance of the deferent's centre from the Earth (bisected eccentricity). & The empty focus lies at twice the distance of the ellipse's centre from the Sun. [Holds: Kepler confirmed the bisection of the eccentricity.] & AS08a, AS07$^{\circ}$ & \textbf{R} \\
AS.I17 & The period of the second anomaly is the synodic period. & The period of the apparent effect of the Earth's motion is the synodic period. & AS10, AS20 & \textbf{H} \\
AS.I18 & The first anomaly depends on the planet's longitude in the zodiac. & The non-uniformity of the planet's orbital motion depends on its heliocentric longitude. & AS09 & \textbf{H} \\
AS.I19 & The mean longitude of Venus and Mercury advances with a period of one year. & The mean geocentric longitude of Venus and Mercury advances with the Earth's sidereal period of one year. & AS21b, AS01d & \textbf{H} \\
AS.I20 & The order of the planets from the Earth is Moon, Mercury, Venus, Sun, Mars, Jupiter, Saturn. & The order of distances from the Earth is Moon, Mercury, Venus, Sun, Mars, Jupiter, Saturn. [Denied: Venus and Mercury are sometimes farther than the Sun; the order about the Sun is Mercury, Venus, Earth, Mars, Jupiter, Saturn.] & AS22$^{\circ}$, AS01b, AS14$^{\circ}$, AS13$^{\circ}$ & \textbf{X} \\
AS.I21 & Saturn completes its circuit of the zodiac in about 29 1/2 years. & Saturn completes its orbit in about 29 1/2 years. & AS21a, AS01a & \textbf{H} \\
AS.I22 & Venus never departs from the Sun by more than about 47 degrees. & Venus never appears more than about 47 degrees from the Sun. & AS23, AS01d, AS13$^{\circ}$ & \textbf{H} \\
AS.I23 & Comets are exhalations below the Moon. & Comets are below the Moon. [Denied on Tycho's parallax of the comet of 1577.] & AS25$^{\circ}$, AS14$^{\circ}$ & \textbf{X} \\
AS.I24 & The tropical year is shorter than the time of the Sun's return to the same fixed star. & The tropical year is shorter than the sidereal year, the difference being due to precession. & AS18, AS19$^{\circ}$, AS13$^{\circ}$ & \textbf{H} \\
AS.I25 & The Sun moves uniformly on an eccentric circle about the Earth. & The Earth moves uniformly on an eccentric circle about the Sun. [Holds to first order; Kepler: an ellipse with the area law.] & AS13$^{\circ}$, AS07$^{\circ}$, AS12$^{\circ}$ & \textbf{R} \\
AS.I26 & The Moon is the planet nearest the Earth. & The Moon is the celestial body nearest the Earth. & AS14$^{\circ}$, AS01c & \textbf{H} \\
AS.I27 & The Sun moves annually along the ecliptic. & The Sun appears to move annually along the ecliptic. & AS13$^{\circ}$, AS17 & \textbf{H} \\
AS.I28 & The Moon's motion has an inequality depending on its position in its own circuit. & The Moon's motion has an inequality depending on its position in its orbit (the equation of centre). & AS06c, AS14$^{\circ}$ & \textbf{H} \\
AS.I29 & The orbs fix the absolute distances of the planets. & Interpretation: nested solid shells fix how far each planet is. T2 fixes relative distances by the geometry of the orbits instead. & AS04b$^{\ast}$ & \textbf{U} \\
AS.I30 & Retrograde motion is a real motion of the planet. & Retrograde motion is a real motion of the planet. [Denied: it is apparent.] & AS11$^{\circ}$ & \textbf{X} \\
AS.I31 & The second anomaly of a superior planet is completed when it returns to the same elongation from the Sun. & The synodic cycle of a superior planet is completed when it returns to the same elongation from the Sun. & AS10, AS23 & \textbf{H} \\
AS.I32 & Eclipses occur only when the Moon is near the ecliptic. & Eclipses occur only when the Moon is near the ecliptic. & AS17, AS14$^{\circ}$ & \textbf{H} \\
AS.I33 & The orb of Mars turns physically uniformly about the equant. & Interpretation: a solid shell rotates uniformly about an off-centre point. T2 has no shells; Kepler took the non-uniformity to be real. & AS08b$^{\ast}$, AS04a$^{\ast}$ & \textbf{U} \\
AS.I34 & The Sun is a planet like the others. & The Sun is a planet like the others. [Denied: it is the centre about which the planets move.] & AS13$^{\circ}$, AS01b & \textbf{X} \\
\bottomrule
\end{inftable}

\Needspace{10\baselineskip}\subsection*{MECH: inferences}
\addcontentsline{toc}{subsection}{MECH: inferences}
\begin{inftable}
ME.I01 & The mass of a body is the same in every frame. & The mass of a body is the same in every frame. [Holds on the rest-mass reading; on the relativistic-mass reading only for $v \ll c$.] & ME01$^{\circ}$ & \textbf{R} \\
ME.I02 & The mass of a composite is the sum of the masses of its parts. & The mass of a composite is the sum of the masses of its parts. [Holds on the relativistic-mass reading; for rest mass only when binding and kinetic energies are negligible.] & ME01$^{\circ}$ & \textbf{R} \\
ME.I03 & Mass is conserved in all interactions. & Mass is conserved in all interactions. [Holds on the relativistic-mass reading; rest mass is not conserved.] & ME01$^{\circ}$ & \textbf{R} \\
ME.I04 & Momentum is the product of mass and velocity. & Momentum is the product of mass and velocity. [Holds exactly on the relativistic-mass reading; on the rest-mass reading only for $v \ll c$.] & ME03$^{\circ}$, ME01$^{\circ}$, ME05$^{\circ}$ & \textbf{R} \\
ME.I05 & Momentum is conserved in an isolated system. & Relativistic momentum is conserved in an isolated system. & ME03$^{\circ}$ & \textbf{H} \\
ME.I06 & Force equals the rate of change of momentum. & The force on a charge equals the rate of change of its relativistic momentum. & ME02a, ME03$^{\circ}$ & \textbf{H} \\
ME.I07 & The gravitational force between two bodies varies as the product of their masses and inversely as the square of their distance. & Interpretation: a force reaches from the Sun to each planet in proportion to their masses and inversely as the square of the distance. In general relativity there is no such force; the Newtonian law is recovered as the weak-field limit of geodesic motion (ME.I08). & ME02b$^{\ast}$, ME01$^{\circ}$, ME11$^{\circ}$ & \textbf{U} \\
ME.I08 & In a gravitational field a body accelerates down the gradient of the potential. & In a weak static field a slow body accelerates down the gradient of $(g_{00}+1)c^2/2$. [Weak-field, slow-motion limit.] & ME12$^{\circ}$, ME21$^{\circ}$ & \textbf{R} \\
ME.I09 & Simultaneity of distant events is the same for all observers. & Simultaneity of distant events is the same in all inertial frames. [Holds as $v\,\Delta x/c^2 \to 0$.] & ME08$^{\circ}$ & \textbf{R} \\
ME.I10 & Velocities compose by addition. & Velocities compose by addition. [Holds for $v \ll c$; in general by the relativistic formula.] & ME05$^{\circ}$, ME10$^{\circ}$ & \textbf{R} \\
ME.I11 & Inertial frames are related by Galilean transformations. & Inertial frames are related by Galilean transformations. [Limit of the Lorentz transformation as $v/c \to 0$.] & ME09$^{\circ}$, ME10$^{\circ}$ & \textbf{R} \\
ME.I12 & The laws of mechanics are the same in all inertial frames. & The laws of physics are the same in all inertial frames. & ME09$^{\circ}$ & \textbf{H} \\
ME.I13 & Absolute space remains immovable and is not acted on by bodies. & Interpretation: a background space exists that is immovable and is not acted on by bodies. T2 recognizes no state of absolute rest; in general relativity the geometry is itself dynamical. & ME06a$^{\ast}$ & \textbf{U} \\
ME.I14 & The concavity of the water in the rotating bucket reveals its true rotation. & The concavity of the water reveals rotation relative to the local inertial structure. & ME06b & \textbf{H} \\
ME.I15 & Absolute time flows equably for all observers. & Interpretation: one time flows at the same rate for every observer. T2 has proper times along worldlines and frame-relative coordinate times. & ME07a$^{\ast}$ & \textbf{U} \\
ME.I16 & Clocks once synchronized remain synchronized whatever their relative motion. & Clocks once synchronized remain synchronized whatever their relative motion. [Holds as $v/c \to 0$ and in weak fields.] & ME07b, ME08$^{\circ}$ & \textbf{R} \\
ME.I17 & The length of a body does not depend on its motion. & The length of a body does not depend on its motion. [Holds as $v/c \to 0$.] & ME14$^{\circ}$ & \textbf{R} \\
ME.I18 & Light travels with a fixed speed relative to the ether. & Interpretation: there is a medium relative to which light has its speed. T2 has no such medium; light has speed $c$ in every inertial frame. & ME15$^{\circ}$, ME16a$^{\ast}$ & \textbf{U} \\
ME.I19 & The measured speed of light depends on the observer's motion. & The measured speed of light depends on the observer's motion. [Denied: it is $c$ in every inertial frame.] & ME15$^{\circ}$, ME05$^{\circ}$ & \textbf{X} \\
ME.I20 & Electromagnetic disturbances in the ether propagate with the speed of light. & Electromagnetic disturbances in the field propagate with speed $c$. & ME16b, ME15$^{\circ}$ & \textbf{H} \\
ME.I21 & A planet attracted by the Sun alone describes a fixed ellipse. & A planet orbiting the Sun alone follows a fixed ellipse. [Holds in the weak-field limit; general relativity adds a perihelion advance.] & ME13$^{\circ}$, ME11$^{\circ}$ & \textbf{R} \\
ME.I22 & Gravitational and inertial mass are proportional. & Gravitational and inertial mass are proportional (the equivalence principle). & ME01$^{\circ}$, ME11$^{\circ}$ & \textbf{H} \\
ME.I23 & Gravitation acts instantaneously across any distance. & Interpretation: gravitation acts at once across any distance. In T2 changes in the gravitational field propagate at the speed of light. & ME23a$^{\ast}$, ME11$^{\circ}$ & \textbf{U} \\
ME.I24 & An impulse given to one end of a rigid body moves the other end at once. & Interpretation: a body could transmit a push instantaneously. T2 excludes signals faster than light, and so excludes perfectly rigid bodies. & ME22a$^{\ast}$ & \textbf{U} \\
ME.I25 & A rod moved slowly from one place to another keeps its length. & A rod moved slowly from place to place keeps its proper length. & ME22b, ME14$^{\circ}$ & \textbf{H} \\
ME.I26 & Kinetic energy is one half of mass times velocity squared. & Kinetic energy is one half of mass times velocity squared. [Leading term of $(\gamma-1)m_0c^2$ for $v \ll c$.] & ME04$^{\circ}$, ME01$^{\circ}$, ME05$^{\circ}$ & \textbf{R} \\
ME.I27 & Energy is conserved in an isolated system. & Energy is conserved in an isolated system. & ME20 & \textbf{H} \\
ME.I28 & The weight of a body is the gravitational force on it. & Interpretation: weight is the pull of gravity on a body. In general relativity no force acts on a freely falling body; what a balance registers is the support force. & ME19a$^{\ast}$, ME02b$^{\ast}$ & \textbf{U} \\
ME.I29 & The centre of mass of an isolated system moves uniformly. & The centre of energy of an isolated system moves uniformly in any inertial frame. & ME18 & \textbf{H} \\
ME.I30 & Acceleration is the same in all inertial frames. & Acceleration is the same in all inertial frames. [Holds as $v/c \to 0$.] & ME21$^{\circ}$, ME09$^{\circ}$ & \textbf{R} \\
ME.I31 & The strength of gravitation is fixed by one universal constant. & The strength of gravitation is fixed by one universal constant $G$. & ME17, ME11$^{\circ}$ & \textbf{H} \\
ME.I32 & A spring balance holding a body at rest on the Earth registers its weight. & A spring balance holding a body at rest on the Earth registers the support force that keeps it from free fall. & ME19b & \textbf{H} \\
ME.I33 & The Earth's motion through the ether should alter the round-trip time of light along perpendicular arms. & Interpretation: the Earth moves through a medium, which should alter the round-trip time of light. In T2 there is no medium to move through, and no alteration is expected. & ME16c$^{\ast}$, ME15$^{\circ}$ & \textbf{U} \\
ME.I34 & Bodies falling freely together show no relative acceleration in a uniform field. & Bodies falling freely together show no relative acceleration in a uniform field. & ME09$^{\circ}$, ME21$^{\circ}$ & \textbf{H} \\
ME.I35 & The centre of the world is at rest in absolute space. & Interpretation: the centre of the solar system is at absolute rest. T2 has no absolute rest. & ME06c$^{\ast}$ & \textbf{U} \\
\bottomrule
\end{inftable}

\Needspace{10\baselineskip}\subsection*{HEAT: inferences}
\addcontentsline{toc}{subsection}{HEAT: inferences}
\begin{inftable}
HE.I01 & Caloric is a subtle fluid whose particles repel one another and are attracted by ordinary matter. & Interpretation: heat is a self-repelling fluid attracted by matter. T2 has no such fluid; heat is energy transferred. & HE01a$^{\ast}$ & \textbf{U} \\
HE.I02 & The total quantity of caloric is conserved. & Interpretation: the quantity of a heat-substance is fixed. In T2 energy is conserved but heat is not: it is converted to and from work. & HE01c$^{\ast}$ & \textbf{U} \\
HE.I03 & When hot and cold water are mixed, the heat lost by the one equals the heat gained by the other. & When hot and cold water are mixed (no work done), the heat lost by the one equals the heat gained by the other. & HE02$^{\circ}$, HE03$^{\circ}$ & \textbf{H} \\
HE.I04 & Heat passes spontaneously from hotter to colder bodies. & Heat passes spontaneously from hotter to colder bodies. & HE02$^{\circ}$, HE03$^{\circ}$ & \textbf{H} \\
HE.I05 & Ice at its melting point absorbs heat without rising in temperature. & Ice at its melting point absorbs heat without rising in temperature. & HE04$^{\circ}$, HE03$^{\circ}$ & \textbf{H} \\
HE.I06 & Latent heat is caloric combined with the substance. & Interpretation: the absorbed heat is held combined in the substance. In T2 it goes into the molecular potential energy of the new phase. & HE04$^{\circ}$, HE07a$^{\ast}$ & \textbf{U} \\
HE.I07 & The temperature rise of a body is the heat added divided by its heat capacity. & The temperature rise of a body is the heat added divided by its heat capacity. & HE05$^{\circ}$, HE02$^{\circ}$, HE03$^{\circ}$ & \textbf{H} \\
HE.I08 & Specific heat is the capacity of a body for caloric. & Interpretation: specific heat measures how much heat-fluid a body can hold. In T2 it is the energy needed per degree per unit mass. & HE05$^{\circ}$, HE01a$^{\ast}$ & \textbf{U} \\
HE.I09 & Air heats on sudden compression and cools on sudden expansion. & Air heats on sudden compression and cools on sudden expansion. & HE10, HE18$^{\circ}$, HE03$^{\circ}$ & \textbf{H} \\
HE.I10 & Sound compresses air adiabatically, so its speed is $\sqrt{\gamma p/\rho}$. & Sound compresses air adiabatically, so its speed is $\sqrt{\gamma p/\rho}$. & HE12, HE11, HE10 & \textbf{H} \\
HE.I11 & The specific heat of a gas at constant pressure exceeds that at constant volume. & The specific heat of a gas at constant pressure exceeds that at constant volume. & HE05$^{\circ}$, HE11, HE18$^{\circ}$ & \textbf{H} \\
HE.I12 & The efficiency of a reversible engine depends only on the temperatures of its reservoirs. & The efficiency of a reversible engine depends only on the temperatures of its reservoirs. & HE13$^{\circ}$, HE14, HE03$^{\circ}$ & \textbf{H} \\
HE.I13 & In a cycle, the heat given to the cold reservoir equals the heat taken from the hot. & In a cycle, the heat given to the cold reservoir equals the heat taken from the hot. [Denied: it is less by the work done.] & HE02$^{\circ}$, HE13$^{\circ}$ & \textbf{X} \\
HE.I14 & Motive power is produced by the fall of caloric through a temperature difference, as water falls through a wheel. & Work is produced by the passage of entropy through a temperature difference, $W = S(T_1 - T_2)$ for a reversible engine. [Holds on Callendar's reading of the token; fails on Kelvin's.] & HE01d, HE13$^{\circ}$, HE03$^{\circ}$ & \textbf{R} \\
HE.I15 & Frictional heat is caloric released from bodies whose capacity is reduced by abrasion. & Interpretation: frictional heat is heat-fluid squeezed out of abraded bodies. In T2 it is work converted into heat. & HE15$^{\circ}$, HE05$^{\circ}$, HE01a$^{\ast}$ & \textbf{U} \\
HE.I16 & Heat is conducted with a flux proportional to the temperature gradient. & Heat is conducted with a flux proportional to the temperature gradient. & HE08, HE03$^{\circ}$ & \textbf{H} \\
HE.I17 & All bodies emit radiant heat at all temperatures, and an enclosure comes to equilibrium by exchange. & All bodies emit thermal radiation at all temperatures, and an enclosure comes to equilibrium by exchange. & HE09$^{\circ}$, HE16 & \textbf{H} \\
HE.I18 & Radiant heat is caloric in free flight. & Interpretation: radiant heat is the heat-fluid travelling freely. In T2 it is electromagnetic radiation. & HE09$^{\circ}$, HE01a$^{\ast}$ & \textbf{U} \\
HE.I19 & Bodies in contact come to a common temperature. & Bodies in contact come to a common temperature. & HE16, HE03$^{\circ}$ & \textbf{H} \\
HE.I20 & Gases at constant pressure expand by the same fraction for the same rise in temperature. & Gases at constant pressure expand by the same fraction for the same rise in temperature. [Holds for ideal gases.] & HE17$^{\circ}$, HE18$^{\circ}$, HE03$^{\circ}$ & \textbf{R} \\
HE.I21 & The elasticity of a gas arises from static repulsion between caloric atmospheres. & Interpretation: the spring of a gas comes from static repulsion of heat-fluid atmospheres. In T2 pressure comes from molecular impacts. & HE19$^{\circ}$, HE01a$^{\ast}$, HE18$^{\circ}$ & \textbf{U} \\
HE.I22 & At constant temperature, the elasticity of a gas varies inversely as its volume. & At constant temperature the pressure of a gas varies inversely as its volume. [Holds for ideal gases.] & HE19$^{\circ}$, HE18$^{\circ}$ & \textbf{R} \\
HE.I23 & Respiration is a slow combustion and the source of animal heat. & Respiration is a slow combustion and the source of body heat. & HE23, HE02$^{\circ}$ & \textbf{H} \\
HE.I24 & The ice calorimeter measures heat by the quantity of ice melted. & The ice calorimeter measures heat by the quantity of ice melted. & HE21, HE02$^{\circ}$, HE04$^{\circ}$ & \textbf{H} \\
HE.I25 & Evaporation absorbs heat. & Evaporation absorbs heat. & HE20$^{\circ}$, HE04$^{\circ}$ & \textbf{H} \\
HE.I26 & The thermometer measures the quantity of free caloric in a body. & Interpretation: the thermometer reads the density of uncombined heat-fluid. In T2 it reads temperature, which is not a quantity of anything. & HE22$^{\circ}$, HE06a$^{\ast}$ & \textbf{U} \\
HE.I27 & In Laplace's calculation the caloric released by compression stays in the compressed parcel. & In Laplace's calculation the heat released by compression stays in the compressed parcel (the compression is adiabatic). & HE01b, HE10 & \textbf{H} \\
HE.I28 & Heat cannot be generated without limit from a finite body. & Heat cannot be generated without limit from a finite body. [Denied: work can be converted into heat without limit, as in Rumford's boring.] & HE02$^{\circ}$, HE15$^{\circ}$ & \textbf{X} \\
HE.I29 & A vapour is a liquid combined with caloric. & Interpretation: a vapour is a liquid chemically combined with heat-fluid. In T2 it is the gaseous phase. & HE20$^{\circ}$, HE07a$^{\ast}$ & \textbf{U} \\
HE.I30 & Thermometers agree on the order of bodies by hotness. & Thermometers agree on the order of bodies by temperature. & HE22$^{\circ}$, HE03$^{\circ}$ & \textbf{H} \\
\bottomrule
\end{inftable}

\Needspace{10\baselineskip}\subsection*{QUAN: inferences}
\addcontentsline{toc}{subsection}{QUAN: inferences}
\begin{inftable}
QU.I01 & An electron in an atom moves on a definite orbit about the nucleus. & Interpretation: the bound electron follows a closed path about the nucleus. In T2 a stationary state assigns no path to the electron. & QU02a$^{\ast}$, QU06$^{\circ}$, QU07$^{\circ}$, QU08$^{\circ}$ & \textbf{U} \\
QU.I02 & An electron revolving in an atom radiates continuously and spirals into the nucleus. & Interpretation: the bound electron loses energy continuously and falls into the nucleus. In T2 the ground state is stable and does not radiate. & QU02a$^{\ast}$, QU10$^{\circ}$, QU08$^{\circ}$ & \textbf{U} \\
QU.I03 & An accelerated charge radiates with power given by the Larmor formula. & An accelerated charge radiates with power given by the Larmor formula. [Holds for macroscopic currents and free charges; not for stationary states.] & QU10$^{\circ}$, QU19$^{\circ}$ & \textbf{R} \\
QU.I04 & The frequencies emitted by an atom are those of its electrons' orbital motions and their harmonics. & Interpretation: the emitted frequencies are those of periodic electron motions. In T2 they are differences of stationary-state energies divided by $h$. & QU11$^{\circ}$, QU02a$^{\ast}$ & \textbf{U} \\
QU.I05 & A very highly excited electron radiates at its orbital frequency. & A circulating Rydberg wave packet radiates at its orbital frequency. [Holds in the limit of high quantum numbers: Bohr's correspondence principle.] & QU02b, QU11$^{\circ}$ & \textbf{R} \\
QU.I06 & Every particle has at every time a definite position and momentum. & Every particle has at every time a definite value of the position and of the momentum observable. [Denied in the standard formulation.] & QU03$^{\circ}$, QU04$^{\circ}$ & \textbf{X} \\
QU.I07 & Position and momentum can be measured together with arbitrary precision. & Position and momentum can be measured together with arbitrary precision. [Denied: $\Delta x\,\Delta p \geq \hbar/2$.] & QU03$^{\circ}$, QU04$^{\circ}$ & \textbf{X} \\
QU.I08 & The trajectory of a body is fixed by its position and momentum at one time together with the forces. & The centre of a localized state moves as the classical equations say. [Ehrenfest's theorem: holds for well-localized states in slowly varying potentials.] & QU01c, QU03$^{\circ}$, QU04$^{\circ}$ & \textbf{R} \\
QU.I09 & A charged particle in supersaturated vapour leaves a track along its trajectory. & A charged particle in supersaturated vapour leaves a nearly straight track of condensed droplets. & QU01b, QU06$^{\circ}$ & \textbf{H} \\
QU.I10 & The electron in an atom can have any energy in a continuous range. & The energy of a bound electron can take any value in a continuous range. [Denied: bound-state energies are discrete.] & QU05$^{\circ}$, QU06$^{\circ}$, QU07$^{\circ}$ & \textbf{X} \\
QU.I11 & In equilibrium every degree of freedom has mean energy $kT/2$. & In equilibrium each quadratic degree of freedom has mean energy $kT/2$. [Holds for $kT$ large compared with the level spacing.] & QU12$^{\circ}$, QU05$^{\circ}$ & \textbf{R} \\
QU.I12 & The spectral density of cavity radiation is proportional to $\nu^2 T$ (Rayleigh--Jeans). & The spectral density of cavity radiation is proportional to $\nu^2 T$. [Holds for $h\nu \ll kT$.] & QU13$^{\circ}$, QU12$^{\circ}$ & \textbf{R} \\
QU.I13 & The atomic heat of every solid element is about 6 calories per degree. & The molar heat capacity of every solid element is about $3R$. [Holds at temperatures high compared with its characteristic temperature.] & QU14$^{\circ}$, QU12$^{\circ}$ & \textbf{R} \\
QU.I14 & The energy of photoelectrons increases with the intensity of the light. & The energy of photoelectrons increases with the intensity of the light. [Denied: it depends on the frequency; the intensity fixes their number.] & QU15$^{\circ}$, QU09$^{\circ}$ & \textbf{X} \\
QU.I15 & Light shows interference and diffraction. & Electromagnetic radiation shows interference and diffraction. & QU09$^{\circ}$ & \textbf{H} \\
QU.I16 & The energy of light is spread continuously over the wave front. & The energy of radiation is spread continuously over the wave front. [Denied: it is exchanged with matter in quanta $h\nu$.] & QU09$^{\circ}$, QU05$^{\circ}$ & \textbf{X} \\
QU.I17 & Electrons have a definite charge and mass. & Electrons have a definite charge and mass. & QU06$^{\circ}$, QU19$^{\circ}$ & \textbf{H} \\
QU.I18 & Charge is conserved and comes in multiples of $e$. & Charge is conserved and comes in multiples of $e$ (for free particles). & QU19$^{\circ}$ & \textbf{H} \\
QU.I19 & Nearly all the mass and the positive charge of an atom are in a small nucleus. & Nearly all the mass and the positive charge of an atom are in a small nucleus. & QU08$^{\circ}$, QU07$^{\circ}$ & \textbf{H} \\
QU.I20 & Alpha particles scattered by a nucleus follow Rutherford's angular distribution. & Alpha particles scattered by a nucleus follow Rutherford's distribution (which the quantum treatment of Coulomb scattering reproduces exactly). & QU21, QU08$^{\circ}$ & \textbf{H} \\
QU.I21 & X-rays are diffracted by crystals as short electromagnetic waves. & X-rays are diffracted by crystals as short electromagnetic waves. & QU20, QU09$^{\circ}$ & \textbf{H} \\
QU.I22 & The angular momentum of a body can take any value. & The angular momentum of a body can take any value. [Denied: its components have discrete eigenvalues.] & QU16$^{\circ}$ & \textbf{X} \\
QU.I23 & The angular momentum of an isolated system is conserved. & The angular momentum of an isolated system is conserved. & QU16$^{\circ}$ & \textbf{H} \\
QU.I24 & An atomic magnetic moment can point in any direction relative to a field. & An atomic magnetic moment can point in any direction relative to a field. [Denied: Stern and Gerlach found discrete components.] & QU17$^{\circ}$, QU16$^{\circ}$ & \textbf{X} \\
QU.I25 & The state of a system is a point in phase space, evolving deterministically. & Interpretation: the full state is a set of exact positions and momenta, evolving deterministically. In T2 the state is a vector in Hilbert space. & QU18a$^{\ast}$ & \textbf{U} \\
QU.I26 & The macroscopic state of a gas is fixed by its pressure, volume and temperature. & The thermodynamic macrostate of a gas is fixed by its pressure, volume and temperature. & QU18b & \textbf{H} \\
QU.I27 & An antenna current oscillating at a given frequency radiates at that frequency. & An antenna current oscillating at a given frequency radiates at that frequency. & QU10$^{\circ}$, QU09$^{\circ}$ & \textbf{H} \\
QU.I28 & The spectral lines of hydrogen satisfy the Balmer--Rydberg formula. & The transition frequencies of hydrogen satisfy the Balmer--Rydberg formula. & QU11$^{\circ}$, QU07$^{\circ}$ & \textbf{H} \\
QU.I29 & A planet's orbit is a definite path in space. & A planet's orbit is a definite path in space. [Holds for macroscopic bodies to any accuracy that can be tested.] & QU01c & \textbf{R} \\
QU.I30 & The electron in an atom traces a continuous trajectory fixed by the Coulomb force of the nucleus. & Interpretation: the bound electron has a continuous path determined by the Coulomb force. In T2 the Coulomb potential fixes the stationary states, which assign no path. & QU01a$^{\ast}$, QU08$^{\circ}$ & \textbf{U} \\
\bottomrule
\end{inftable}

\Needspace{10\baselineskip}\subsection*{BIOL: inferences}
\addcontentsline{toc}{subsection}{BIOL: inferences}
\begin{inftable}
BI.I01 & Species are immutable. & Species do not change. [Denied: species descend with modification.] & BI01$^{\circ}$ & \textbf{X} \\
BI.I02 & Varieties differ from species in not being permanent. & Varieties differ from species in not being permanent. [Denied: varieties are incipient species, and the difference is one of degree.] & BI02$^{\circ}$, BI01$^{\circ}$ & \textbf{X} \\
BI.I03 & The variation of a species is confined within fixed limits about its type. & The heritable variation within a species is confined within fixed limits. [Denied.] & BI23$^{\circ}$, BI01$^{\circ}$ & \textbf{X} \\
BI.I04 & The natural system arranges organisms in groups subordinate to groups. & Genealogical classification arranges organisms in groups subordinate to groups. & BI04$^{\circ}$, BI03 & \textbf{H} \\
BI.I05 & The natural system expresses the plan of the Creator. & Interpretation: the hierarchy of groups expresses a design. In T2 it expresses community of descent. & BI04$^{\circ}$, BI05a$^{\ast}$ & \textbf{U} \\
BI.I06 & The members of each great class share a unity of plan. & The members of each great class share an ancestral pattern of structure. & BI05b, BI04$^{\circ}$ & \textbf{H} \\
BI.I07 & The archetype is an idea in the divine mind. & Interpretation: the pattern is an idea in the mind of a designer. In T2 it is the structure of an ancestor. & BI06a$^{\ast}$ & \textbf{U} \\
BI.I08 & All vertebrate skeletons are modifications of one archetypal pattern. & All vertebrate skeletons are modifications of the skeleton of their common ancestor. & BI06b, BI07 & \textbf{H} \\
BI.I09 & The forelimbs of man, bat and porpoise are homologous. & The forelimbs of man, bat and porpoise are homologous. & BI07 & \textbf{H} \\
BI.I10 & The wings of birds and insects are analogous and not homologous. & The wings of birds and insects are analogous and not homologous. & BI08, BI07 & \textbf{H} \\
BI.I11 & Rudimentary organs exist to complete the plan. & Interpretation: useless organs are present to complete a design. In T2 they are inherited remnants of organs useful to ancestors. & BI11$^{\circ}$, BI05a$^{\ast}$ & \textbf{U} \\
BI.I12 & Rudimentary organs are useless or nearly so. & Vestigial organs are useless or nearly so. & BI11$^{\circ}$ & \textbf{H} \\
BI.I13 & Embryos of different species in a class resemble each other more than the adults do. & Embryos of different species in a class resemble each other more than the adults do. & BI12 & \textbf{H} \\
BI.I14 & Every organism is perfectly adapted to its conditions of life. & Every organism is perfectly adapted to its conditions. [Denied: adaptation is relative to competitors and often imperfect.] & BI09$^{\circ}$ & \textbf{X} \\
BI.I15 & Organs such as the eye are fitted to their functions. & Organs such as the eye are adapted to their functions. & BI09$^{\circ}$ & \textbf{H} \\
BI.I16 & Each species arose in a single region and spread from it. & Each species arose in a single region and spread from it. & BI10b, BI01$^{\circ}$ & \textbf{H} \\
BI.I17 & Each species was created by a separate act in its centre of creation. & Interpretation: each species was brought into being by a separate act. T2 recognizes no such acts; species arise by descent. & BI10a$^{\ast}$, BI01$^{\circ}$ & \textbf{U} \\
BI.I18 & Whole faunas were destroyed by sudden catastrophes and replaced. & Whole faunas were destroyed by sudden catastrophes and replaced. [Denied in Darwin's gradualism.] & BI13$^{\circ}$ & \textbf{X} \\
BI.I19 & Fossils appear in a determinate order in successive strata. & Fossils appear in a determinate order in successive strata. & BI15 & \textbf{H} \\
BI.I20 & Extinct and living forms in the same continent are allied. & Extinct and living forms in the same continent are allied. & BI14, BI01$^{\circ}$ & \textbf{H} \\
BI.I21 & Living beings form a single linear scale without gaps. & Interpretation: organisms form one gapless linear series. T2 has a branching tree with extinct lineages. & BI16a$^{\ast}$ & \textbf{U} \\
BI.I22 & Domestic races vary more than wild species. & Domestic breeds vary more than wild species. & BI17, BI23$^{\circ}$, BI01$^{\circ}$ & \textbf{H} \\
BI.I23 & Hybrids between distinct species are generally sterile. & Hybrids between distinct species are generally sterile. & BI18$^{\circ}$, BI01$^{\circ}$ & \textbf{H} \\
BI.I24 & Hybrid sterility is a special provision to prevent the confusion of species. & Hybrid sterility is a special provision against the confusion of species. [Denied: it is incidental on other differences (Origin, ch.~8).] & BI18$^{\circ}$ & \textbf{X} \\
BI.I25 & More individuals are born than can survive. & More individuals are born than can survive. & BI19 & \textbf{H} \\
BI.I26 & Each species keeps its station in a designed balance of nature. & Each species keeps its station in a designed balance. [Denied: the relations shift as species change and go extinct.] & BI20$^{\circ}$, BI01$^{\circ}$ & \textbf{X} \\
BI.I27 & Instincts are fixed endowments, perfect for their purpose. & Instincts are fixed endowments, perfect for their purpose. [Denied: they vary and are shaped by selection.] & BI21$^{\circ}$ & \textbf{X} \\
BI.I28 & Sports occasionally appear in domestic breeds. & Large discontinuous variations occasionally appear in domestic breeds. & BI22, BI17 & \textbf{H} \\
BI.I29 & Higher organisms have more differentiated parts. & Organisms with more differentiated parts are those called higher. & BI24 & \textbf{H} \\
BI.I30 & Extinct species do not reappear. & Extinct species do not reappear. & BI13$^{\circ}$, BI01$^{\circ}$ & \textbf{H} \\
BI.I31 & The species of a genus resemble each other more than species of different genera. & The species of a genus resemble each other more than species of different genera. & BI03, BI01$^{\circ}$ & \textbf{H} \\
\bottomrule
\end{inftable}

\Needspace{10\baselineskip}\subsection*{PHYS: inferences}
\addcontentsline{toc}{subsection}{PHYS: inferences}
\begin{inftable}
GA.I01 & Blood is made continuously in the liver from the chyle of digested food. & Blood is made continuously in the liver from chyle. [Denied: the same blood circulates; Harvey's quantitative argument.] & GA06$^{\circ}$, GA16$^{\circ}$, GA01$^{\circ}$ & \textbf{X} \\
GA.I02 & Venous blood is distributed from the liver through the veins and consumed by the parts it nourishes. & Venous blood flows from the liver through the veins to the parts and is consumed there. [Denied: it flows back to the heart.] & GA01$^{\circ}$, GA24$^{\circ}$, GA19$^{\circ}$ & \textbf{X} \\
GA.I03 & Blood passes from the right to the left ventricle through pores in the septum. & Interpretation: blood crosses the septum through hidden channels. T2 recognizes none; blood passes from right to left through the lungs. & GA08a$^{\ast}$, GA07$^{\circ}$ & \textbf{U} \\
GA.I04 & The arterial vein carries blood from the right ventricle to the lungs. & The pulmonary artery carries blood from the right ventricle to the lungs. & GA09$^{\circ}$, GA20$^{\circ}$ & \textbf{H} \\
GA.I05 & The venous artery carries air from the lungs to the left ventricle. & The pulmonary vein carries air from the lungs to the left ventricle. [Denied: it carries blood.] & GA10$^{\circ}$, GA20$^{\circ}$ & \textbf{X} \\
GA.I06 & Vital spirit is elaborated in the left ventricle from air and blood. & Interpretation: an active principle is made in the left ventricle from air and blood. T2 recognizes no such principle. & GA03a$^{\ast}$, GA07$^{\circ}$ & \textbf{U} \\
GA.I07 & The arteries carry blood mixed with vital spirit to the whole body. & Interpretation: the arteries distribute blood with that principle. T2 accepts the distribution of arterial blood but not the principle. & GA02$^{\circ}$, GA03a$^{\ast}$ & \textbf{U} \\
GA.I08 & Arterial blood is brighter than venous blood. & Arterial blood is brighter than venous blood. & GA02$^{\circ}$, GA01$^{\circ}$ & \textbf{H} \\
GA.I09 & The active motion of the heart is diastole, by which it attracts blood. & Interpretation: the heart draws blood in by an active power of dilation. In T2 the active phase is systole. & GA11$^{\circ}$, GA18a$^{\ast}$ & \textbf{U} \\
GA.I10 & In systole the heart expels blood into the aorta. & In systole the heart expels blood into the aorta. & GA12, GA25 & \textbf{H} \\
GA.I11 & The arteries pulsate by a faculty conveyed from the heart through their coats. & Interpretation: the arteries pulsate by a power of their own conveyed through their walls. In T2 they are distended by the blood the heart ejects. & GA13$^{\circ}$, GA18a$^{\ast}$ & \textbf{U} \\
GA.I12 & The valves of the veins retard the flow of blood to the extremities. & The venous valves retard the flow of blood to the extremities. [Denied: they prevent flow away from the heart.] & GA17$^{\circ}$ & \textbf{X} \\
GA.I13 & The heart is the seat of the innate heat. & Interpretation: the heart is the seat of a vital heat. T2 recognizes body heat but not a vital principle seated in the heart. & GA07$^{\circ}$, GA14a$^{\ast}$ & \textbf{U} \\
GA.I14 & Respiration cools the heart and supplies the matter of vital spirit. & Interpretation: breathing cools the vital heat and supplies the matter of a vital principle. Harvey left the function of respiration open; T2 as continued explains it otherwise. & GA21$^{\circ}$, GA20$^{\circ}$, GA03a$^{\ast}$ & \textbf{U} \\
GA.I15 & Sooty vapours from the left ventricle are expelled through the venous artery to the lungs. & Interpretation: waste passes back from the left ventricle to the lungs. T2 excludes such backflow: the mitral valve prevents it. & GA15a$^{\ast}$, GA10$^{\circ}$ & \textbf{U} \\
GA.I16 & Expired air carries away waste products. & Expired air carries away waste products. & GA15b, GA21$^{\circ}$ & \textbf{H} \\
GA.I17 & Animal spirit is elaborated from vital spirit in the rete mirabile. & Interpretation: a further principle is elaborated in a network at the base of the human brain. T2 recognizes neither the principle nor the human network. & GA05a$^{\ast}$, GA26a$^{\ast}$, GA03a$^{\ast}$ & \textbf{U} \\
GA.I18 & What passes along the nerves from the brain produces motion of the muscles. & Nerve conduction from the brain produces motion of the muscles. & GA05b & \textbf{H} \\
GA.I19 & The ox has a network of arteries at the base of its brain. & The ox has a network of arteries at the base of its brain (the carotid rete). & GA26b & \textbf{H} \\
GA.I20 & Natural spirit is elaborated in the liver and carried by the veins. & Interpretation: a principle of nutrition is elaborated in the liver. T2 recognizes no such principle. & GA04a$^{\ast}$, GA06$^{\circ}$ & \textbf{U} \\
GA.I21 & Arteries and veins communicate at their extremities. & Arteries and veins communicate at their extremities (through capillaries, seen by Malpighi in 1661). & GA22 & \textbf{H} \\
GA.I22 & The portal vein carries chyle from the intestines to the liver. & The hepatic portal vein carries chyle to the liver. [Denied: it carries blood with absorbed nutrients; the chyle passes by the lacteals.] & GA23$^{\circ}$, GA16$^{\circ}$, GA06$^{\circ}$ & \textbf{X} \\
GA.I23 & The aorta carries arterial blood from the left ventricle to the body. & The aorta carries arterial blood from the left ventricle to the body. & GA25, GA02$^{\circ}$ & \textbf{H} \\
GA.I24 & A ligature on an artery stops the pulse beyond it. & A ligature on an artery stops the pulse beyond it. & GA13$^{\circ}$ & \textbf{H} \\
GA.I25 & Each part attracts the nourishment it needs by its attractive faculty. & Interpretation: each part draws nourishment by a power of its own. T2 recognizes no such power. & GA18a$^{\ast}$, GA19$^{\circ}$ & \textbf{U} \\
GA.I26 & The parts are nourished by matter carried to them in the blood. & The parts are nourished by matter carried to them in the blood. & GA19$^{\circ}$, GA01$^{\circ}$ & \textbf{H} \\
GA.I27 & The lungs are nourished by blood brought by the arterial vein. & The lungs are nourished by blood brought by the pulmonary artery. [Denied: that blood passes through the lungs to the left heart.] & GA09$^{\circ}$, GA20$^{\circ}$, GA19$^{\circ}$ & \textbf{X} \\
GA.I28 & The living body is warm, and grows cold at death. & The living body is warm, and grows cold at death. & GA14b & \textbf{H} \\
GA.I29 & Blood flows out of a cut vein and a cut artery. & Blood flows out of a cut vein and a cut artery. & GA01$^{\circ}$, GA02$^{\circ}$ & \textbf{H} \\
GA.I30 & In systole the heart becomes harder and smaller. & In systole the heart becomes harder and smaller. & GA12, GA07$^{\circ}$ & \textbf{H} \\
\bottomrule
\end{inftable}

\Needspace{10\baselineskip}\subsection*{OPTI: inferences}
\addcontentsline{toc}{subsection}{OPTI: inferences}
\begin{inftable}
OP.I01 & Light consists of corpuscles emitted by luminous bodies. & Interpretation: light is a stream of small bodies. The wave theory recognizes none. & OP02a$^{\ast}$ & \textbf{U} \\
OP.I02 & Rays travel in straight lines in a uniform medium. & Light travels in straight lines in a uniform medium. [Holds when apertures and obstacles are large compared with the wavelength.] & OP01$^{\circ}$ & \textbf{R} \\
OP.I03 & Refraction toward the normal arises from the attraction of the corpuscles by the denser medium. & Interpretation: refraction is due to an attraction of light-particles by the denser medium. In T2 it is due to the lower wave speed in that medium. & OP02a$^{\ast}$, OP05$^{\circ}$ & \textbf{U} \\
OP.I04 & Light travels faster in water than in air. & Light travels faster in water than in air. [Denied: slower, as Foucault found in 1850.] & OP06$^{\circ}$ & \textbf{X} \\
OP.I05 & Sunlight is a mixture of rays of different refrangibility. & Sunlight is a superposition of light of different wavelengths, refracted differently. & OP13, OP05$^{\circ}$, OP01$^{\circ}$ & \textbf{H} \\
OP.I06 & Homogeneal light keeps its colour and refrangibility under reflection and refraction. & Monochromatic light keeps its wavelength in reflection and refraction (its frequency is unchanged). & OP14, OP04, OP05$^{\circ}$ & \textbf{H} \\
OP.I07 & Each corpuscle passes alternately into fits of easy reflection and of easy transmission. & Interpretation: each light-particle alternates between two dispositions. In T2 the periodicity belongs to the wave, not to particles. & OP03a$^{\ast}$, OP02a$^{\ast}$ & \textbf{U} \\
OP.I08 & The thicknesses of the film at the bright rings are as the odd numbers 1, 3, 5, \ldots & The film thicknesses at the bright rings in reflected light are as the odd numbers 1, 3, 5 (in quarter-wavelengths). & OP11, OP03b & \textbf{H} \\
OP.I09 & The interval of the fits is greater for red than for violet light. & The wavelength is greater for red than for violet light. & OP03b, OP04 & \textbf{H} \\
OP.I10 & Light passing near an edge is inflected and forms coloured fringes outside the shadow. & Light diffracted near an edge forms coloured fringes outside the shadow. & OP07$^{\circ}$, OP16$^{\circ}$ & \textbf{H} \\
OP.I11 & Inflection is caused by forces of the edge acting on the corpuscles. & Interpretation: the edge exerts forces on passing light-particles. In T2 fringes arise from the interference of waves. & OP07$^{\circ}$, OP02a$^{\ast}$ & \textbf{U} \\
OP.I12 & Light does not bend into the geometrical shadow. & Light does not enter the geometrical shadow. [Denied: the bright spot at the centre of a disk's shadow.] & OP16$^{\circ}$, OP01$^{\circ}$ & \textbf{X} \\
OP.I13 & The sides of a ray determine its refraction in a second Iceland crystal. & The polarization of a beam determines its refraction in a second calcite crystal. & OP08b, OP10 & \textbf{H} \\
OP.I14 & The sides are original properties of the corpuscles. & Interpretation: the sides are properties of the particles. In T2 polarization is the orientation of a transverse vibration. & OP08a$^{\ast}$ & \textbf{U} \\
OP.I15 & Light reflected from glass at a certain angle is polarized. & Light reflected from glass at a certain angle is polarized (Brewster's angle). & OP09$^{\circ}$, OP15$^{\circ}$ & \textbf{H} \\
OP.I16 & Iceland crystal divides a beam into two beams polarized at right angles. & Calcite divides a beam into two beams polarized at right angles. & OP10, OP09$^{\circ}$ & \textbf{H} \\
OP.I17 & Polarization consists in the alignment of the axes of the light particles. & Interpretation: polarization aligns the axes of light-particles. T2 has no particles to align. & OP09$^{\circ}$, OP02a$^{\ast}$ & \textbf{U} \\
OP.I18 & Starlight is aberrated by the ratio of the Earth's speed to the speed of light. & Starlight is aberrated by the ratio of the Earth's speed to the speed of light. & OP12 & \textbf{H} \\
OP.I19 & Intensity diminishes as the inverse square of the distance. & Intensity diminishes as the inverse square of the distance. & OP17 & \textbf{H} \\
OP.I20 & Reflection is caused by a repulsive power acting at a distance from the surface. & Interpretation: a repulsive force near the surface turns the particles back. In T2 reflection follows from the boundary conditions on the wave. & OP15$^{\circ}$, OP02a$^{\ast}$ & \textbf{U} \\
OP.I21 & The angle of reflection equals the angle of incidence. & The angle of reflection equals the angle of incidence. & OP15$^{\circ}$, OP01$^{\circ}$ & \textbf{H} \\
OP.I22 & For homogeneal light the sine of incidence bears a constant ratio to the sine of refraction. & For monochromatic light the ratio of the sines of incidence and refraction is constant. & OP05$^{\circ}$, OP14, OP01$^{\circ}$ & \textbf{H} \\
OP.I23 & A prism spreads sunlight into a spectrum from red, least refracted, to violet, most refracted. & A prism disperses sunlight into a spectrum from long wavelengths, least refracted, to short, most refracted. & OP18, OP13, OP04 & \textbf{H} \\
OP.I24 & The rings in transmitted light are complementary in colour to those in reflected light. & The rings in transmitted light are complementary to those in reflected light. & OP11, OP04 & \textbf{H} \\
OP.I25 & The colours of thin plates depend on the thickness of the plate. & The colours of thin films depend on their thickness through the path difference of the reflected waves. & OP11, OP03b & \textbf{H} \\
\bottomrule
\end{inftable}

\Needspace{10\baselineskip}\subsection*{MIAS: inferences}
\addcontentsline{toc}{subsection}{MIAS: inferences}
\begin{inftable}
MI.I01 & Epidemic diseases are caused by miasmata from decaying organic matter. & Interpretation: a poisonous emanation of decay, breathed in, causes epidemics. T2 recognizes no such single poison; epidemics are caused by specific organisms. & MI01a$^{\ast}$ & \textbf{U} \\
MI.I02 & Intermittent fevers are most prevalent in marshy districts and diminish when the marshes are drained. & Malaria is most prevalent in marshy districts and diminishes when the marshes are drained (which removes the breeding places of the mosquitoes that carry Plasmodium). & MI09, MI01b & \textbf{H} \\
MI.I03 & Zymotic diseases are produced by a process in the blood analogous to fermentation. & Interpretation: a chemical fermentation in the blood produces these diseases. In T2 they are produced by the multiplication of organisms in the body. & MI02$^{\circ}$, MI03c$^{\ast}$ & \textbf{U} \\
MI.I04 & Smallpox is communicated by contact with the sick or their clothing. & Smallpox is transmitted by contact with the sick or their clothing. & MI04, MI05a & \textbf{H} \\
MI.I05 & Contagia are living organisms. & Pathogenic agents are living organisms. & MI05b & \textbf{H} \\
MI.I06 & Contagia are chemical poisons that multiply by inducing decomposition. & Interpretation: the agents are chemical poisons multiplying by catalysis. T2 has no such poisons; the agents reproduce as organisms. & MI05c$^{\ast}$ & \textbf{U} \\
MI.I07 & Yeast is necessary for alcoholic fermentation. & Yeast cells are necessary for alcoholic fermentation. & MI03a, MI20$^{\circ}$ & \textbf{H} \\
MI.I08 & Fermentation is a decomposition propagated from a decomposing body. & Interpretation: fermentation is decomposition propagated from a decomposing body. In T2 it is the metabolism of living cells or the action of their enzymes. & MI20$^{\circ}$, MI03c$^{\ast}$ & \textbf{U} \\
MI.I09 & Putrefaction arises in dead organic matter from the action of air alone. & Microbial decomposition arises in dead organic matter from the action of air alone. [Denied: it requires organisms, which air may carry; sterile matter in filtered air does not putrefy.] & MI06$^{\circ}$ & \textbf{X} \\
MI.I10 & Microscopic organisms arise in infusions without parents. & Interpretation: organisms arise without parents. T2 (Pasteur, 1862) denies that they do. & MI07a$^{\ast}$ & \textbf{U} \\
MI.I11 & Boiled hay infusions can teem with organisms after cooling. & Boiled hay infusions can teem with organisms after cooling (grown from spores that survive boiling). & MI07b & \textbf{H} \\
MI.I12 & Epidemics follow constitutions of the atmosphere. & Interpretation: epidemics follow states of the air. T2 recognizes no such states as causes, though weather affects vectors and transmission. & MI08a$^{\ast}$ & \textbf{U} \\
MI.I13 & Removal of filth reduces the mortality of epidemic disease. & Removal of organic refuse reduces the mortality of infectious disease. & MI11$^{\circ}$ & \textbf{H} \\
MI.I14 & Disease arises from breathing the emanations of filth. & Interpretation: the smell of filth causes disease. T2 denies that smells cause disease; filth harbours pathogens. & MI11$^{\circ}$, MI01a$^{\ast}$ & \textbf{U} \\
MI.I15 & Sewer gas causes typhoid. & Interpretation: a gas from drains causes typhoid. In T2 typhoid is caused by Salmonella Typhi, transmitted by contaminated water and food. & MI12a$^{\ast}$ & \textbf{U} \\
MI.I16 & Cholera is communicated by morbid matter in the evacuations, swallowed in water. & Cholera is transmitted by Vibrio cholerae in the evacuations, swallowed in water. & MI13a, MI21$^{\circ}$ & \textbf{H} \\
MI.I17 & The cholera poison must mature in porous soil before it causes disease. & Interpretation: the cholera agent must ripen in the soil. T2 recognizes no such ripening. & MI13b$^{\ast}$ & \textbf{U} \\
MI.I18 & Cholera mortality in London in 1849 was lower in districts of higher elevation. & Cholera mortality in London in 1849 was lower in districts of higher elevation. & MI21$^{\circ}$ & \textbf{H} \\
MI.I31 & The cholera miasma was most concentrated in the low-lying districts by the Thames. & Vibrio cholerae was most prevalent in the low-lying districts by the Thames. [Holds for London in 1849, where those districts drew contaminated water.] & MI01c, MI21$^{\circ}$ & \textbf{R} \\
MI.I19 & Agents that destroy offensive smells prevent disease by destroying miasma. & Interpretation: deodorants prevent disease by destroying an emanation. In T2 some deodorants are disinfectants and prevent disease by killing microorganisms. & MI14$^{\circ}$, MI01a$^{\ast}$ & \textbf{U} \\
MI.I20 & Washing the hands in chloride of lime reduced mortality from puerperal fever. & Washing the hands in hypochlorite reduced mortality from puerperal sepsis. & MI14$^{\circ}$, MI17$^{\circ}$ & \textbf{H} \\
MI.I21 & Puerperal fever arises from an epidemic influence in the ward atmosphere. & Interpretation: an epidemic influence of the air causes puerperal fever. In T2 it is carried by the hands and instruments of attendants. & MI17$^{\circ}$, MI08a$^{\ast}$ & \textbf{U} \\
MI.I22 & Vaccination with cowpox protects against smallpox. & Vaccination with cowpox protects against smallpox. & MI15 & \textbf{H} \\
MI.I23 & Hospital gangrene is more common in crowded wards. & Necrotizing wound infection is more common in crowded wards. & MI16 & \textbf{H} \\
MI.I24 & Each zymotic disease breeds true: smallpox produces smallpox and never another disease. & Each infectious disease breeds true: variola produces smallpox and never another disease. & MI02$^{\circ}$, MI05a & \textbf{H} \\
MI.I25 & Continued fever is a single disease. & Typhus and typhoid are one disease. [Denied: distinguished clinically by William Jenner (1849--1851), aetiologically later.] & MI10$^{\circ}$ & \textbf{X} \\
MI.I26 & Poverty and intemperance predispose to epidemic disease. & Poverty and intemperance increase susceptibility to infectious disease. & MI19 & \textbf{H} \\
MI.I27 & Infection can be carried by clothes and bedding. & Infection can be carried by clothes and bedding. & MI18, MI04 & \textbf{H} \\
MI.I28 & Intermittent fever is cured by cinchona bark. & Malaria is cured by cinchona bark (quinine). & MI09 & \textbf{H} \\
MI.I29 & Diastase converts starch into sugar. & Amylase converts starch into sugar. & MI03b & \textbf{H} \\
MI.I30 & Putrid matter introduced into the blood produces fever. & Putrid matter introduced into the blood produces fever (by infection). & MI06$^{\circ}$, MI18 & \textbf{H} \\
\bottomrule
\end{inftable}

\Needspace{10\baselineskip}\subsection*{HERE: inferences}
\addcontentsline{toc}{subsection}{HERE: inferences}
\begin{inftable}
HR.I01 & Every cell throws off gemmules, which collect in the reproductive elements. & Interpretation: particles shed by every cell collect in the gametes. T2 has no such particles; the hereditary material is in the chromosomes of the germ line. & HR01a$^{\ast}$, HR17$^{\circ}$ & \textbf{U} \\
HR.I02 & The offspring blends the characters of the parents, and the blend is not undone in later generations. & Quantitative characters of offspring are intermediate, and the intermediacy persists in later generations. [Denied: segregation restores parental types.] & HR02$^{\circ}$, HR11$^{\circ}$ & \textbf{X} \\
HR.I03 & A single favourable variation is swamped by crossing within a few generations. & A rare favourable variant is diluted to nothing by crossing within a few generations. [Denied: Mendelian alleles are not diluted.] & HR02$^{\circ}$, HR10$^{\circ}$ & \textbf{X} \\
HR.I04 & In crosses one parent is often prepotent in transmitting a character. & In crosses one allele is often dominant in its effect on a trait. & HR03, HR11$^{\circ}$ & \textbf{H} \\
HR.I05 & Crossing distinct pure breeds of pigeon can produce the colouring of the rock pigeon. & Crossing distinct homozygous breeds of pigeon can produce the rock-pigeon phenotype (by recombination). & HR04$^{\circ}$, HR07$^{\circ}$ & \textbf{H} \\
HR.I06 & Reversion is due to dormant gemmules transmitted through generations. & Reversion is due to recessive alleles transmitted unexpressed through generations. [Holds for reversion by recessives; reversions by recombination are explained otherwise.] & HR04$^{\circ}$, HR01b & \textbf{R} \\
HR.I07 & Characters acquired by use or disuse are inherited. & Interpretation: effects of use and disuse are transmitted. T2 (with Weismann) recognizes no mechanism for this. & HR05a$^{\ast}$, HR11$^{\circ}$ & \textbf{U} \\
HR.I08 & A previous sire influences the offspring of later matings. & Interpretation: a previous sire affects later offspring. T2 recognizes no route for such influence. & HR06a$^{\ast}$ & \textbf{U} \\
HR.I09 & Pure breeds breed true. & Homozygous lines breed true. & HR07$^{\circ}$ & \textbf{H} \\
HR.I10 & The offspring of hybrids are variable, some resembling each grandparent. & The offspring of heterozygotes are variable, some resembling each grandparent (segregation). & HR08$^{\circ}$, HR11$^{\circ}$ & \textbf{H} \\
HR.I11 & The offspring of hybrids remain uniformly intermediate. & The offspring of heterozygotes are uniformly intermediate. [Denied: they segregate.] & HR08$^{\circ}$, HR02$^{\circ}$ & \textbf{X} \\
HR.I12 & Sports are sometimes inherited. & Mutations are sometimes inherited. & HR09 & \textbf{H} \\
HR.I13 & Individual variation is heritable in some degree. & Individual variation is heritable in some degree. [Holds on the reading heritable variation; fails for variation induced by the environment.] & HR10$^{\circ}$ & \textbf{R} \\
HR.I14 & Selection within a pure breed shifts its mean. & Selection within a homozygous line shifts its mean. [Denied: Johannsen's pure lines (1903).] & HR10$^{\circ}$, HR07$^{\circ}$ & \textbf{X} \\
HR.I15 & Each offspring receives hereditary tendencies from both parents. & Each offspring receives its genotype from both parents. & HR12, HR17$^{\circ}$ & \textbf{H} \\
HR.I16 & The offspring of exceptional parents regress toward the mean of the race. & The offspring of exceptional parents regress toward the population mean. & HR14 & \textbf{H} \\
HR.I17 & Parents contribute one half of the heritage, grandparents one quarter, and so on. & Interpretation: the heritage is composed of fixed fractions from each ancestor. T2 has no such fractions; each parent transmits one allele at each locus. & HR13a$^{\ast}$ & \textbf{U} \\
HR.I18 & Filial stature regresses on mid-parental stature with a coefficient of about two thirds. & Filial stature regresses on mid-parental stature with a coefficient of about two thirds. & HR13b, HR14 & \textbf{H} \\
HR.I19 & The hereditary substance is continuous and is not affected by the body. & The germ-line hereditary material is continuous and is not altered by the body. & HR15 & \textbf{H} \\
HR.I20 & A character can lie latent for a generation and reappear. & A recessive trait can skip a generation and reappear. & HR16, HR11$^{\circ}$ & \textbf{H} \\
HR.I21 & Each parent contributes one reproductive element to the offspring. & Each parent contributes one gamete to the offspring. & HR17$^{\circ}$ & \textbf{H} \\
HR.I22 & The intermediate stature of children of tall and short parents is due to blending. & The intermediate stature of children of tall and short parents is due to quantitative inheritance. [Holds: many factors with additive effect give intermediate means.] & HR02$^{\circ}$, HR11$^{\circ}$ & \textbf{R} \\
HR.I23 & A dilute rare variation in a large population cannot be raised in frequency by selection. & A rare variant in a large population cannot be raised in frequency by selection. [Denied: selection raises the frequency of an undiluted allele.] & HR02$^{\circ}$, HR10$^{\circ}$ & \textbf{X} \\
\bottomrule
\end{inftable}



\begin{thebibliography}{99}\small
\bibitem[Andersen et~al.(2006)]{AndersenBarkerChen2006} Andersen, H., Barker, P. and Chen, X. (2006). \emph{The Cognitive Structure of Scientific Revolutions}. Cambridge University Press.
\bibitem[Bird(2000)]{Bird2000} Bird, A. (2000). \emph{Thomas Kuhn}. Acumen.
\bibitem[Bird(2008)]{Bird2008} Bird, A. (2008). Incommensurability naturalized. In L. Soler, H. Sankey and P. Hoyningen-Huene (eds.), \emph{Rethinking Scientific Change and Theory Comparison: Stabilities, Ruptures, Incommensurabilities?}, 21--39. Springer.
\bibitem[Buchwald(1989)]{Buchwald1989} Buchwald, J.~Z. (1989). \emph{The Rise of the Wave Theory of Light}. University of Chicago Press.
\bibitem[Callendar(1911)]{Callendar1911} Callendar, H.~L. (1911). The caloric theory of heat and Carnot's principle. \emph{Proceedings of the Physical Society of London}, 23.
\bibitem[Chang(2012)]{Chang2012} Chang, H. (2012). \emph{Is Water H$_2$O? Evidence, Realism and Pluralism}. Springer.
\bibitem[Darwin(1859)]{Darwin1859} Darwin, C. (1859). \emph{On the Origin of Species by Means of Natural Selection}. John Murray.
\bibitem[Davidson(1974)]{Davidson1974} Davidson, D. (1974). On the very idea of a conceptual scheme. \emph{Proceedings and Addresses of the American Philosophical Association}, 47, 5--20.
\bibitem[Devitt(1979)]{Devitt1979} Devitt, M. (1979). Against incommensurability. \emph{Australasian Journal of Philosophy}, 57(1), 29--50.
\bibitem[Feyerabend(1962)]{Feyerabend1962} Feyerabend, P.~K. (1962). Explanation, reduction, and empiricism. In H. Feigl and G. Maxwell (eds.), \emph{Scientific Explanation, Space, and Time}, Minnesota Studies in the Philosophy of Science 3, 28--97. University of Minnesota Press.
\bibitem[Feyerabend(1975)]{Feyerabend1975} Feyerabend, P.~K. (1975). \emph{Against Method}. New Left Books.
\bibitem[Field(1973)]{Field1973} Field, H. (1973). Theory change and the indeterminacy of reference. \emph{The Journal of Philosophy}, 70(14), 462--481.
\bibitem[Fox(1971)]{Fox1971} Fox, R. (1971). \emph{The Caloric Theory of Gases from Lavoisier to Regnault}. Clarendon Press.
\bibitem[French(1994)]{French1994} French, R. (1994). \emph{William Harvey's Natural Philosophy}. Cambridge University Press.
\bibitem[Fisher(1918)]{Fisher1918} Fisher, R.~A. (1918). The correlation between relatives on the supposition of Mendelian inheritance. \emph{Transactions of the Royal Society of Edinburgh}, 52, 399--433.
\bibitem[Hacking(1993)]{Hacking1993} Hacking, I. (1993). Working in a new world: the taxonomic solution. In P. Horwich (ed.), \emph{World Changes: Thomas Kuhn and the Nature of Science}. MIT Press.
\bibitem[Hanson(1958)]{Hanson1958} Hanson, N.~R. (1958). \emph{Patterns of Discovery}. Cambridge University Press.
\bibitem[Hoyningen-Huene(1993)]{HoyningenHuene1993} Hoyningen-Huene, P. (1993). \emph{Reconstructing Scientific Revolutions: Thomas S. Kuhn's Philosophy of Science}. University of Chicago Press.
\bibitem[Hoyningen-Huene and Sankey(2001)]{HoyningenHueneSankey2001} Hoyningen-Huene, P. and Sankey, H. (eds.) (2001). \emph{Incommensurability and Related Matters}. Kluwer.
\bibitem[Jammer(1966)]{Jammer1966} Jammer, M. (1966). \emph{The Conceptual Development of Quantum Mechanics}. McGraw-Hill.
\bibitem[Kitcher(1978)]{Kitcher1978} Kitcher, P. (1978). Theories, theorists and theoretical change. \emph{The Philosophical Review}, 87(4), 519--547.
\bibitem[Kitcher(1993)]{Kitcher1993} Kitcher, P. (1993). \emph{The Advancement of Science: Science without Legend, Objectivity without Illusions}. Oxford University Press.
\bibitem[Kuhn(1957)]{Kuhn1957} Kuhn, T.~S. (1957). \emph{The Copernican Revolution}. Harvard University Press.
\bibitem[Kuhn(1962)]{Kuhn1962} Kuhn, T.~S. (1962). \emph{The Structure of Scientific Revolutions}. University of Chicago Press.
\bibitem[Kuhn(1970)]{Kuhn1970} Kuhn, T.~S. (1970). \emph{The Structure of Scientific Revolutions}. 2nd edition. University of Chicago Press.
\bibitem[Kuhn(1977)]{Kuhn1977} Kuhn, T.~S. (1977). Objectivity, value judgment, and theory choice. In \emph{The Essential Tension: Selected Studies in Scientific Tradition and Change}, 320--339. University of Chicago Press.
\bibitem[Kuhn(1983)]{Kuhn1983} Kuhn, T.~S. (1983). Commensurability, comparability, communicability. In P.~D. Asquith and T. Nickles (eds.), \emph{PSA 1982}, vol.~2, 669--688. Philosophy of Science Association.
\bibitem[Kuhn(1991)]{Kuhn1991} Kuhn, T.~S. (1991). The road since \emph{Structure}. In A. Fine, M. Forbes and L. Wessels (eds.), \emph{PSA 1990}, vol.~2, 3--13. Philosophy of Science Association.
\bibitem[Kuhn(1993)]{Kuhn1993} Kuhn, T.~S. (1993). Afterwords. In P. Horwich (ed.), \emph{World Changes: Thomas Kuhn and the Nature of Science}, 311--341. MIT Press.
\bibitem[Lakatos(1970)]{Lakatos1970} Lakatos, I. (1970). Falsification and the methodology of scientific research programmes. In I. Lakatos and A. Musgrave (eds.), \emph{Criticism and the Growth of Knowledge}, 91--196. Cambridge University Press.
\bibitem[Laudan(1977)]{Laudan1977} Laudan, L. (1977). \emph{Progress and Its Problems: Towards a Theory of Scientific Growth}. University of California Press.
\bibitem[Laudan(1984)]{Laudan1984} Laudan, L. (1984). \emph{Science and Values: The Aims of Science and Their Role in Scientific Debate}. University of California Press.
\bibitem[McMullin(1983)]{McMullin1983} McMullin, E. (1983). Values in science. In P.~D. Asquith and T. Nickles (eds.), \emph{PSA 1982}, vol.~2, 3--28. Philosophy of Science Association.
\bibitem[Musgrave(1976)]{Musgrave1976} Musgrave, A. (1976). Why did oxygen supplant phlogiston? Research programmes in the chemical revolution. In C. Howson (ed.), \emph{Method and Appraisal in the Physical Sciences}. Cambridge University Press.
\bibitem[Okasha(2011)]{Okasha2011} Okasha, S. (2011). Theory choice and social choice: Kuhn versus Arrow. \emph{Mind}, 120(477), 83--115.
\bibitem[Putnam(1975)]{Putnam1975} Putnam, H. (1975). The meaning of \enquote{meaning}. In K. Gunderson (ed.), \emph{Language, Mind, and Knowledge}, Minnesota Studies in the Philosophy of Science 7, 131--193. University of Minnesota Press.
\bibitem[Putnam(1981)]{Putnam1981} Putnam, H. (1981). \emph{Reason, Truth and History}. Cambridge University Press.
\bibitem[Quine(1960)]{Quine1960} Quine, W.~V.~O. (1960). \emph{Word and Object}. MIT Press.
\bibitem[Ruse(1979)]{Ruse1979} Ruse, M. (1979). \emph{The Darwinian Revolution}. University of Chicago Press.
\bibitem[Salmon(1990)]{Salmon1990} Salmon, W.~C. (1990). Rationality and objectivity in science, or Tom Kuhn meets Tom Bayes. In C.~W. Savage (ed.), \emph{Scientific Theories}, Minnesota Studies in the Philosophy of Science 14, 175--204. University of Minnesota Press.
\bibitem[Sankey(1994)]{Sankey1994} Sankey, H. (1994). \emph{The Incommensurability Thesis}. Avebury.
\bibitem[Sankey(1997)]{Sankey1997} Sankey, H. (1997). \emph{Rationality, Relativism and Incommensurability}. Ashgate.
\bibitem[Scheffler(1967)]{Scheffler1967} Scheffler, I. (1967). \emph{Science and Subjectivity}. Bobbs-Merrill.
\bibitem[Shapere(1966)]{Shapere1966} Shapere, D. (1966). Meaning and scientific change. In R.~G. Colodny (ed.), \emph{Mind and Cosmos: Essays in Contemporary Science and Philosophy}, 41--85. University of Pittsburgh Press.
\bibitem[Thagard(1992)]{Thagard1992} Thagard, P. (1992). \emph{Conceptual Revolutions}. Princeton University Press.
\end{thebibliography}

\end{document}
